% concatenated from main.tex, math.tex
\documentclass[11pt]{article}

\usepackage[utf8]{inputenc}
\usepackage[T1]{fontenc}
\usepackage[margin=1in]{geometry}
\usepackage{amsmath,amsfonts,amssymb,amsthm}
\usepackage{bm}
\usepackage{microtype}
\usepackage{graphicx}
\usepackage{float}
\usepackage[table]{xcolor}
\usepackage[para,online,flushleft]{threeparttable}
\usepackage{booktabs}
\usepackage{enumitem}
\usepackage{algorithm}
\usepackage{algorithmic}
\usepackage{placeins}
\usepackage{bbding}
\usepackage{makecell}
\usepackage{hhline}
\usepackage{natbib}
\usepackage{hyperref}

\hypersetup{
  hidelinks,
  pdfauthor={Chengchang Liu and Luo Luo},
  pdftitle={Randomized Quasi-Gauss--Newton Methods for Solving General Nonlinear Equations},
  pdfsubject={Quasi-Newton methods for nonlinear equations},
  pdfkeywords={quasi-Newton methods, nonlinear equations, nonlinear least squares, randomized algorithms, local convergence}
}

\theoremstyle{plain}
\newtheorem{theorem}{Theorem}
\newtheorem{lemma}[theorem]{Lemma}
\newtheorem{proposition}[theorem]{Proposition}
\newtheorem{corollary}[theorem]{Corollary}
\newtheorem{assumption}[theorem]{Assumption}
\theoremstyle{definition}
\newtheorem{definition}[theorem]{Definition}
\theoremstyle{remark}
\newtheorem{remark}[theorem]{Remark}
\theoremstyle{plain}

\newenvironment{keywords}
  {\par\medskip\noindent\small\textbf{Keywords:} }
  {\par\medskip}

% Mark sections of captions for referring to divisions of figures
\newcommand{\figleft}{{\em (Left)}}
\newcommand{\figcenter}{{\em (Center)}}
\newcommand{\figright}{{\em (Right)}}
\newcommand{\figtop}{{\em (Top)}}
\newcommand{\figbottom}{{\em (Bottom)}}
\newcommand{\captiona}{{\em (a)}}
\newcommand{\captionb}{{\em (b)}}
\newcommand{\captionc}{{\em (c)}}
\newcommand{\captiond}{{\em (d)}}

% Highlight a newly defined term
\newcommand{\newterm}[1]{{\bf #1}}


% Figure reference, lower-case.
\def\figref#1{figure~\ref{#1}}
% Figure reference, capital. For start of sentence
\def\Figref#1{Figure~\ref{#1}}
\def\twofigref#1#2{figures \ref{#1} and \ref{#2}}
\def\quadfigref#1#2#3#4{figures \ref{#1}, \ref{#2}, \ref{#3} and \ref{#4}}
% Section reference, lower-case.
\def\secref#1{section~\ref{#1}}
% Section reference, capital.
\def\Secref#1{Section~\ref{#1}}
% Reference to two sections.
\def\twosecrefs#1#2{sections \ref{#1} and \ref{#2}}
% Reference to three sections.
\def\secrefs#1#2#3{sections \ref{#1}, \ref{#2} and \ref{#3}}
% Reference to an equation, lower-case.
%\def\eqref#1{equation~(\ref{#1})}
% Reference to an equation, upper case
\def\Eqref#1{Equation~\ref{#1}}
% A raw reference to an equation---avoid using if possible
\def\plaineqref#1{\ref{#1}}
% Reference to a chapter, lower-case.
\def\chapref#1{chapter~\ref{#1}}
% Reference to an equation, upper case.
\def\Chapref#1{Chapter~\ref{#1}}
% Reference to a range of chapters
\def\rangechapref#1#2{chapters\ref{#1}--\ref{#2}}
% Reference to an algorithm, lower-case.
\def\algref#1{algorithm~\ref{#1}}
% Reference to an algorithm, upper case.
\def\Algref#1{Algorithm~\ref{#1}}
\def\twoalgref#1#2{algorithms \ref{#1} and \ref{#2}}
\def\Twoalgref#1#2{Algorithms \ref{#1} and \ref{#2}}
% Reference to a part, lower case
\def\partref#1{part~\ref{#1}}
% Reference to a part, upper case
\def\Partref#1{Part~\ref{#1}}
\def\twopartref#1#2{parts \ref{#1} and \ref{#2}}

\def\ceil#1{\lceil #1 \rceil}
\def\floor#1{\left\lfloor #1 \right\rfloor}
\def\1{\bf{1}}
\newcommand{\train}{\mathcal{D}}
\newcommand{\valid}{\mathcal{D_{\mathrm{valid}}}}
\newcommand{\test}{\mathcal{D_{\mathrm{test}}}}
\newcommand{\Norm}[1]{\left\| #1 \right\|}
\newcommand{\norm}[1]{\left\| #1 \right\|_2}
\def\inner#1#2{\langle #1, #2 \rangle}
\def\Inner#1#2{\left\langle #1, #2 \right\rangle}
\def\eps{{\varepsilon}}
\def \abs#1{\left|#1\right|}
% Vectors
\def\vzero{{\bf{0}}}
\def\vone{{\bf{1}}}
\def\vmu{{\bm{\mu}}}
\def\vphi{{\bm{\phi}}}
\def\veta{{\bm{\eta}}}
\def\valpha{{\bm{\alpha}}}
\def\vbeta{{\bm{\beta}}}
\def\vtheta{{\bm{\theta}}}
\def\va{{\bf{a}}}
\def\vb{{\bf{b}}}
\def\vc{{\bf{c}}}
\def\vd{{\bf{d}}}
\def\ve{{\bf{e}}}
\def\vf{{\bf{f}}}
\def\vg{{\bf{g}}}
\def\vh{{\bf{h}}}
\def\vi{{\bf{i}}}
\def\vj{{\bf{j}}}
\def\vk{{\bf{k}}}
\def\vl{{\bf{l}}}
\def\vm{{\bf{m}}}
\def\vn{{\bf{n}}}
\def\vo{{\bf{o}}}
\def\vp{{\bf{p}}}
\def\vq{{\bf{q}}}
\def\vr{{\bf{r}}}
\def\vs{{\bf{s}}}
\def\vt{{\bf{t}}}
\def\vu{{\bf{u}}}
\def\vv{{\bf{v}}}
\def\vw{{\bf{w}}}
\def\vx{{\bf{x}}}
\def\vy{{\bf{y}}}
\def\vz{{\bf{z}}}


% Graph
\def\fA{{\mathcal{A}}}
\def\fB{{\mathcal{B}}}
\def\fC{{\mathcal{C}}}
\def\fD{{\mathcal{D}}}
\def\fE{{\mathcal{E}}}
\def\fF{{\mathcal{F}}}
\def\fG{{\mathcal{G}}}
\def\fH{{\mathcal{H}}}
\def\fI{{\mathcal{I}}}
\def\fJ{{\mathcal{J}}}
\def\fK{{\mathcal{K}}}
\def\fL{{\mathcal{L}}}
\def\fM{{\mathcal{M}}}
\def\fN{{\mathcal{N}}}
\def\fO{{\mathcal{O}}}
\def\fP{{\mathcal{P}}}
\def\fQ{{\mathcal{Q}}}
\def\fR{{\mathcal{R}}}
\def\fS{{\mathcal{S}}}
\def\fT{{\mathcal{T}}}
\def\fU{{\mathcal{U}}}
\def\fV{{\mathcal{V}}}
\def\fW{{\mathcal{W}}}
\def\fX{{\mathcal{X}}}
\def\fY{{\mathcal{Y}}}
\def\fZ{{\mathcal{Z}}}

% Sets
\def\sA{{\mathbb{A}}}
\def\sB{{\mathbb{B}}}
\def\sC{{\mathbb{C}}}
\def\sD{{\mathbb{D}}}
\def\BE{{\mathbb{E}}}
% Don't use a set called E, because this would be the same as our symbol
% for expectation.
\def\BF{{\mathbb{F}}}
\def\BG{{\mathbb{G}}}
\def\BH{{\mathbb{H}}}
\def\BI{{\mathbb{I}}}
\def\BJ{{\mathbb{J}}}
\def\BK{{\mathbb{K}}}
\def\BL{{\mathbb{L}}}
\def\BM{{\mathbb{M}}}
\def\BN{{\mathbb{N}}}
\def\BO{{\mathbb{O}}}
\def\BP{{\mathbb{P}}}
\def\BQ{{\mathbb{Q}}}
\def\BR{{\mathbb{R}}}
\def\BS{{\mathbb{S}}}
\def\BT{{\mathbb{T}}}
\def\BU{{\mathbb{U}}}
\def\BV{{\mathbb{V}}}
\def\BW{{\mathbb{W}}}
\def\BX{{\mathbb{X}}}
\def\BY{{\mathbb{Y}}}
\def\BZ{{\mathbb{Z}}}

\def\mA {{\bf A}}
\def\mB {{\bf B}}
\def\mC {{\bf C}}
\def\mD {{\bf D}}
\def\mE {{\bf E}}
\def\mF {{\bf F}}
\def\mG {{\bf G}}
\def\mH {{\bf H}}
\def\mI {{\bf I}}
\def\mJ {{\bf J}}
\def\mK {{\bf K}}
\def\mL {{\bf L}}
\def\mM {{\bf M}}
\def\mN {{\bf N}}
\def\mO {{\bf O}}
\def\mP {{\bf P}}
\def\mQ {{\bf Q}}
\def\mR {{\bf R}}
\def\mS {{\bf S}}
\def\mT {{\bf T}}
\def\mU {{\bf U}}
\def\mV {{\bf V}}
\def\mW {{\bf W}}
\def\mX {{\bf X}}
\def\mY {{\bf Y}}
\def\mZ {{\bf Z}}
\def \tmx {{T_M(x)}}
\def\AUC{{\rm{AUC}}}
\def\SVI{{\rm{SVI}}}
\def\VI{{\rm{VI}}}
\def\ApproxMax{{\rm{ApproxMax}}}


% The true underlying data generating distribution
\newcommand{\pdata}{p_{\rm{data}}}
% The empirical distribution defined by the training set
\newcommand{\ptrain}{\hat{p}_{\rm{data}}}
\newcommand{\Ptrain}{\hat{P}_{\rm{data}}}
% The model distribution
\newcommand{\pmodel}{p_{\rm{model}}}
\newcommand{\Pmodel}{P_{\rm{model}}}
\newcommand{\ptildemodel}{\tilde{p}_{\rm{model}}}
% Stochastic autoencoder distributions
\newcommand{\pencode}{p_{\rm{encoder}}}
\newcommand{\pdecode}{p_{\rm{decoder}}}
\newcommand{\precons}{p_{\rm{reconstruct}}}

\newcommand{\laplace}{\mathrm{Laplace}} % Laplace distribution

\newcommand{\E}{\mathbb{E}}
\newcommand{\Ls}{\mathcal{L}}
\newcommand{\R}{\mathbb{R}}
\newcommand{\emp}{\tilde{p}}
\newcommand{\lr}{\alpha}
\newcommand{\reg}{\lambda}
\newcommand{\rect}{\mathrm{rectifier}}
\newcommand{\softmax}{\mathrm{softmax}}
\newcommand{\sigmoid}{\sigma}
\newcommand{\softplus}{\zeta}
\newcommand{\KL}{D_{\mathrm{KL}}}
\newcommand{\Var}{\mathrm{Var}}
\newcommand{\standarderror}{\mathrm{SE}}
\newcommand{\Cov}{\mathrm{Cov}}
% Wolfram Mathworld says $L^2$ is for function spaces and $\ell^2$ is for vectors
% But then they seem to use $L^2$ for vectors throughout the site, and so does
% wikipedia.
\newcommand{\normlzero}{L^0}
\newcommand{\normlone}{L^1}
\newcommand{\normltwo}{L^2}
\newcommand{\normlp}{L^p}
\newcommand{\normmax}{L^\infty}

\newcommand{\parents}{Pa} % See usage in notation.tex. Chosen to match Daphne's book.

\DeclareMathOperator*{\argmax}{arg\,max}
\DeclareMathOperator*{\argmin}{arg\,min}

\DeclareMathOperator{\sign}{sign}
\DeclareMathOperator{\Tr}{Tr}
\let\ab\allowbreak

\DeclareMathOperator{\diag}{diag}
\DeclareMathOperator{\spn}{span}
\DeclareMathOperator{\prox}{prox}

\makeatletter
\def\Ddots{\mathinner{\mkern1mu\raise\p@
\vbox{\kern7\p@\hbox{.}}\mkern2mu
\raise4\p@\hbox{.}\mkern2mu\raise7\p@\hbox{.}\mkern1mu}}
\makeatother

\def\pr#1{\sP\left( #1 \right)}

\makeatletter
\newcommand*{\rom}[1]{\expandafter\@slowromancap\romannumeral #1@}
\makeatother

\newcommand{\Hnswer}[1]{\begin{quote}{\color{blue}#1}\end{quote}}
\newcommand{\HND}{\wedge}
%\newcommand{\OR}{\vee}
\newcommand{\ra}{\rightarrow}
\newcommand{\lra}{\leftrightarrow}
 
\def\vx {{{\bf x}}}
\def\vw {{{\bf w}}}
\def\vy {{{\bf y}}}
\def\va {{{\bf a}}}
\def\BR{{\mathbb{R}}}
\def\Hl{{{\widetilde{H}}}}

\def\A{{\bf A}}
\def\B{{\bf B}}
\def\bb{{\bf b}}
\def\C{{\bf C}}
\def\c{{\bf c}}
\def\D{{\bf D}}
\def\E{{\bf E}}
\def\e{{\bf e}}
\def\F{{\bf F}}
\def\f{{\bf f}}
\def\g{{\bf g}}
\def\h{{\bf h}}
\def\G{{\bf G}}
\def\H{{\bf H}}
\def\I{{\bf I}}
\def \J{{\bf J}}
\def\K{{\bf K}}
\def\k{{\bf k}}
\def\LL{{\bf L}}
\def\l{{\bf l}}
\def\M{{\bf M}}
\def\m{{\bf m}}
\def\N{{\bf N}}
\def\n{{\bf n}}
\def\oo{{\bf o}}
\def\PP{{\bf P}}
\def\pp{{\bf p}}
\def\Q{{\bf Q}}
\def\q{{\bf q}}
\def\R{{\bf R}}
\def\rr{{\bf r}}
\def\S{{\bf S}}
\def\s{{\bf s}}
\def\T{{\bf T}}
\def\tt{{\bf t}}
\def\U{{\bf U}}
\def\u{{\bf u}}
\def\V{{\bf V}}
\def\v{{\bf v}}
\def\W{{\bf W}}
\def\w{{\bf w}}
\def\X{{\bf X}}
\def\x{{\bf x}}
\def\Y{{\bf Y}}
\def\y{{\bf y}}
\def\Z{{\bf Z}}
\def\z{{\bf z}}
\def\0{{\bf 0}}
\def\1{{\bf 1}}
\def\bu{{\bar \vu}}
\def\bH{{\widetilde{\H}}}
\def\HM{{\mathcal A}}
\def\CM{{\mathcal C}}
\def\DM{{\mathcal D}}
\def\EM{{\mathcal E}}
\def\GM{{\mathcal G}}
\def\FM{{\mathcal F}}
\def\IM{{\mathcal I}}
\def\JM{{\mathcal J}}
\def\KM{{\mathcal K}}
\def\LM{{\mathcal L}}
\def\NM{{\mathcal N}}
\def\OM{{\mathcal O}}
\def\PM{{\mathcal P}}
\def\SM{{\mathcal S}}
\def\TM{{\mathcal T}}
\def\UM{{\mathcal U}}
\def\VM{{\mathcal V}}
\def\WM{{\mathcal W}}
\def\XM{{\mathcal X}}
\def\YM{{\mathcal Y}}
\def\ZM{{\mathcal Z}}
\def\CB{{\mathbb C}}
\def\RB{{\mathbb R}}
\def\EB{{\mathbb E}}
\def\PB{{\mathbb P}}
\def\BH{{\widetilde{\H}}}
\def \Bg{{\widetilde{\G}}}
\def \hH{{\hat{\H}}}
\def \hG{{\hat{\G}}}
\def \hn{{\hat{n}}}
\def \hu{{\hat{\vu}}}

\def\TT{{\rm{T}}}
\def \hgamma {{\gamma(\z_k;\nu,2\fH_\nu)}}
\def \adagamma {{\gamma(\z_k;\nu,H_k2^{i_k})}}
\def \padagamma {{\gamma_{p}(\z_k;\nu,H_k2^{i_k})}}
\def \unigamma {{\gamma(\z_k;1, H_k2^{i_k})}}
\def \punigamma{{\gamma_{p}(\z_k;1,H_k2^{i_k})}}

\def\Holder{{H\oo lder}}
\def\defeq{\overset{\text{def}}{=}}
\def \EBP #1{\EB\left[#1\right]}
\def \tr #1{\text{\rm tr}\left(#1\right)}
\def\Holder{{H\oo lder}}
\def \oo{ö}

\def \JJ #1{{\left(\J(#1)^{\top}\J(#1)\right)}}
\def \JJO #1{{\J(#1)^{\top}\J(#1)}}
\def \srk{{\text{\rm SR1}}}
\def \brk{{\text{BR$k$}}}
\def \bfgs{{\text{\rm BFGS}}}
\def \fbfgs{{\text{\rm F-BFGS}}}
\def \dfp{{\text{BlockDFP}}}
\def \L{{\bf L}}
%\def\rddots{\cdot^{\cdot^{\cdot}}


\begin{document}

\title{Randomized Quasi-Gauss--Newton Methods for Solving General Nonlinear Equations}
\author{Chengchang Liu \and Luo Luo}
\date{}
\maketitle
\begin{abstract}
This paper considers the local convergence for solving general nonlinear equations. 
We establish randomized quasi-Gauss--Newton methods based on the approximation of the Gram matrix, addressing both the underdetermined and the overdetermined settings.
For the underdetermined case, we show that our methods achieve local superlinear convergence to the optimal solution.
For the overdetermined case, we show that our methods achieve local condition-number-free convergence to the stationary point of the nonlinear least-square formulation of the problem.
In contrast, existing quasi-Newton methods mostly focus on square systems and additionally require the initial Jacobian estimate to be sufficiently close to the exact one.
\end{abstract}

\begin{keywords}
quasi-Newton methods, nonlinear equations, nonlinear least squares, randomized algorithms, explicit superlinear convergence
\end{keywords}

% {\color{red}
% To-do List:
% \begin{itemize}
%     \item check the proof, add the reference, and determine the size of the local region ($\delta$)
%     \item write the introduction and preliminaries (all related papers).
%     \item experiment.
%     \item add globally convergence technique?
% \end{itemize}
% }
\section{Introduction}
We study the problem of solving the nonlinear equation
\begin{align}\label{eq:nonlinear_obj}
    \F(\x) = \0,
\end{align}
where $\F(\cdot)\triangleq [F_1(\cdot);\cdots;F_n(\cdot)]^{\top}$ is a general mapping such that $\F:\RB^{d}\to\RB^n$ with differentiable $F_i:\RB^d\to\RB$ for all $i\in[n]$. 
We also consider its nonlinear least square formulation
\begin{align}
    \label{eq:nonlinear_least_obj}
    \min_{\x\in\RB^d} \phi(\x)\triangleq\frac{1}{2}\|\F(\x)\|_2^2.
\end{align}
% When the solution of above equations system does not exist, we also consider the following nonlinear least squares problem
% \begin{align}
%     \label{eq:nonlinear_least_obj}
%     \min_{\x\in\RB^d} \phi(\x)\triangleq\frac{1}{2}\|\F(\x)\|_2^2.
% \end{align}
% In the remaining context, we use $\x^*$ to denote the solution of a nonlinear equations system such that $\F(\x^*)=\0$ or the solution of nonlinear least squares such that $\partial\F(\x^*)^{\top}\F(\x^*)=\0$ when the prior one does not exist.
Numerous applications in machine learning~\citep{scieur2018nonlinear,alizamir2020comparative,song2021accelerating}, game theory~\citep{frehse1984nonlinear,nourian2013e}, economics~\citep{albu2006non}, and control systems~\citep{more1989collection,berthier2021fast} can be formulated as above problems, which have been extensively studied for decades and be considered one of the most important in scientific computing~\citep{nesterov2007modified}.

The Gauss--Newton method is very popular for solving nonlinear equations system, and it iterates according to
\begin{align}  
\label{update:gn}
    \x_{+} = \x-\J(\x)^{\dag}\F(\x),
\end{align}
where we use $\J(\x)\in\BR^{n\times d}$ to denote the Jacobian of $\F(\cdot)$ at $\vx$ and the notation $(\,\cdot\,)^{\dag}$ presents the pseudoinverse.
We are interested in the case that the Jacobian in the local region of the solution is full rank, which implies 
\begin{align}\label{update:gnJ}
\J(\x)^{\dag}=\begin{cases}
    \J(\x)^{\top}(\J(\x)\J(\x)^{\top})^{-1}, &\text{if }n<d,\\
	 (\J(\x)^{\top}\J(\x))^{-1}\J(\x)^{\top}, &\text{if }n\geq d.    
\end{cases}
\end{align}
It is well-known that the Gauss--Newton method enjoys faster local convergence rates than simply applying gradient descent to solve problem (\ref{eq:nonlinear_least_obj})~\citep{gratton2007approximate,ortega2000iterative,nesterov2018lectures}. 
%It enjoys local linear or superlinear convergence rate when the Jacobian at a solution of \eqref{eq:nonlinear_obj} or \eqref{eq:nonlinear_least_obj} has full rank~\citep{gratton2007approximate,ortega2000iterative}. 
%As a comparison, simply applying first-order methods to solve problem (\ref{eq:nonlinear_least_obj}) only achieves a sublinear rate for finding an approximate stationary point~\citep{nesterov2018lectures}.
However, accessing the exact Gram matrix $\J(\x)^{\top}\J(\x)$ 
or~$\J(\x)\J(\x)^{\top}$ and 
computing its inverse require the computational cost of 
$\OM(\min\{n,d\}nd)$, which is impractical for large-scale problems.
% In the remaining context, we use $\x^*$ to denote the solution of a nonlinear equations system such that $\F(\x^*)=\0$ or the solution of nonlinear least squares such that $\partial\F(\x^*)^{\top}\F(\x^*)=\0$ when the prior one does not exist.
Therefore, we desire to reduce the cost of each iteration while maintaining a good convergence behavior.

Quasi-Newton methods are popular to solve nonlinear equations~\citep{yuan2011recent,yuan2022sketched,yuan2009subspace,broyden1967quasi,martinez1991quasi,broyden1967quasi,davidon1991variable,broyden1970convergence,broyden1970convergence2,davidon1991variable,fletcher1963rapidly}. 
They use low-rank modification to approximate the Jacobian or its inverse, with relatively low computational costs, and often enjoy fast local convergence rates.
%The study of quasi-Newton methods dates back to the 1960s. 
%The early proposed methods include SR1 method~\citep{broyden1967quasi,davidon1991variable}, BFGS method~\citep{broyden1970convergence,broyden1970convergence2}, and DFP method~\citep{davidon1991variable,fletcher1963rapidly}.
%\citet{li1999globally,gu2002descent} applied the BFGS update to approximate $\J(\x)^{\top}\J(\x)$ for nonlinear equations with symmetric Jacobian.
However, classical quasi-Newton methods such as Broyden's methods \citep{broyden2000discovery}, 
SR1 method~\citep{broyden1967quasi,davidon1991variable}, BFGS method~\citep{broyden1970convergence,broyden1970convergence2}, and DFP method~\citep{davidon1991variable,fletcher1963rapidly} only focus on the square systems. %, and require the Jacobian be symmetric or even positive-definite.
%However, for the general nonlinear equations~\eqref{eq:nonlinear_obj} or for the nonlinear least squares problem~\eqref{eq:nonlinear_least_obj}, the Jacobian of $\F(\x)$ and the Hessian of $\nabla^2\phi(\x)$ are not necessarily positive definite or even symmetric.
% In our early version~\citep{liu2022quasi}, we applied randomized quasi-Newton updates~\citep{rodomanov2021greedy,lin2022explicit} to approximate the square of the Hessian in minimax optimization, which can be extended to solve nonlinear equations when $\J(\x)$ is square.
% The Broyden's methods~\citep{broyden2000discovery}, including the good and the bad schemes, are popular quasi-Newton methods to solve nonlinear equations with square Jacobians.
% They use Broyden updates to approximate $\J(\x)$ or its inverse directly and enjoy local superlinear rates when the Jacobian at a solution is nondegenerate.
The non-asymptotic local convergence of these methods has been established in recent years \citep{rodomanov2021greedy,rodomanov2021new,rodomanov2021rates,ye2022towards,liu2022quasi,lin2022explicit,liu2023symmetric,liu2023block}, while exiting results cannot be applied to more general settings.
% In recent years, the non-asymptotic local convergence of these methods have been extensively studied \citep{rodomanov2021greedy,rodomanov2021new,rodomanov2021rates,ye2022towards,liu2022quasi,lin2022explicit,liu2023block}, while how to extend the theory to more general nonlinear systems is still unclear.

There are also several works that consider general nonlinear equations. 
Specifically, \citet{martinez1991quasi,vater2024convergence} proposed quasi-Newton methods for the underdetermined systems (i.e., the case of $n<d$) by generalizing Broyden's updates to approximate the pseudoinverse of the Jacobian, but they only achieve the asymptotic local convergence.
In addition, their convergence guarantees require the assumption that the initial estimator of the Jacobian (or its inverse) to be sufficiently close to the exact one.
Although this assumption is widely adopted in the analysis~\citep{asl2024aj,martinez1991quasi,vater2024convergence,liu2023block}, 
achieving a sufficiently accurate initial estimator may be so expensive.
For the overdetermined case (i.e., the case of $n\geq d$), the nonlinear equation~\eqref{eq:nonlinear_obj} may has no solution so that we typically consider its nonlinear least square formulation~\eqref{eq:nonlinear_least_obj}.
For example, \citet{yabe1991factorized,eriksson1999quasi,zhou2010global} adopt Broyden family updates to approximate the Hessian for the objective of problem \eqref{eq:nonlinear_least_obj}.
However, these methods still require accessing the exact Gram matrix at the early stage of the iterations, leading to expensive per-iteration cost like standard Gauss--Newton methods.

We summarize the main limitations of the aforementioned quasi-Newton methods for nonlinear systems as follows.
\begin{itemize}[leftmargin=*,label={}]
\item \emph{Limitation 1.} Existing non-asymptotic superlinear convergence guarantees for quasi-Newton methods are limited to squared systems (i.e., the case of $n=d$), whereas how to establish the superlinear convergence for the general case is unclear.
\item \emph{Limitation 2.} Existing quasi-Newton methods for general nonlinear systems require either a sufficiently accurate initial Jacobian estimator or accessing the exact Gram matrix, resulting quite expensive computational cost.
\end{itemize}



\begin{table}[t]
  \caption{We summarize the local convergence results and initial conditions of quasi-Newton methods for solving nonlinear equations, where the dash indicates that the algorithm has no convergence guarantee under the corresponding setting.}
  \label{tbl:compare}\vskip-0.1cm
    \resizebox{\linewidth}{!}{
    \centering
    %\renewcommand{\arraystretch}{1.3}
    \begin{tabular}{cccccc} \hline\
    Method  & $n=d$ & $n<d$ & $n>d$ & Non-Asymptotic & Initial Condition \\ \hline\hline 
    \makecell{Broyden \\ 
         \footnotesize \citep{lin2021explicit} \\
         \footnotesize \citep{ye2021greedy} \\ 
         \footnotesize \citep{liu2023block}} & 
    superlinear & -- & -- & \Checkmark & $\G_0\approx\J(\x_0)$ \\\hline 
    \makecell{Generalized Broyden \\
         \footnotesize \citep{martinez1991quasi} \\
         \footnotesize \citep{vater2024convergence}
         } & 
    superlinear & superlinear & -- &  \XSolidBrush & $\G_0\approx\J(\x_0)$ \\\hline 
    \makecell{GN-BFGS \\
         \footnotesize \citep{gu2002descent} \\
         \footnotesize \citep{li1999globally}
         } &
    superlinear & -- & -- & \XSolidBrush & $\mG_0=c\mI_m$ \\ \hline
    \makecell{RQGN \\
         \footnotesize Algorithm \ref{alg:RQGN}} &
         ~superlinear~ & ~superlinear~ & ~linear~ & \Checkmark & $\mG_0=c\mI_m$ \\ \hline
    \end{tabular}} \vskip-0.1cm
  % \begin{tablenotes}
  %   \footnotesize
  %   \item[$\dag$] These methods only consider a symmetric Jacobian.
  % \end{tablenotes}
  % \end{threeparttable}
\end{table}

% \begin{table}[t]
%   \centering
%   \begin{threeparttable}
%   \small
%   \setlength{\tabcolsep}{1.5pt}
%   \renewcommand{\arraystretch}{1.12}
%   \begin{tabular*}{\textwidth}{@{\extracolsep{\fill}}ccccc@{}}
%     \toprule
%     \begin{tabular}[c]{@{}c@{}}\textbf{Method}\end{tabular}
%     & \begin{tabular}[c]{@{}c@{}}\textbf{$n=d$}\end{tabular}
%     & \begin{tabular}[c]{@{}c@{}}\textbf{$n<d$}\end{tabular}
%     & \begin{tabular}[c]{@{}c@{}}\textbf{$n>d$}\end{tabular}
%     & \begin{tabular}[c]{@{}c@{}}\textbf{Initial condition}\end{tabular} \\
%     \midrule
%     \begin{tabular}[c]{@{}c@{}}Broyden\\[-1pt]
%     {\tiny \citep{lin2021explicit,ye2021greedy}}\\[-2pt]
%     {\tiny \citep{liu2023block}}\end{tabular}
%       & \begin{tabular}[c]{@{}c@{}}\textbf{non-asymptotic}\\\textbf{superlinear}\end{tabular}
%       & --
%       & --
%       & $\G_0\approx\J(\x_0)$ \\[2pt]
%     \begin{tabular}[c]{@{}c@{}}Generalized Broyden\\[-1pt]
%     {\tiny \citep{martinez1991quasi}}\\[-2pt]
%     {\tiny \citep{vater2024convergence}}\end{tabular}
%       & \begin{tabular}[c]{@{}c@{}}\textbf{asymptotic}\\\textbf{superlinear}\end{tabular}
%       & \begin{tabular}[c]{@{}c@{}}\textbf{asymptotic}\\\textbf{superlinear}\end{tabular}
%       & --
%       & $\G_0\approx\J(\x_0)$ \\[2pt]
%     \begin{tabular}[c]{@{}c@{}}GN-BFGS\tnote{$\dag$}\\[-1pt]
%     {\tiny \citep{gu2002descent}}\\[-2pt]
%     {\tiny \citep{li1999globally}}\end{tabular}
%       & \begin{tabular}[c]{@{}c@{}}\textbf{asymptotic}\\\textbf{superlinear}\end{tabular}
%       & --
%       & --
%       & $\G_0=c\I$ \\[2pt]
%     \midrule
%     \begin{tabular}[c]{@{}c@{}}QN-SPP\\[-1pt]
%     {\tiny Conference version}\\[-2pt]
%     {\tiny \citep{liu2022quasi}}\end{tabular}
%       & \begin{tabular}[c]{@{}c@{}}\textbf{non-asymptotic}\\\textbf{superlinear}\end{tabular}
%       & --
%       & --
%       & $\G_0=c\I$ \\[2pt]
%     \begin{tabular}[c]{@{}c@{}}RQGN\\[-1pt]{\tiny Algorithm~\ref{alg:RQGN}}\end{tabular}
%       & \begin{tabular}[c]{@{}c@{}}\textbf{non-asymptotic}\\\textbf{superlinear}\end{tabular}
%       & \begin{tabular}[c]{@{}c@{}}\textbf{non-asymptotic}\\\textbf{superlinear}\end{tabular}
%       & \begin{tabular}[c]{@{}c@{}}\textbf{non-asymptotic}\\\textbf{linear}\end{tabular}
%       & $\G_0=c\I$ \\[2pt]
%     \bottomrule
%   \end{tabular*}
%   \caption{Local convergence rates and matrix-initialization requirements for quasi-Newton methods under square, underdetermined, and overdetermined nonlinear equations. A dash indicates that the cited result does not cover the corresponding setting.}
%   \label{tbl:compare}
%   \begin{tablenotes}
%     \footnotesize
%     \item[$\dag$] These methods only consider a symmetric Jacobian.
%   \end{tablenotes}
%   \end{threeparttable}
% \end{table}



In this paper, we develop randomized quasi-Gauss–Newton (RQGN) methods to address both the above two limitations.
In contrast to existing methods that directly approximate the (pseudo) inverse of the Jacobian~\citep{lin2021explicit,ye2021greedy,liu2023block,martinez1991quasi,vater2024convergence},
we approximate the Gram matrix in the Gauss--Newton iteration by Broyden family updates with randomized direction,
which leads to the non-asymptotic local convergence guarantees for both overdetermined and underdetermined cases.
We compare our results with existing quasi-Newton methods in Table~\ref{tbl:compare} and summarize our main theoretical contributions as follows.
\begin{itemize}[topsep=0.08cm,itemsep=0.08cm]
\item Our proposed RQGN methods work for both underdetermined and overdetermined nonlinear equations, and they requires neither the exact Gram matrix nor a sufficiently accurate initial Jacobian approximation so its Gram estimator can be initialized by a scaled identity matrix.
\item For the underdetermined case (i.e., $n<d$), we show that our RQGN enjoy explicit non-asymptotic local superlinear convergence rates without assuming uniqueness of the solution.
\item For the overdetermined case (i.e., $n\geq d$), we establish explicit local condition-number-free linear convergence rates for finding stationary points of the nonlinear least square formulation \eqref{eq:nonlinear_least_obj}, and our results hold even if the system \eqref{eq:nonlinear_obj} has no solution.
\end{itemize}

\paragraph{Paper Organization}
In Section~\ref{sec:pre}, we formalize the notations and assumptions throughout this paper and introduce the background of quasi-Newton updates. 
In Section~\ref{sec:rqn}, we propose our randomized quasi-Gauss--Newton (RQGN) method for solving the general nonlinear equations.
In Sections~\ref{sec:underde} and~\ref{sec:overfit}, we establish the non-asymptotic local convergence guarantees of our RQGN for both underdetermined and overdetermined cases, respectively.
% In Section \ref{sec:conference}, we discuss the improvements of our paper over its conference version \citep{liu2022quasi}. 
In Section~\ref{sec:exp}, we conduct numerical experiments to show the superiority of the proposed methods.
We conclude our work in Section \ref{sec:conclu}.

% Appendix~\ref{app:proofs} contains the supplementary proofs, Appendix~\ref{app:fbfgs} gives the F-BFGS implementation, and Appendix~\ref{app:additional_experiments} reports additional experiments.


\section{Notation and Preliminaries} 
\label{sec:pre}

We first introduce the notations used in this paper. We use $\Norm{\,\cdot\,}$ to present the spectral norm and the Euclidean norm of the matrix and the vector, respectively. 
Given a matrix $\S$, we write $\S^{\dag}$ for its Moore--Penrose inverse and $\sigma_{i}(\S)$ for its $i$th largest singular value; when $\S$ is square, 
we use~$\tr{\S}$ to denote its trace.
For two positive definite matrices $\G$ and $\H$ with the same size, we write
$\langle\G,\H\rangle\triangleq\text{tr}(\G\H)$ for their inner product.
Additionally, we denote the Jacobian of the mapping $\F:\BR^d\to\BR^n$ as $\J:\BR^d\to\BR^{n\times d}$ and define the Gram matrix~as
\begin{align}\label{eq:H-Gram}
    \H(\x)\triangleq
\begin{cases}
\J(\x)\J(\x)^{\top}, & \text{if}~~n<d, \\    
\J(\x)^{\top}\J(\x), & \text{if}~~n\geq d.
\end{cases}
\end{align}
%For $n<d$, define the Gram matrix by $\H(\x)\triangleq\J(\x)\J(\x)^{\top}$; for $n\geq d$, define it by $\H(\x)\triangleq\J(\x)^{\top}\J(\x)$. 
For the ease of presentation, we let $m\triangleq\min\{n,d\}$ so that we can write $\H(\x)\in\RB^{m\times m}$.
Accordingly, we denote the standard basis for the Euclidean space $\BR^m$ by $\{\ve_1,\dots,\ve_m\}$ and let $\I_{m}$ be the identity matrix with the size of $m\times m$.

We then formalized the assumptions of our problems (\ref{eq:nonlinear_obj}) and (\ref{eq:nonlinear_least_obj}). 
We assume the nonlinear mapping $\F:\BR^d\to\BR^n$ 
has the bounded and Lipschitz continuous Jacobian $\mJ:\BR^d\to\BR^{n\times d}$.
\begin{assumption}
\label{ass:Jlip}
    We assume that there exist constants $L>0$ and $L_2>0$ such that
    \begin{align}
    \label{eq:lip}
    \|\J(\x)\|\leq L\qquad\text{and}\qquad\| \J(\x)-\J(\y)\|\leq L_2\|\x-\y\| 
    \end{align}
    for all $\vx,\vy\in\BR^d$. 
\end{assumption}
We also assume that the Jacobian Jacobian $\mJ:\BR^d\to\BR^{n\times d}$ is always nonsingular for some non-empty set.
\begin{assumption}
\label{ass:nonsingular}
We assume that there exist some non-empty set $\Omega\subseteq\RB^d$ 
% such that $\x^*\in\Omega$
% $\Omega(\delta) = \{\x~~\text{s.t.}~~~\|\x-\x^*\|\leq \delta \}$ where $\delta>0$ and $\x^*$ is the solution such that $\F(\x^*)=\0$ for $n<d$ and the solution such that $\J(\x^*)^{\top}\F(\x^*)=\0$ for $n\geq d$, 
and a constant $\mu>0$ such that $\sigma_m(\J(\x))\geq \mu$ holds for all $\x\in\Omega$, where $m=\min\{n,d\}$.
%${\rm rank}(\J(\x))=m$.
%In addition, we assume that $\sigma_m(\J(\x)^{\top})\geq \mu>0$ if $n<d$ and $\sigma_m(\J(\x))\geq \mu>0$ if $n\geq d$.
\end{assumption}
% {\textcolor{red}{
% \begin{remark}
% In the convergence analysis, we will explicitly define the set in Assumption~\ref{ass:nonsingular}.    
% \end{remark}
% }}

The above assumptions indicate the following properties for the nonlinear mapping $\F(\cdot)$ and the Gram matrix $\H(\cdot)$.
\begin{proposition}
\label{prop:gram_lip}
    Under Assumption~\ref{ass:Jlip}, the nonlinear mapping $\F(\cdot)$ is $L$-Lipschitz continuous and the Gram matrix $\H(\cdot)$ is $2LL_2$-Lipschitz continuous with $\H(\x)\preceq L^2\I_m$ for all $\x\in\RB^{d}$.
\end{proposition}
% \begin{proof}
% The $L$-Lipschitz continuity of $\F(\cdot)$ follows directly from equation \eqref{eq:lip}.
% For the Lipschitz continuity of $\H(\cdot)$, we focus on the case $n\geq d$. 
% Under Assumption~\ref{ass:Jlip}, it holds
% \begin{align*}
% \|\H(\x)\|
% &=\big\|\J(\x)^{\top}\J(\x)\big\|
% \leq \big\|\J(\x)^{\top}\big\|\|\J(\x)\|
% \leq L^2,
% \end{align*}
% which means $\H(\x)\preceq L^2\I$.
% Besides, we also have
% \begin{align*}
%     & \|\H(\x)-\H(\y)\|
% \leq \big\|\J(\x)^{\top}\J(\x) - \J(\x)^{\top}\J(\y) + \J(\x)^{\top}\J(\y) - \J(\y)^{\top}\J(\y)\big\|  \\
% \leq & \big\|\J(\x)^{\top}\big\|\Norm{\J(\x)-\J(\y)}+\big\|\J(\x)^{\top}-\J(\y)^{\top}\big\|\Norm{\J(\y)}\leq 2LL_2\|\x-\y\|.
% \end{align*}
% The proof for the case $n<d$ is almost identical to the case $n\geq d$, so we omit it here.
% \end{proof}
\begin{proposition}
\label{prop:gram_lower}
    Under Assumption~\ref{ass:nonsingular}, it holds  $\H(\x)\succeq \mu^2\I_m$ for all $\x\in\Omega$.
\end{proposition}
The Gram matrix $\mH(\cdot)$ also holds the following partial relation.
\begin{proposition}
\label{prop:Gram_concor}
    Under Assumptions \ref{ass:Jlip} and \ref{ass:nonsingular}, it holds 
    \begin{align}
    \label{eq:Gram_bound_0}
        \frac{\H(\x)}{1+M\|\y-\x\|}\preceq\H(\y) \preceq (1+M\|\y-\x\|)\H(\x).
    \end{align}
    for all $\x,\y\in\Omega$, where  $M\triangleq2LL_2/\mu^2$.    
\end{proposition}
% \begin{proof}
% According to Proposition~\ref{prop:gram_lip}, it holds
%     \begin{align*}
%     -2L_2L \|\x-\y\|\cdot\I \preceq   \H(\y)-\H(\x)\preceq 2L_2L \|\x-\y\|\cdot\I.
%     \end{align*}
% Combining above results with Proposition~\ref{prop:gram_lower}, we obtain
%     \begin{align*}
%      -\frac{2L_2L}{\mu^2}\|
%     \x-\y\| \H(\y)\preceq   \H(\y)-\H(\x)\preceq \frac{2L_2L}{\mu^2}\|\x-\y\|\H(\x).
%     \end{align*}
% Rearranging above inequalities directly indicates equation \eqref{eq:Gram_bound_0}.
% \end{proof}
% \begin{proof}
% We consider the case $n\geq d$.
% Under Assumption~\ref{ass:nonsingular}, we have $\sigma_m(\H(\x))=\sigma_m^{2}(\J(\x))\geq\mu^2$, which means $\mu^2\I\preceq\H(\x)$. 
% The proof for the case $n<d$ is almost identical to the case $n\geq d$, so we omit it here.
% \end{proof}

We now introduce some popular quasi-Newton updates to approximate the target matrix~$\H\in\RB^{m\times m}$ by 
the current estimator $\G\in\RB^{m\times m}$ and the update direction $\u\in\RB^{m}$ in the following definition.

%Let $\H\in\RB^{m\times m}$ be the target matrix that is positive definite, $\G\in\RB^{m\times m}$ be an estimator of $\H$, and $\u\in\RB^{m}$ be the update direction. 
\begin{definition}[{\citet{lin2022explicit,rodomanov2021greedy}}]
\label{def:qnupdate}
Given two positive-definite matrices $\G,\H\in\RB^{m\times m}$ and vector $\u\in\RB^m$, we define the following quasi-Newton updates.
\begin{itemize}
\item The BFGS update is defined by
\begin{align*}
\bfgs(\G,\H,\u)
&\triangleq \G-\frac{\G\u\u^{\top}\G}{\u^{\top}\G\u}
 +\frac{\H\u\u^{\top}\H}{\u^{\top}\H\u}.
\end{align*}
\item The fast BFGS update is defined by
\begin{align*}
\fbfgs(\G,\H,\u)
&\triangleq\bfgs\big(\G,\H,\L^{\top}\u\big),
\end{align*}
where the matrix $\L\in\RB^{m\times m}$ satisfies $\G^{-1}=\L^{\top}\L$.
\item The SR1 update is defined by
\begin{align*}
\srk(\G,\H,\u)
&\triangleq\G-\frac{(\G-\H)\u\u^{\top}(\G-\H)}
{\u^{\top}(\G-\H)\u}.
\end{align*}
\end{itemize}
\end{definition}
% \begin{remark}
% For $k=1$, these updates are exactly the same as BFGS, fast BFGS, and SR1 updates in \citet{lin2022explicit}.
% For $k>1$, these updates are introduced in \citet{gower2017randomized,liu2023symmetric}. 
% \end{remark}
\begin{remark}
The fast BFGS update requires to maintain $\L_t$ such that $\G_t^{-1}=\L_t^{\top}\L_t$, which can be implemented by a closed form update with the complexity of $\OM(m^2)$ flops per iteration.
Please refer to Algorithm~\ref{alg:fbfgs} in Appendix~\ref{app:fbfgs}.
\end{remark}

We describe the convergence of above quasi-Newton updates by introducing the quantities
\begin{align*}
    \sigma_{\H}(\G) \triangleq \tr{\H^{-1}(\G-\H)}\qquad\text{and}\qquad\tau_{\H}(\G)\triangleq\tr{\G-\H}
\end{align*}
to describe the difference between positive-definite matrices $\G\in\RB^{m\times m}$ and $\H\in\RB^{m\times m}$. 
By taking random vector $\u\in\RB^{m}$ according to Gaussian distribution, it holds the following in-expectation convergence guarantees \citep{lin2022explicit}.
%We use the following quantities to show the difference between two positive definite matrices


\begin{lemma}[{\citet[Theorems~6, 13, and~14]{lin2022explicit}}]
\label{lm:QN_one_iter}
    Let $\H,\G\in\RB^{m\times m}$ be positive definite with $\H\preceq\G$. We suppose $\mu^2\I_m\preceq\H\preceq L^2\I_m$ for some constants $L,\mu>0$ such that~$L\geq\mu$. By taking $\vu\sim\fN(\vzero,\mI_m)$, the quasi-Newton updates has the following properties.
\begin{itemize}
\item For $\G_{+}=\bfgs(\G,\H,\u)$, it holds  $\EBP{\sigma_{\H}(\G_{+})}\leq(1-\mu^2/(mL^2))\sigma_{\H}(\G)$;
\item For $\G_{+}=\fbfgs(\G,\H,\u)$, it holds  $\EBP{\sigma_{\H}(\G_{+})}\leq(1-1/m)\sigma_{\H}(\G)$;
\item For $\G_{+}=\srk(\G,\H,\u)$, it holds  $\EBP{\tau_{\H}(\G_{+})}\leq(1-1/m)\tau_{\H}(\G)$.
\end{itemize}
\end{lemma}

\begin{lemma}[{\citet[Lemma~5]{lin2022explicit}}]
\label{lm:QN_order}
    Let $\H,\G\in\RB^{m\times m}$ be positive definite with $\H\preceq\G\preceq\eta\H$ for some $\eta\geq1$. If the matrix $\G_{+}$ is obtained by either BFGS update, fast BFGS update, or SR1 update by following Definition~\ref{def:qnupdate}, it holds $\H\preceq\G_{+}\preceq\eta\H$.
\end{lemma}



% \begin{remark}
%     For $k=1$, the results can be found in \citet{rodomanov2021greedy} and \citet{lin2022explicit}.
% \end{remark}

\begin{algorithm}[t]
\caption{Randomized Quasi-Gauss--Newton Methods (RQGN)}\label{alg:RQGN}
\begin{algorithmic}[1]
\STATE \textbf{Input:} $\x_0\in\Omega$, $\G_0\succeq\H_0\in\RB^{m\times m}$, $M\geq0$ \label{line:input} \\[0.15cm]
\STATE \textbf{for} $t=0,1,\ldots$ \\[0.15cm]
 		\STATE \quad
 		    $\x_{t+1}=\begin{cases}
               \x_t-\J(\x_t)^{\top}\G_t^{-1}\F(\x_t), &\text{if }n<d,\\
	 	    \x_t-\G_t^{-1}\J(\x_t)^{\top}\F(\x_t), &\text{if }n\geq d
	 	    \end{cases}$ \label{eq:line-update-x} \\[0.15cm]
	 	\STATE \quad $r_t=\|\x_{t+1}-\x_{t}\|$ \\[0.15cm]
 		\STATE \quad $\widetilde{\G}_{t}=(1+Mr_t)\G_t$ \\[0.15cm]
        \STATE \quad $\u_t\sim\fN(\0,\I_m)$ \\[0.15cm]
 		\STATE \quad [Option I]~~ $\G_{t+1}=\bfgs(\widetilde{\G}_t,\H(\x_{t+1}),\u_t)$ \label{eq:line-update-BFGS} \\[0.15cm]
 		\STATE \quad [Option II]~ $\G_{t+1}=\fbfgs(\widetilde{\G}_t,\H(\x_{t+1}),\u_t)$ \label{eq:line-update-F-BFGS} \\[0.15cm]
 		\STATE \quad [Option III] $\G_{t+1}=\srk(\widetilde{\G}_t,\H(\x_{t+1}),\u_t)$ \label{eq:line-update-SR1} \\[0.15cm]
\STATE \textbf{end for}
\end{algorithmic}
\end{algorithm}

\section{Randomized Quasi-Gauss--Newton Methods}
\label{sec:rqn}
In this section, we introduce randomized quasi-Gauss--Newton (RQGN) methods for solving general nonlinear equations.
The key idea of our algorithm design is constructing the square matrix~$\G$ to approximate the Gram matrix $\mJ(\vx)\mJ(\vx)^\top$ or $\mJ(\vx)^\top\mJ(\vx)$ in Gauss--Newton formulations~\eqref{update:gn} and \eqref{update:gnJ}, which leads to the iteration scheme
\begin{align}
\label{eq:qNupdate}
    \x_{+}=\begin{cases}
      \x-\J(\x)^{\top}\G^{-1}\F(\x), &\text{if }n<d,\\
      \x-\G^{-1}\J(\x)^{\top}\F(\x), &\text{if }n\geq d.
    \end{cases}
\end{align}
Furthermore, we apply quasi-Newton updates introduced in Definition \ref{def:qnupdate} to construct the square matrix~$\mG^+\in\BR^{m\times m}$ to approximate the Gram matrix at next point $\vx_+\in\BR^m$ by given the current estimator $\G\in\BR^{m\times m}$, i.e.,
\begin{align}
\label{eq:update_G}
\G_{+}
=\begin{cases}
\bfgs(\widetilde{\G},\H(\x_{+}),\u), & \text{Option I},\\
\fbfgs(\widetilde{\G},\H(\x_{+}),\u), & \text{Option II},\\
\srk(\widetilde{\G},\H(\x_{+}),\u), & \text{Option III},
\end{cases}
\end{align}
where $\widetilde{\G}\triangleq (1+M\|\x_{+}-\x\|)\G$, $\u\sim\fN(\0,\I_m)$, $M>0$, and $\mH(\cdot)$ follows equation \eqref{eq:H-Gram}.
Combining the iteration scheme~\eqref{eq:qNupdate} and the construction~\eqref{eq:update_G}, we formally present our randomized quasi-Gauss--Newton methods in Algorithm~\ref{alg:RQGN}.

Note that the initialization of RQGN (Line \ref{line:input} of Algorithm \ref{alg:RQGN}) only requires that the initial Gram estimator satisfies $\G_0\succeq\H_0$. Therefore, Proposition \ref{prop:gram_lip} indicates that we can simply set $\mG_0=c\mI_m$ for some $c\geq L^2$ and maintain the square matrix $\mG_t$ (or its inverse) with the complexity of $\OM(m^2)$ flops per iteration, since the quasi-Newton updates introduced in Definition \ref{def:qnupdate} only involves rank-one or rank-two modification.
Furthermore, we emphasize that the updates of $\vx_{t+1}$ and $\mG_{t+1}$ (Lines \ref{eq:line-update-x} and \ref{eq:line-update-BFGS}--\ref{eq:line-update-SR1}) can be efficiently implemented by accessing the Jacobian-vector product so that it is unnecessary to explicitly compute the Jacobian $\mJ(\cdot)$.




Additionally, our RQGN methods with appropriate scaled parameter $M$ can guarantee the partial relation 
\begin{align}
\label{eq:scale_G00}
    {\tilde \G}_t \succeq \H(\x_{t+1})
\end{align}
holds for all $t\geq 0$ (see Appendix \ref{app:proof_lemma_gram_approximate}), which is crucial to our later convergence analysis .
In contrast, simply applying quasi-Newton updates on matrix $\mG_t$ without scaling cannot achieve~\eqref{eq:scale_G00} for general nonlinear equations.


% {\color{red}This means we can apply quasi-Newton updates as introduced in Section~\ref{sec:pre} on $\widetilde{\G}$ instead of $\G$ to obtain an estimator for $\H(\x_{+})$:
% % However, $\G$ may not satisfy the partial order relation that $\G\succeq \H(\x_{+})$ even if we assume $\G\succeq \H(\x)$.
% Following the notation \eqref{eq:H-Gram}, we establish the partial order relation for the Gram matrix as follows. 



% apply similar strategies as in minimizing strongly self-concordant functions~\citep{rodomanov2021greedy}, that is, we always multiply a scale factor on $\G$ before doing quasi-Newton updates.
% The following properties of $\H(\x)$ analogize the strongly self-concordance in convex optimization.
% When $\G\succeq\H(\x)$, Proposition~\ref{prop:Gram_concor} suggests choosing the scale factor $(1+M\|\x_{+}-\x\|)$ in the quasi-Newton updates, so that we have
% \begin{align}
% \label{eq:scale_G}
%     \widetilde{\G}\triangleq (1+M\|\x_{+}-\x\|)\G\succeq (1+M\|\x_{+}-\x\|)\H(\x)\overset{\eqref{eq:Gram_bound_0}}{\succeq} \H(\x_{+}).
% \end{align}
% This means we can apply quasi-Newton updates as introduced in Section~\ref{sec:pre} on $\widetilde{\G}$ instead of $\G$ to obtain an estimator for $\H(\x_{+})$:}




% {\color{red}
% The following lemma incorporates Lemmas~\ref{lm:QN_one_iter} and~\ref{lm:QN_order} and Proposition~\ref{prop:Gram_concor} to show how well $\G_{t}$ approximates $\H(\x_{t})$.
% Let $\mathcal F_t=\sigma(\u_0,\ldots,\u_{t-1})$ denote the history before the random update at iteration $t$, and write $\mathbb E_t[\,\cdot\,]=\mathbb E[\,\cdot\mid\mathcal F_t]$.
% \begin{lemma}[{Modified from \citet[Lemma 25]{lin2022explicit}}]
%     \label{lm:Gram_approximate}
%     Under Assumptions~\ref{ass:Jlip} and \ref{ass:nonsingular}, run Algorithm~\ref{alg:RQGN} with independent random directions $\u_t\in\RB^m$. Assume $\x_t\in\Omega$ for every $t$ and $\H(\x_0)\preceq\G_0\preceq\eta_0\H(\x_0)$. Let $r_t\triangleq\|\x_{t+1}-\x_t\|$. For each update option, there is a random sequence $\delta_t$ such that $\H(\x_t)\preceq\G_t\preceq(1+\delta_t)\H(\x_t)$ and
%     \begin{align}
%     \label{eq:delta-recur}
%         \mathbb E_t[\delta_{t+1}]\leq (1-\gamma)(1+Mr_t)^2(\delta_t+2\tilde{c}mMr_t)
%     \end{align}
%     for every $t$, with the following parameters.
%     \begin{itemize}
%     \item Run Algorithm~\ref{alg:RQGN} with the BFGS update (Option I). Then
%     \begin{align*}
%     \delta_t&=\sigma_{\H(\x_t)}(\G_t),&
%     \gamma&=\mu^2/(mL^2),&
%     \tilde c&=1,&
%     \delta_0&\leq m(\eta_0-1).
%     \end{align*}
%     \item Run Algorithm~\ref{alg:RQGN} with the fast BFGS update (Option II). Then
%     \begin{align*}
%     \delta_t&=\sigma_{\H(\x_t)}(\G_t),&
%     \gamma&=1/m,&
%     \tilde c&=1,&
%     \delta_0&\leq m(\eta_0-1).
%     \end{align*}
%     \item Run Algorithm~\ref{alg:RQGN} with the SR1 update (Option III). Then
%     \begin{align*}
%     \delta_t&=\frac{mL^2\tau_{\H(\x_t)}(\G_t)}
%     {\mu^2\tr{\H(\x_t)}},&
%     \gamma&=1/m,&
%     \tilde c&=L^2/\mu^2,&
%     \delta_0&\leq mL^2(\eta_0-1)/\mu^2.
%     \end{align*}
%     \end{itemize}
% \end{lemma}
% \begin{proof}
%     We first prove that $\G_t\succeq\H(\x_t)$ for every $t$. Since $\G_0\succeq\H(\x_0)$, equation~\eqref{eq:scale_G} gives $\widetilde{\G}_0\succeq\H(\x_1)$. For each of Options I--III, Lemma~\ref{lm:QN_order} therefore gives $\G_1\succeq\H(\x_1)$. Repeating this argument yields $\G_t\succeq\H(\x_t)$ for every $t$.
%     The remaining proof is modified from \citet[Lemma 25]{lin2022explicit} and is presented in Appendix~\ref{app:proof_lemma_gram_approximate} for completeness.
% \end{proof}
% }

\section{The Convergence Analysis}\label{sec:local_convergence}

We analyze the local convergence behavior of the randomized quasi-Gauss--Newton method (Algorithm~\ref{alg:RQGN}) for both the underdetermined case and the overdetermined cases in Section~\ref{sec:underde} and~\ref{sec:overfit}, respectively.

For the ease of presentation, we introduce the notations
\begin{align*}
r_t\triangleq \|\x_{t+1}-\x_t\| 
\qquad\text{and}\qquad \H_t\triangleq\H(\x_t).    
\end{align*}
In addition, the initial condition $\x_0\in\Omega$ of Algorithm~\ref{alg:RQGN} and Assumption \ref{ass:nonsingular} indicate $\mH_0\succ \vzero$, so that we can define $\eta_0\in\BR$ as the smallest scalar satisfying $\G_0\preceq\eta_0\H_0$, i.e.,
\begin{align}\label{eq:dfn-eta0}
    \eta_0\triangleq \min_{\{\eta: 
 \H_0\preceq\G_0\preceq \eta\H_0\}}\eta.
\end{align}
Furthermore, we let $\mathcal F_t\triangleq\sigma(\u_0,\ldots,\u_{t-1})$ be the 
sigma filed generated by the random variable sequence $\{\vu_0,\dots,\vu_{t-1}\}$ and define $\mathbb E_t[\,\cdot\,]\triangleq\mathbb E[\,\cdot\mid\mathcal F_t]$.

For our later analysis, we first incorporate Lemmas~\ref{lm:QN_one_iter} and~\ref{lm:QN_order} and Proposition~\ref{prop:Gram_concor} to show how well the square matrix~$\G_{t}$ approximates the Gram matrix $\H_t$ in the following lemma, which extends Lemma 25 of \citet{lin2022explicit}.

% We use the following notation to simplify the presentation: $r_t\triangleq \|\x_{t+1}-\x_t\|$ and $\H_t\triangleq\H(\x_t)$.
% The following lemma incorporates Lemmas~\ref{lm:QN_one_iter} and~\ref{lm:QN_order} and Proposition~\ref{prop:Gram_concor} to show how well $\G_{t}$ approximates $\H(\x_{t})$.
\begin{lemma}
    \label{lm:Gram_approximate}
    Under Assumptions \ref{ass:Jlip} and \ref{ass:nonsingular}, we suppose points $\{\vx_t\}$ generated by Algorithm~\ref{alg:RQGN} holds $\x_t\in\Omega$ for all $t$. Then there exists non-negative randomized sequence $\{\delta_t\}$ such that %$\H_t\preceq\G_t\preceq(1+\delta_t)\H_t$ and
    \begin{align}
    \label{eq:delta-recur}
        \H_t\preceq\G_t\preceq(1+\delta_t)\H_t
        \qquad\text{and}\qquad
        \mathbb E_t[\delta_{t+1}]\leq (1-\gamma)(1+Mr_t)^2(\delta_t+2\tilde{c}mMr_t)
    \end{align}
    for all $t$, where $\delta_t$, $\gamma$, and $\tilde c$ are determined as follows.
    \begin{itemize}
    \item For Algorithm~\ref{alg:RQGN} with the BFGS update (Option I), the result \eqref{eq:delta-recur} holds by taking
    \begin{align*}
    \delta_t=\sigma_{\H_t}(\G_t),\quad
    \gamma=\frac{\mu^2}{mL^2},\quad
    \tilde c=1,
    \quad\text{and}\quad\delta_0\leq m(\eta_0-1).
    \end{align*}
    \item For Algorithm~\ref{alg:RQGN} with the fast BFGS update (Option II), the result \eqref{eq:delta-recur} holds by taking
    \begin{align*}
    \delta_t=\sigma_{\H_t}(\G_t),\quad
    \gamma=\frac{1}{m},\quad
    \tilde c=1,
    \quad\text{and}\quad\delta_0\leq m(\eta_0-1).
    \end{align*}
    \item For Algorithm~\ref{alg:RQGN} with the SR1 update (Option III), the result \eqref{eq:delta-recur} holds by taking
    \begin{align*}
    \delta_t=\frac{mL^2\tau_{\H_t}(\G_t)}
    {\mu^2\tr{\H_t}},\quad
    \gamma=\frac{1}{m},\quad
    \tilde c=\frac{L^2}{\mu^2},
    \quad\text{and}\quad
    \delta_0\leq \frac{mL^2(\eta_0-1)}{\mu^2}.
    \end{align*}
    \end{itemize}
\end{lemma}

\begin{remark}
The results of the above Lemma \ref{lm:Gram_approximate} require the condition $\x_t\in \Omega$ holds for all $t$, 
which can be guaranteed by taking appropriate $\Omega$ for both underdetermined and overdetermined settings.
Please refer to Theorems \ref{thm:linear_under} and \ref{thm:linear1n>d}.    
\end{remark}


The results of Lemma~\ref{lm:Gram_approximate} suggest  $\H_t\preceq\G_t\preceq(1+\delta_t)\H_t$ holds for some  $\delta_t\geq 0$.
Hence, we can extend the definition of $\eta_0$ to every iteration by denoting
\begin{align}
\label{eq:def_eta}
    \eta_t\triangleq \min_{\{\eta: 
 \H_t\preceq\G_t\preceq \eta\H_t\}}\eta.
\end{align}


\subsection{The Local Convergence for the Underdetermined Case}
\label{sec:underde}
In this section, we analyze the convergence of our RQGN method for solving the nonlinear equations  \eqref{eq:nonlinear_obj} in the case of $n<d$.
% In this section, we analyze the convergence of our RQGN method for solving underdetermined nonlinear equations (i.e., the case of $n<d$), where the solution of nonlinear equations~\eqref{eq:nonlinear_obj} exists but may not be unique.
We suppose the singular value of the Jacobian is lower bounded on a local residual region.
% We make the following assumption to guarantee the existence of the solution $\x^*$ of \eqref{eq:nonlinear_obj}.
\begin{assumption}
\label{ass:Omegan<d}
    We assume that Assumption~\ref{ass:nonsingular} holds with
    \begin{align*}
    \Omega
    &=\Omega_{\mathrm{under}}
    \triangleq\left\{\x:\|\F(\x)\|\leq\frac{\mu^2}{2L_2}\right\},
    \end{align*}
    i.e., there exists a constant $\mu>0$ such that $\sigma_m(\J(\x))\geq \mu$ holds for all $\x\in\Omega_{\mathrm{under}}$.
\end{assumption}
\begin{remark}
   In the global analysis of a modified Gauss--Newton method, \citet{nesterov2007modified} require the singular value of the Jacobian is lower bounded for all $\x\in\RB^d$.
   Here, our Assumption~\ref{ass:Omegan<d} only requires the lower bound holds on the region $\Omega_{\mathrm{under}}$ since we focus on the local convergence analysis.
\end{remark}

%We target to make the norm of $\F(\cdot)$ be close to zero. 

Recall that the classical analysis of quasi-Newton methods for minimizing the strongly convex function $f:\BR^d\to\BR$ \citep{rodomanov2021greedy} describes the convergence rates by the scaled gradient norm 
\begin{align*}
\lambda(\vx)\triangleq\sqrt{\nabla f(\x)^{\top}(\nabla^2f(\x))^{-1}\nabla f(\x)}     
\end{align*}
where $\nabla f(\cdot)$ and $\nabla^2 f(\cdot)$ are the gradient and the Hessian of $f(\cdot)$, respectively.
A direct extension of the above metric to nonlinear linear equations in the view of taking~$\mF(\cdot)=\nabla f(\cdot)$ and $\mJ(\cdot)=\nabla^2 f(\vx)$ is
\begin{align}
\label{eq:ordi-measure}
    \hat{\lambda}(\x) \triangleq \sqrt{\F(\x)^{\top}(\J(\x))^{-1}\F(\x)}.
\end{align}
However, for the underdetermined nonlinear equations, the Jacobian $\J(\cdot)$ may non-invertible or even non-square. 

For nonlinear equation problems,
it is worth noting that Proposition~\ref{prop:gram_lower} indicates the Gram matrix $\H(\cdot)$ is positive definite on the local region.
This motivates us to construct the measure of the weighted norm with respect to $\H(\cdot)$, i.e.,
\begin{align*}
    \lambda_{\mathrm{under}}(\x) \triangleq \sqrt{\F(\x)^{\top}(\H(\x))^{-1}\F(\x)}.
\end{align*}
Based on the above metric, we establish the following lemma shows the convergence for the general iteration \eqref{eq:qNupdate} in the underdetermined case.
\begin{lemma}
\label{lm:linear-quadra_n<d}
Under Assumptions~\ref{ass:Jlip} and \ref{ass:Omegan<d}, we additionally suppose the iteration \eqref{eq:qNupdate} in the case of $n<d$ holds
  \begin{align}
  \label{eq:G_condi_2}
      \H(\x)\preceq \G\preceq\eta\H(\x)
      \qquad\text{and}\qquad
      \lambda_{\mathrm{under}}(\x)\leq \frac{\mu^2}{4LL_2}
  \end{align}
for some $\eta\geq1$, then it holds $\x,\x_{+}\in\Omega_{\mathrm{under}}$, $\|\x_{+}-\x\|\leq\lambda_{\mathrm{under}}(\x)$, and
  \begin{align}
  \label{eq:linear-quadra-n<d}
      \lambda_{\mathrm{under}}(\x_{+})\leq \left(1-\frac{1}{\eta}\right)\lambda_{\mathrm{under}}(\x) + \frac{3LL_2}{\mu^2}(\lambda_{\mathrm{under}}(\x))^2.
  \end{align}
\end{lemma}

We now combine results of Lemma~\ref{lm:linear-quadra_n<d}, the partial relation in Proposition~\ref{prop:Gram_concor}, and the partial order-preserving property in Lemma~\ref{lm:QN_order} to establish the following explicit linear convergence rate for our methods.
%while showing that all iterates remain in $\Omega_{\mathrm{under}}$.
% The linear quadratic rate indicates the following linear rate of Algorithm~\ref{alg:RQGN} when $\lambda_{\mathrm{under}}(\x_0)$ is small.
\begin{theorem}
\label{thm:linear_under}
    Under Assumptions~\ref{ass:Jlip} and \ref{ass:Omegan<d}, we run Algorithm~\ref{alg:RQGN} with the initial point $\x_0$ such that
    \begin{align}
        \label{eq:linear_initial_condi_2}
           \lambda_{\mathrm{under}}(\x_0)\leq \frac{\mu^2}{20\eta_0LL_2}.
    \end{align}
    Then it holds $\x_t\in\Omega_{\mathrm{under}}$ for all $t$. Moreover, we also have
    \begin{align}
\label{eq:induc_lambda_n<d}
   \H_t \preceq \G_t\preceq \eta_0\exp{\left(2\sum_{i=0}^{t-1}Mr_i\right)}\H_t \preceq \frac{3\eta_0}{2}\H_t~\text{and}~   \lambda_{\mathrm{under}}(\x_t)\leq \left(1-\frac{1}{2\eta_0}\right)^t\lambda_{\mathrm{under}}(\x_0).
\end{align}
% where we define $\lambda_{\mathrm{under},t} \triangleq \lambda_{\mathrm{under}}(\x_t)$.
\end{theorem}
\begin{remark}
Unlike existing quasi-Newton methods for solving underdetermined case that require the initial Jacobian estimator to be sufficiently close to the exact Jacobian~\citep{vater2024convergence,martinez1991quasi},
the initial condition $\G_0\succeq\H_0$ in Algorithm~\ref{alg:RQGN} can be easily satisfied by taking $\G_0=L^2\I_n$, which leads to $\eta_0\leq L^2/\mu^2$.  
\end{remark}
Theorem~\ref{thm:linear_under} guarantees $\x_t\in\Omega_{\mathrm{under}}$ holds for all $t$, verifying the trajectory condition required by Lemma~\ref{lm:Gram_approximate} with $\Omega=\Omega_{\mathrm{under}}$. 
Hence, the sequence $\{\eta_t\}$ follows equation \eqref{eq:def_eta} is well-defined along the iterates.
Combining Lemma~\ref{lm:Gram_approximate} and Theorem~\ref{thm:linear_under}, we can show that that the sequence $\{\eta_t\}$ converges to $1$ in expectation and Algorithm~\ref{alg:RQGN} has the following explicit local superlinear convergence for the underdetermined case.

\begin{theorem}
\label{thm:n<dsuperlinear}
    Under Assumptions~\ref{ass:Jlip} and \ref{ass:Omegan<d}, run Algorithm~\ref{alg:RQGN} with $\x_0\in\Omega_{\rm under}$ satisfying
    \begin{align}
    \label{eq:initial_condi_n<d}
        \lambda_{\mathrm{under}}(\x_0)\leq \frac{\mu^2(1-\gamma)}{20\tilde{c}nLL_2\eta_0},
    \end{align}
where the values $\gamma$ and $\tilde{c}$ for each of Options I--III follow Lemma~\ref{lm:Gram_approximate}.
Then we have
    \begin{align}
    \label{eq:expeta_tlambda_t}
        \EBP{\eta_t-1}\leq \left(1-\gamma\right)^tq_{\text{\rm under}}\qquad\text{and}\qquad\EBP{\frac{\lambda_{\mathrm{under}}(\x_{t+1})}{\lambda_{\mathrm{under}}(\x_t)}}\leq (1-\gamma)^tq_{\text{\rm under}},
    \end{align}
where $q_{\text{\rm under}}\triangleq {\rm e}^{11/10}(\delta_0+7(1-\gamma)/(20\eta_0))$
and the value of $\delta_0$ for each of Options I--III follows Lemma~\ref{lm:Gram_approximate}.
Moreover, for all $p\in(0,1)$, there exists $t_0=\OM(\log(1+q_{\text{\rm under}}/(p\gamma^2))/\gamma)$ such that it holds
    \begin{align*}
        \lambda_{\mathrm{under}}(\x_{t_0+t})\leq \left(1-\frac{\gamma}{\gamma+1}\right)^{\frac{t(t-1)}{2}}\left(\frac{1}{2}\right)^{t}\left(1-\frac{1}{2\eta_0}\right)^{t_0}\lambda_{\mathrm{under}}(\x_0)
    \end{align*}
with probability at least $1-p$ for all $t\geq1$.
\end{theorem}
% Theorem~\ref{thm:n<dsuperlinear} presents the explicit local superlinear rates to solve nonlinear equations.
% The initial condition of $\G_0$ in \eqref{eq:initial_condi_n<d} can also be easily satisfied by letting $\G_0 = L^2\I_{n}$.


\subsection{The Local Convergence for the Overdetermined Case}
\label{sec:overfit}

This section considers the overdetermined case of $n\geq d$ by focusing on the nonlinear least square formulation \eqref{eq:nonlinear_least_obj}. 
Note that the gradient and the Hessian of the objective
$\phi(\x)=(1/2)\|\F(\x)\|^2$ can be written as
\begin{align*}
\nabla \phi(\vx)=\mJ(\vx)^\top\mF(\vx)
\qquad\text{and}\qquad
\nabla^2\phi(\x) =\H(\x) + \Q(\x),
\end{align*}
where $\mH(\vx)=\mJ(\vx)^\top\mJ(\vx)$ and $\Q(\x)\triangleq \sum_{i=1}^n F_i(\x)\nabla^2 F_i(\x)$.
We suppose that the existence of the stationary point  $\x_*\in\BR^d$ such that 
$\nabla\phi(\vx_*)=\mJ(\x_*)^{\top}\F(\x_*)=\0$
and the lower boundedness for the singular value of the Jacobian $\J(\x_*)$.
\begin{assumption}
\label{ass:nonsingular_J^*}
We assume there exists $\x_*\in\BR^d$ such that $\J(\x_*)^{\top}\F(\x_*)=\0$ and the Jacobian at $\x_*$ satisfies $\sigma_m(\J(\x_*))\geq \mu_*$ for some $\mu_*>0$, where $m=\min\{n,d\}$.
\end{assumption}
Combining Assumptions~\ref{ass:Jlip} and \ref{ass:nonsingular_J^*} suggests  Assumption~\ref{ass:nonsingular} can be satisfied by taking a proper region $\Omega$, which is introduced in the following proposition. 
%by the above assumption on $\J(\x_*)$.
\begin{proposition}
\label{prop:local_nonsingularity}
    Under Assumptions~\ref{ass:Jlip} and \ref{ass:nonsingular_J^*}, Assumption~\ref{ass:nonsingular} is satisfied with
    \begin{align*}
    \Omega
    &=\Omega_{\mathrm{over}}
    \triangleq\left\{\x:\|\x-\x_*\|\leq \frac{\mu^2_*}{4LL_2}\right\}
    \quad\text{and}\quad
    \mu=\frac{\mu_*}{\sqrt{2}}.
    \end{align*}
\end{proposition}

We then impose the following smoothness assumption with respect to the stationary point $\vx_*$.
\begin{assumption}
\label{ass:x_star_philip}
    We assume there exists a constant $\rho\geq 0$ such that 
    \begin{align}
        \label{eq:x_star_philip}
        \big\|\J(\x)^{\top}\F(\x)-(\H(\x_*)+\Q(\x_*))(\x-\x_*)\big\| \leq \rho \|\x-\x_*\|^2.
    \end{align}
    for all $\vx\in\BR^d$, where $\vx_*\in\BR^d$ is the stationary point of $\phi(\cdot)$ such that $\mJ(\vx_*)^\top\mF(\vx_*)=\vzero$.
\end{assumption}
The following proposition verifies that our Assumption \ref{ass:x_star_philip} is more general than several classical smoothness assumptions used in prior studies for nonlinear equations \citep{gould2019convergence,zhou2010global,vater2024convergence,liu2023block}.
\begin{proposition}
\label{prop:assumption-sati}
    Assumption~\ref{ass:x_star_philip} can be verified in one of the following three cases.
    \begin{enumerate}
        \item  Suppose Assumption~\ref{ass:Jlip} holds and the Hessian $\nabla^2 F_i(\cdot)$ is $L_3$-Lipschitz continuous with respect to the stationary point of $\phi(\cdot)$ for all $i\in\BR^n$
         \citep[Assumption AS.1]{gould2019convergence} \citep[Lemma 2.9--2.11]{zhou2010global}, i.e., there exists constant $L_3>0$ such~that
        \begin{align*}
        \|\nabla^2F_i(\x)-\nabla^2 F_i(\x_*)\|
        \leq L_3\|\x-\x_*\|    
        \end{align*}
        for all $\vx\in\BR^d$, where $\vx_*\in\BR^d$ is the stationary point of $\phi(\cdot)$ such that $\mJ(\vx_*)^\top\mF(\vx_*)=\vzero$.
        Then for all $\x\in\BR^d$ satisfying $\|\x-\x_*\|\leq\delta$, Assumption~\ref{ass:x_star_philip} holds with
        \begin{align}\label{eq:rho-cond-1}
        \rho
        =\frac{1}{2}\Big(2LL_2+L\sum_{i=1}^{n}\|\nabla^2 F_i(\x_*)\|
        +L_3\sqrt{n}L\delta+L_3\|\F(\x_*)\|_1\Big),
        \end{align}
        where $\|\cdot\|_1$ is the $\ell_1$-norm.
        \item  Suppose the Hessian of $\phi(\cdot)$ is $\rho_0$-Lipschitz continuous \citep[Assumption B]{zhou2010global} , i.e., there exists a constant $\rho_0>0$ such that
        \begin{align}\label{eq:phi-2-Lip-xy}
        \|\nabla^2\phi(\x)-\nabla^2\phi(\y)\|
        \leq\rho_0\|\x-\y\|
        \end{align}
        for all $\vx,\vy\in\BR^d$.
        Then Assumption~\ref{ass:x_star_philip} holds with $\rho=\rho_0/2$.
        \item Suppose Assumption~\ref{ass:Jlip} holds and the point $\x_*\in\BR^d$ is the solution of the nonlinear equations such that $\F(\x_*)= \0$ \citep{vater2024convergence, liu2023block}, then Assumption~\ref{ass:x_star_philip} holds with $\rho = {3LL_2}/{2}$.
        % \item If the condition in \citet{gratton2007approximate} holds.
    \end{enumerate}
\end{proposition}
% \begin{remark}
%     The conditions of Proposition~\ref{prop:assumption-sati} has been widely applied in the prior works. 
%     The first condition weakens the Lipschitz condition of $\nabla^2 F_i$, which has been used in  Assumption AS.1 in \citet{gould2019convergence} and Lemma 2.9--2.11 in \citet{zhou2010global}.
%     The second condition is used in Assumption B in\citet{zhou2010global} to guarantee the local superlinear rate of their GN-BFGS method.
%     The third condition is widely applied to solve the square nonlinear equations systems~.
% \end{remark}


We also make the following assumption on $\H(\x_*)$ and $\Q(\x_*)$, which is widely used and necessary to guarantee the local linear convergence of the Gauss--Newton method~\citep{gratton2007approximate,ortega2000iterative,wedin1974gauss}.
\begin{assumption}
\label{ass:deltacondi}
We suppose that
\begin{align}
\label{eq:condiDelta}
\theta\triangleq\Norm{\H(\x_*)^{-1/2}\Q(\x_*)\H(\x_*)^{-1/2}}\in[0,1),
\end{align}
where $\mH(\cdot)=\mJ(\cdot)^\top\mJ(\cdot)$, $\Q(\cdot)\triangleq \sum_{i=1}^n F_i(\cdot)\nabla^2 F_i(\cdot)$, and $\vx_*$ is the stationary point follows Assumption \ref{ass:nonsingular_J^*}.
\end{assumption}
% Throughout the overdetermined-case analysis, when $\theta=0$ we interpret the condition $c\in(1,1/\theta)$ as simply $c>1$.
Note that Assumption~\ref{ass:deltacondi} indicates the positive definiteness of $\nabla^2\phi(\x_*)$ for the stationary point, while it does not guarantee the positive definiteness of $\nabla^2\phi(\x)$ always holds even for~$\vx\in\BR^d$ in some local region.
However, the following lemma shows that we can guarantee the positive definiteness of the matrix
\begin{align*}
    \tilde{\nabla}^2\phi(\x)\triangleq \H(\x) + \Q(\x_*),
\end{align*}
which can be regarded as an approximation of the Hessian $\nabla^2\phi(\x) =\H(\x) + \Q(\x)$.
\begin{lemma}
\label{lm:pd_approximate_phi}
Under Assumptions~\ref{ass:Jlip}, \ref{ass:nonsingular_J^*}, and~\ref{ass:deltacondi}, 
for all $c\in(1,1/\theta)$ and $\x,\vy\in\BR^d$ such that~$\|\x-\x_*\|\leq(c-1)\mu_*^2/(4cLL_2)$, it holds
\begin{align}\label{eq:approx_phi_positive}
     (1-c\theta)\mu^2 \I_d \preceq (1-c\theta)\H(\x)\preceq  \tilde{\nabla}^2\phi(\x)\preceq (1+c\theta)\H(\x)\preceq (1+c\theta)L^2\I_d.
\end{align}
and
\begin{align}\label{eq:concor_phi}
    (1-\widehat M\|\y-\x\|)\tilde{\nabla}^2\phi(\x)
        \preceq\tilde{\nabla}^2\phi(\y)
        \preceq(1+\widehat M\|\y-\x\|)\tilde{\nabla}^2\phi(\x),
\end{align}
where $\widehat M\triangleq 2LL_2/((1-c\theta)\mu^2)$ and $\mu=\mu_*/\sqrt{2}$. 
\end{lemma}

Note that Gauss--Newton iteration \eqref{update:gn} for the case of $n\geq d$ can be written as
\begin{align}\label{eq:update_n>d}
\begin{split}    
    \vx_+= &\x-(\J(\x)^{\top}\J(\x))^{-1}\J(\x)^{\top}\mF(\vx)  \\
    = & \x- \mH^{-1}\nabla \phi(\x),
\end{split}
\end{align}
since we have $\mH(\vx)=\J(\x)^{\top}\J(\x)$ and $\nabla \phi(\x)=\J(\x)^{\top}\F(\x)$.
Recall that classical linear-quadratic convergence rate for iteration \eqref{eq:update_n>d} is established based on the distance $\|\x-\x_*\|$ \citep{gratton2007approximate,dennis1996numerical}.
However, this metric does not work for proposed RQGN methods because our iteration \eqref{eq:qNupdate} only accesses the Gram estimator $\mG$ rather than the exact $\mH(\vx)$.

For analyzing the convergence of our RQGN, we construct the weighted norm 
\begin{align}\label{eq:our_measure}
\begin{split}\    
    \lambda_{\mathrm{over}}(\x)
    \triangleq & \big\|(\H(\x)+\Q(\x_*))^{1/2}(\x-\x_*)\big\| \\
    = & \big\|(\tilde{\nabla}^2\phi(\x))^{1/2}(\x-\x_*)\big\|
\end{split}
\end{align}
to address the challenge of the approximation error between $\G$ and $\H(\x)$.
Note that Lemma~\ref{lm:pd_approximate_phi} shows that the matrix $\tilde{\nabla}^2\phi(\x)$ is positive-definite in the local region, which guarantees the metric $\lambda_{\mathrm{over}}(\cdot)$ is well-defined in our local convergence analysis.
Consequently, we establish the partial relation and explicit linear quadratic rate for iteration~\eqref{eq:update_n>d} as follows.
\begin{lemma}\label{lm:GNlinearquadra}
Under Assumptions~\ref{ass:Jlip}, \ref{ass:nonsingular_J^*}, \ref{ass:x_star_philip}, and \ref{ass:deltacondi}, we additionally suppose the iteration~\eqref{eq:qNupdate} in the case of $n\geq d$ holds 
\begin{align}\label{eq:condi_x_linear_quadra}
    \H(\x)\preceq\G\preceq\eta\H(\x),\quad
    \|\x-\x_*\|\leq\frac{(c-1)\mu_*^2}{4cLL_2},\quad
    \text{and}\quad\lambda_{\mathrm{over}}(\x)\leq
    \frac{(1-c\theta)^{3/2}\mu_*^5}{16D},
\end{align}
for some $\eta\geq1$ and $c\in(1,1/\theta)$, where
$D\triangleq L^3L_2+LL_2^2\|\F(\x_*)\|$.
Then it holds
\begin{align}
& \lambda_{\mathrm{over}}(\x_{+}) \leq \left(1-\frac{1-c\theta}{\eta}\right)\lambda_{\mathrm{over}}(\x) +C_1(\lambda_{\mathrm{over}}(\x))^2, \label{eq:overfit_linear_quadra}\\
&
\|\x_{+}-\x_*\| \leq \frac{1}{\sqrt{1-c\theta}\mu}
\left(\left(1-\frac{1-c\theta}{\eta}\right)
\lambda_{\mathrm{over}}(\x) 
+C_1(\lambda_{\mathrm{over}}(\x))^2\right), \label{eq:overfit_distance_lambda} \\
& \|\x_{+}-\x\| \leq\frac{L^2+L_2\|\F(\x_*)\|}{\sqrt{1-c\theta}\mu^3}
\lambda_{\mathrm{over}}(\x), \label{eq:movedistance_lambda}
\end{align}
where 
\begin{align}\label{eq:mu-C1}
\mu=\frac{\mu_*}{\sqrt{2}} 
\quad\text{and}\quad
C_1
&\triangleq\frac{2D}{(1-c\theta)^{3/2}\mu^5}
+\left(\frac{4(1+c\theta)}{(1-c\theta)^{3/2}}+\frac{2c}{c-1}\right)\frac{LL_2}{\mu^3}
+\frac{2\sqrt{1+c\theta}\,\rho}{(1-c\theta)\mu^3}.
\end{align}
\end{lemma}




We now combine Lemma~\ref{lm:GNlinearquadra}, Proposition~\ref{prop:Gram_concor}, and Lemma~\ref{lm:QN_order} to establish a linear convergence rate and additionally show that all iterates remain in the local region 
$\Omega_{\mathrm{over}}$.
\begin{theorem}
\label{thm:linear1n>d}
Under Assumptions~\ref{ass:Jlip}, \ref{ass:nonsingular_J^*}, \ref{ass:x_star_philip}, and \ref{ass:deltacondi}, 
we run Algorithm~\ref{alg:RQGN} with the initial point $\x_0$ such that
\begin{align}
\label{eq:linear_initial_condi}
\|\x_0-\x_*\|\leq\frac{(c-1)\mu_*^2}{4cLL_2}
\quad\text{and}\quad
\lambda_{\mathrm{over}}(\x_0)\leq \frac{1-c\theta}{10\eta_0 C_1},
\end{align}
for some $c\in(1,1/\theta)$, where $C_1$ follows the definition in Lemma \ref{lm:GNlinearquadra}.
Then it holds $\x_t\in\Omega_{\mathrm{over}}$ for all $t$. Moreover, we also have
\begin{align}
\label{eq:induc_lambda}
\H_t\preceq\G_t\preceq\eta_0{\rm e}^{2M\sum_{i=0}^{t-1}r_i}\H_t
\preceq\frac{3\eta_0}{2}\H_t
\quad\text{and}\quad
\lambda_{\mathrm{over}}(\x_t)\leq\left(1-\frac{1-c\theta}{2\eta_0}\right)^t\lambda_{\mathrm{over}}(\x_0).
\end{align}
\end{theorem}
Theorem~\ref{thm:linear1n>d} guarantees that $\x_t\in\Omega_{\mathrm{over}}$ for all $t$, verifying the trajectory condition required by Lemma~\ref{lm:Gram_approximate} with $\Omega=\Omega_{\mathrm{over}}$.
Hence, the sequence $\{\eta_t\}$ defined in~\eqref{eq:def_eta} is well-defined along the iterates like the underdetermined case.
Combining Lemma~\ref{lm:Gram_approximate} and Theorem~\ref{thm:linear1n>d}, we establish finite-time bounds on the Gram approximation error and a two-stage local linear convergence rate.
\begin{theorem}
\label{thm:two-stage-n>d}
Under Assumptions~\ref{ass:Jlip}, \ref{ass:nonsingular_J^*}, \ref{ass:x_star_philip}, and \ref{ass:deltacondi}, run Algorithm~\ref{alg:RQGN} with the initial point~$\x_0$ satisfying
\begin{align}
\label{eq:inicondi_final_n>d}
\|\x_0-\x_*\|\leq\frac{(c-1)\mu_*^2}{4cLL_2}
\qquad\text{and}\qquad
\lambda_{\mathrm{over}}(\x_0)\leq
\frac{(1-\gamma)(1-c\theta)}{10(2-\gamma)\eta_0C_2},
\end{align}
for some $c\in(1,1/\theta)$, where
$C_2\triangleq C_1+4\tilde{c}dD/(\sqrt{1-c\theta}\mu^5)$,
the values of $\gamma$ and $\tilde{c}$ for each of Options I--III follow Lemma~\ref{lm:Gram_approximate}, 
and the values of $D$, $\mu$, and $C_1$ follow Lemma~\ref{lm:GNlinearquadra}.
Then we have
\begin{align}
\label{eq:eta-ratio-over}
\EBP{\eta_t-1}
&\leq q_{\text{\rm over}}(t+1)(1-\nu)^t
\quad\text{and}\quad
\EBP{\frac{\lambda_{\mathrm{over}}(\x_{t+1})}{\lambda_{\mathrm{over}}(\x_t)}}
\leq c\theta+q_{\text{\rm over}}(t+1)(1-\nu)^t,
\end{align}
where $\nu\triangleq\min\{\gamma,(1-c\theta)/(2\eta_0)\}$,
$q_{\text{\rm over}}\triangleq {\rm e}^{1/5}
(\delta_0+(1-\gamma)(1-c\theta)/(10(2-\gamma)\eta_0))$, and the
value of~$\delta_0$ for each of Options I--III follows
Lemma~\ref{lm:Gram_approximate}.
Moreover, for every $p\in(0,1)$, there exists
$t_0=\OM\big(\log(1+q_{\text{\rm over}}/(p(1-c\theta)\nu^2))/\nu\big)$
such that, with probability at least $1-p$, we have
\begin{align}
\label{eq:n>dtwostage}
\lambda_{\mathrm{over}}(\x_{t_0+t})
\leq\left(\frac{1+c\theta}{2}\right)^t
\left(1-\frac{1-c\theta}{2\eta_0}\right)^{t_0}\lambda_{\mathrm{over}}(\x_0),
\end{align}
simultaneously for all $t\geq0$.
\end{theorem}

\begin{remark}
Note that even for the standard Gauss--Newton iteration \eqref{eq:update_n>d},
exact the Gram matrix~$\H(\cdot)=\J(\cdot)^{\top}\J(\cdot)$ does not include the term $\mQ(\cdot)$ in the Hessian  $\nabla^2\phi(\x)=\H(\x)+\Q(\x)$.
Naturally, our RQGN methods also does not consider the term $\mQ(\cdot)$ in the construction of the Gram estimator.
This results the expected ratio~\eqref{eq:eta-ratio-over} 
for shown in Theorem~\ref{thm:two-stage-n>d} contains an additional term $c\theta$, so that the overdetermined case cannot always enjoy the superlinear convergence rate like that of Theorem~\ref{thm:n<dsuperlinear} for the underdetermined case.
\end{remark}



If we strengthen the conditions of Theorem \ref{thm:linear1n>d} by assuming that Assumption \ref{ass:deltacondi} satisfies with $\theta=0$, we can show that our RQGN methods enjoys the superlinear local convergence.

\begin{corollary}
\label{col:theta=0n>d}
Under the same notation and setting as Theorem~\ref{thm:two-stage-n>d}, suppose Assumption~\ref{ass:deltacondi} holds with $\theta=0$ and $c>1$.
% Under Assumptions~\ref{ass:Jlip}, \ref{ass:nonsingular_J^*}, \ref{ass:x_star_philip}, and \ref{ass:deltacondi} with $\theta=0$, run Algorithm~\ref{alg:RQGN} with the initial point $\x_0$ satisfying condition~\eqref{eq:inicondi_final_n>d} with $\theta=0$, where the values of $\gamma$ and $\tilde{c}$ for each of Options I--III follow Lemma~\ref{lm:Gram_approximate}, $D$, $\mu$, and $C_1$ follow the definitions in Lemma~\ref{lm:GNlinearquadra} with $\theta=0$, and $C_2\triangleq C_1+4\tilde{c}dD/\mu^5$.
%Let $q_{\text{\rm over}}\triangleq{\rm e}^{1/5}\big(\delta_0+(1-\gamma)/(10(2-\gamma)\eta_0)\big)$, where the value of $\delta_0$ for each of Options I--III follows Lemma~\ref{lm:Gram_approximate}.
Then, for all $p\in(0,1)$, there exists $t_0=\OM(\log(1+q_{\text{\rm over}}/(p\gamma^2))/\gamma)$ such that it holds
\begin{align*}
\lambda_{\mathrm{over}}(\x_{t_0+t})\leq \left(1-\frac{\gamma}{\gamma+1}\right)^{\frac{t(t-1)}{2}}\left(\frac{1}{2}\right)^t\left(1-\frac{1}{2\eta_0}\right)^{t_0}\lambda_{\mathrm{over}}(\x_0)
\end{align*}
with probability at least $1-p$ for all $t\geq1$.
\end{corollary}

\begin{remark}
    If there exists $\vx^*\in\BR^d$ such that $\F(\x^*)=\0$ (case 3 in Proposition~\ref{prop:assumption-sati}), then Assumption~\ref{ass:deltacondi} holds with $\theta = 0$.            
    This implies our local superlinear convergence result in Corollary~\ref{col:theta=0n>d} can apply to several classical problem settings such as  square nonlinear equations~\citep{lin2021explicit,liu2023block} and strongly-convex-strongly-concave minimax optimization~\citep{liu2022quasi}.
\end{remark}
 


% \section{Globally Convergent Randomized Quasi-Gauss--Newton Method}
% In this section, we provide a globally convergent randomized quasi-Gauss--Newton method by using a step size that is proportional to $\sqrt{1/\lambda_{\mathrm{under}}(\x)}$.

% \begin{algorithm}[t]
% \caption{Globally Convergent Randomized Quasi-Gauss--Newton Method}\label{alg:GLOBAL-RQGN}
% \begin{algorithmic}[1]
% \STATE \textbf{Input:} $\G_0$ and $M$ such that $\G_0\succeq \H(\x_0)$ and $M\geq0$, ${\rm QN} \in \left\{\bfgs,\fbfgs,\srk\right\}$ \\[0.15cm]
% \STATE \textbf{for} $t=0,1\dots$\\[0.15cm]
%             \STATE \quad $\alpha_t = \min\{1,\sqrt{\frac{\mu^2}{\eta_0\|\F(\x_t)\|L_2}}\}$
%  		\STATE \quad
%  		    $\x_{t+1}=
%                     \x_t-\alpha_t\J(\x_t)^{\top}\G_t^{-1}\F(\x_t)$ \\[0.15cm]
%  		\STATE \quad $r_t=\|\x_{t+1}-\x_{t}\|$ \\[0.15cm]
%  		\STATE \quad
%  	        $\widetilde{\G}_{t}=(1+Mr_t)\G_t$\\[0.15cm]
%             \STATE \quad 
%             $\left[\U_{t}\right]_{ij}\overset{\rm{i.i.d}}{\sim}{\fN(0,1)}$ \\[0.15cm]
%  		\STATE \quad 
%  		    $\G_{t+1}= {\rm QN}(\widetilde{\G}_t,\H(\x_{t+1}),\U_t)$\\[0.15cm]
% \STATE \textbf{end for}
% \end{algorithmic}
% \end{algorithm}

% \begin{proof}
%     We modify the proof of Lemma~\ref{lm:linear-quadra_n<d}, where we have
%     \begin{align}
%     r_t = \|\alpha_t\J(\x_t)\G_t^{-1}\F(\x_t) \| = \alpha_t\|\J(\x_t)^{\top}\G_t^{-1}\F(\x_t)\|\leq \alpha_t\lambda_{\mathrm{under}}(\x_t),
%     \end{align}
%     and we modify \eqref{eq:HFx+} by 
% \begin{align}
% \label{eq:HFx+global}
% \begin{split}
%      \Norm{(\H(\x_t))^{-1/2}\F(\x_{t+1})}
%      \overset{\eqref{eq:G_condi_2},\eqref{eq:lip}}{\leq} \left(1-\frac{1}{\eta_t}\right)\lambda_{\mathrm{under},t}+\Norm{\H(\x_t)^{-1/2}}\frac{L_2\alpha_t^2r_t^2}{2}\overset{\eqref{eq:boundr_lambda}}{\leq} \left(1-\frac{1}{\eta}\right)\lambda_{\mathrm{under},t}+\frac{L_2}{2\mu}(\lambda_{\mathrm{under},t})^2.
%      \end{split}
% \end{align}
% Thus, it holds that
% \begin{align*}
%     \lambda_{\mathrm{under},t+1} \leq 
% \end{align*}
% \end{proof}

% \section{Improvements over the Conference Version}\label{sec:conference}. 
% {\color{blue}
% In our earlier version, we propose randomized quasi-Gauss--Newton methods for minimax optimization $\min_{\x_1}\max_{\x_2} f(\x)$, where $\x=[\x_1^{\top},\x_2^{\top}]^{\top}$, which can be generalized to solve nonlinear equations for the case $n=d$. The iteration can be written as
% \begin{align*}
%     \x_{+} = \x-\G^{-1}\nabla^2 f(\x)\nabla f(\x),
% \end{align*}
% where $\G$ is an approximation to the square of the Hessian $(\nabla^2f(\x))^{2}$.  
% % Since $\nabla^2 f(\x)$ is symmetric, this iteration is essentially the same as that for $n\geq d$ in \eqref{update:qn} taking $\F(\x)=\nabla f(\x)$.
% The main differences between this paper and the conference version are threefold.
% Firstly, we generalize the methods in the conference version to both overdetermined and underdetermined cases.
% In addition, our previous version used a simple metric $\|\F(\x)\|_2 $ to describe convergence, which is not suitable for general cases. 
% For example, the solution of $\F(\x)=\0$ may not exist when the nonlinear equations system is overdetermined.
% We carefully design the metrics for overdetermined and underdetermined cases separately and provide non asymptotic convergence analysis for both cases on these new designed metrics. 

% The established linear quadratic rate also provides an explicit convergence guarantee for the Gauss--Newton methods and can be viewed as a generalization of the result in our conference version~\citep{liu2022quasi}.
% \begin{itemize}
% \item For the Gauss--Newton method to solve overdetermined nonlinear equations, we have $\eta=1$, and the linear quadratic rate in \eqref{eq:overfit_linear_quadra} becomes $\lambda_{\mathrm{over}}(\x_{+})\leq c\theta \lambda_{\mathrm{over}}(\x)+C_1\lambda_{\mathrm{over}}(\x)^2$. The diminished constant $c\theta$ before the linear term is only related to the constant $\theta$ and is condition-number-free.
% \item For quasi-Newton methods to solve saddle point problems as considered in our conference version, we have $\theta=0$ in Assumption~\ref{ass:deltacondi}. Thus, the linear quadratic rate in \eqref{eq:overfit_linear_quadra} becomes
% \begin{align*}
% \lambda_{\mathrm{over}}(\x_{+})
% &\leq\left(1-\frac{1}{\eta}\right)\lambda_{\mathrm{over}}(\x)
% +C_1\lambda_{\mathrm{over}}(\x)^2.
% \end{align*}
% The linear factor $1-\frac{1}{\eta}$ is the same as that in \citet[Lemma 3.14]{liu2022quasi}.
% \end{itemize}
% % The local convergence rates can cover the one in our conference version when $n=d$ under weaker assumptions.

% Theorem~\ref{thm:two-stage-n>d} shows that RQGN first converges at a linear rate whose contraction factor depends on $\eta_0$, and then converges at a faster linear rate whose contraction factor is independent of $\eta_0$.
% Quasi-Newton methods for strongly convex optimization~\citep{rodomanov2021greedy,lin2022explicit} or saddle point problems~\citep{liu2022quasi} also have two-stage convergence, which is linear in the first stage but superlinear in the second stage since these problems can always guarantee $\theta=0$ in Assumption~\ref{ass:deltacondi}.
% In the following corollary, we show that our RQGN method can also enjoy local superlinear rate under this case.
% }


\section{Numerical Experiments}
\label{sec:exp}
In this section, we conduct experiments on scientific computing and machine learning applications to verify the advantage of proposed methods for solving both underdetermined and overdetermined nonlinear equations.
We run our RQGN (Algorithm \ref{alg:RQGN}) with SR1 and fast BFGS updates, denoted by RQGN-SR1 and RQGN-FBFGS, respectively.
%The deta (see implementation in Appendix~\ref{app:fbfgs})
We set~$\G_0=L^2\I_m$ and tune parameters $L,M$ over $\{1,10,100,1000\}$ for the algorithm.
All experiments are performed on a PC with Apple M1 and all algorithms are implemented in Python 3.8.12.



\subsection{Empirical Results for the Underdetermined Case}
We first consider the problem of finding eigenvalues of the symmetric matrix $\A\in\RB^{n\times n}$ \citep{vater2024convergence}, which can be formulated by solving the underdetermined nonlinear equation  \eqref{eq:nonlinear_obj} with $\mF:\BR^{n+1}\to\BR^n$ such that
\begin{align}
\label{eq:eigenfinding}
\F(\x)\triangleq (\A-\lambda\I_{n})\u
\qquad\text{and}\qquad
\x= \begin{bmatrix}
    \lambda \\ \u
\end{bmatrix}\in\RB^{n+1}    
\end{align}
We follow the setting of \citet{vater2024convergence} by letting $\A=[a_{i,j}]$ be a tridiagonal matrix with $a_{i,i}=10$ for $i\in\{1,\ldots,n\}$ and $a_{i,i-1}=a_{i-1,i}=-1$ for $i\in\{2,\ldots,n\}$. 

We compare our RQGN methods against baselines gradient descent (GD)~\citep{nesterov2018lectures} and generalized Broyden method~\citep{vater2024convergence}.
For GD, we perform the iteration
$\x_{t+1}=\x_t - \gamma \J(\x_t)^{\top}\F(\x_t)$ in the view of the nonlinear least-square formulation and tune the stepsize $\gamma$ from~$\{0.01,0.1,1,10,100\}$.
For generalized Broyden method, we use the scaled-identity matrix as the initialization like our RQGN.
We report empirical results on iterations and CPU time for $n\in\{300,400,500\}$ in Figure~\ref{fig:eigen_find}.
We observe that our proposed RQGN methods outperform the baseline algorithm in most cases.

\begin{figure}[ht]
\centering
\begingroup
\setlength{\tabcolsep}{1pt}
\renewcommand{\arraystretch}{0.96}
\begin{tabular}{@{}ccc@{}}
\includegraphics[width=0.327\textwidth,trim=11pt 6pt 11pt 6pt,clip]{Eigen_300.res.pdf} &
\includegraphics[width=0.327\textwidth,trim=11pt 6pt 11pt 6pt,clip]{Eigen_400.res.pdf} &
\includegraphics[width=0.327\textwidth,trim=11pt 6pt 11pt 6pt,clip]{Eigen_500.res.pdf} \\[-1pt]
{\footnotesize (a) $n=300$ (iteration)} & {\footnotesize (b) $n=400$ (iteration)} & {\footnotesize (c) $n=500$ (iteration)} \\[10pt]
\includegraphics[width=0.327\textwidth,trim=11pt 6pt 11pt 6pt,clip]{Eigen_300.time.pdf} &
\includegraphics[width=0.327\textwidth,trim=11pt 6pt 11pt 6pt,clip]{Eigen_400.time.pdf} &
\includegraphics[width=0.327\textwidth,trim=11pt 6pt 11pt 6pt,clip]{Eigen_500.time.pdf} \\[-1pt]
{\footnotesize (d) $n=300$ (time)} & {\footnotesize (e) $n=400$ (time)} & {\footnotesize (f) $n=500$ (time)}
\end{tabular}
\endgroup
\caption{Empirical results of iteration counts and CPU time (seconds) against $\|\F(\x)\|$ for the underdetermined problem with nonlinear mapping \eqref{eq:eigenfinding}.
}
\label{fig:eigen_find}
\end{figure}

%Appendix~\ref{app:additional_under} compares RQGN with the generalized Broyden method~\citep{vater2024convergence} under exact and scaled-identity initial estimators.


\subsection{Empirical Results for the Overdetermined Cases}
We then consider the generalized linear model \citep{shao2025convergence}, which can be formulated by  the nonlinear least square problem \eqref{eq:nonlinear_least_obj} with $\mF:\BR^d\to\BR^{2d}$ such that
\begin{align}
\label{eq:overfit}
    \F(\x)  \triangleq \begin{bmatrix}
        \nabla f_1(\x)\\
        \alpha\nabla f_2(\x)
    \end{bmatrix},
\end{align}
where $\alpha>0$ is the scaled parameter and $f_1:\BR^d\to\BR$ and $f_2:\BR^d\to\BR$ are objectives of log-sum-exp and logistic regression problems  \citep{rodomanov2021greedy} which are defined as 
\begin{align*}
f_1(\x)\triangleq \ln\left(\sum_{j=1}^N\exp(\va_j^{\top}\x-b_j)\right)+\frac{1}{2}\sum_{j=1}^N(\va_j^\top\x)^2
~\text{and}~f_2(\x)\triangleq \sum_{j=1}^N
 \ln\left(1+\exp\left(-b_j\va_j^{\top}\x\right)\right),
\end{align*}
where $\{(\va_i, b_i)\}_{i=1}^N$ is the training set with the feature~$\va_j\in\BR^d$ and the corresponding label~$b_j\in\BR$.
Here, we set $\alpha = 0.1$ in our experiments.

We compare our RQGN methods against the baseline GD by tuning the stepsize over $0.01$, $0.1$, $1$, $10$, and~$100$.
We perform the experiments on three real-world data sets, ``a1a'' ($N=1,605$, $d=123$), ``w1a'' ($N=2,477$, $d=300$), and ``splice'' ($N=1,000$, $d=60$), which can be downloaded from the LIBSVM repository~\citep{CC01a}.
We report empirical results on iterations and CPU time for all three data sets in Figure~\ref{fig:overfit}.
We can observe that our RQGN methods enjoy faster local convergence behavior and take less running time than the baseline.

\begin{figure}[!tbp]
\centering
\begingroup
\setlength{\tabcolsep}{1pt}
\renewcommand{\arraystretch}{0.96}
\begin{tabular}{@{}ccc@{}}
\includegraphics[width=0.327\textwidth,trim=11pt 6pt 11pt 6pt,clip]{hybrid_logit.a1a.res.pdf} &
\includegraphics[width=0.327\textwidth,trim=11pt 6pt 11pt 6pt,clip]{hybrid_logit.w1a.res.pdf} &
\includegraphics[width=0.327\textwidth,trim=11pt 6pt 11pt 6pt,clip]{hybrid_logit.splice.res.pdf} \\[-1pt]
{\footnotesize (a) ``a1a'' (iteration)} & {\footnotesize (b) ``w1a'' (iteration)} & {\footnotesize (c) ``splice'' (iteration)} \\[10pt]
\includegraphics[width=0.327\textwidth,trim=11pt 6pt 11pt 6pt,clip]{hybrid_logit.a1a.time.pdf} &
\includegraphics[width=0.327\textwidth,trim=11pt 6pt 11pt 6pt,clip]{hybrid_logit.w1a.cputime.pdf} &
\includegraphics[width=0.327\textwidth,trim=11pt 6pt 11pt 6pt,clip]{hybrid_logit.splice.cputime.pdf} \\[-1pt]
{\footnotesize (d) ``a1a'' (time)} & {\footnotesize (e) ``w1a'' (time)} & {\footnotesize (f) ``splice'' (time)}
\end{tabular}
\endgroup
\caption{Empirical results of iteration counts and CPU time (seconds) against $\|\J(\x)^{\top}\F(\x)\|$ for the overdetermined problem with nonlinear mapping \eqref{eq:overfit}.}
\label{fig:overfit}
\end{figure}



\subsection{Numerical Results for the Square Case}
We also consider the problem of fairness-aware machine learning, which can be formulated as the minimax optimization problem
\begin{align}
\label{eq:minimax_fair}
\min_{\u\in\RB^{d-1}}\max_{v\in\RB}    f(\u,v) \triangleq \frac{1}{N}\sum_{i=1}^N(l_1(\va_i,b_i,\u)-\beta l_2(\va_i^\top\u,c_i,v))+\lambda\|\u\|^2-\gamma v^2,
\end{align}
where $(\vu,v)\in\BR^d\times\BR$ is the weights of the model, $\{(\va_i,b_i,c_i)\}_{i=1}^N$ is the training set such that $\va_i\in\RB^{d}$ contains all the input variables, $b_i\in\RB$ is the output, $c_i\in\RB$ is the input variable that we want to make it unbiased~\citep{lowd2005adversarial,zhang2018mitigating}, and~$\lambda,\beta,\gamma>0$ are the hyperparameters.
Here, the loss $l_1$ and $l_2$ are the logit functions, defined as ${\rm logit}(\va,b,\w)= \log(1+\exp(-b\va^\top\w))$.
The minimax optimization problem (\ref{eq:minimax_fair}) can be viewed as solving the square nonlinear equation $\mF(\vx)=\vzero$ such that
\begin{align}\label{eq:min-max}
\x = \begin{bmatrix}
\u \\ v    
\end{bmatrix}
\qquad\text{and}\qquad
\F(\x)\triangleq \begin{bmatrix}
~~\nabla_{\u}f(\u,v)\\
-\nabla_v f(\u,v)
\end{bmatrix}.
\end{align}

We compare our RQGN methods with the extragradient method (EG)~\citep{korpelevich1976extragradien,tseng1995linear}, which is the optimal first-order method for strongly-convex strongly-concave optimization.
% Note that our RQGN methods with $k=1$ correspond to the algorithms in our conference version~\citep{liu2022quasi}.
For EG, we tune the step size from $\{0.01,0.05, 0.1, 0.5, 1\}$.
We set the parameters $\beta,\lambda$, and $\gamma$ as $0.5$, $0.001$, and $0.001$, respectively.
We perform the experiments on datasets ``heart'' ($N=270$, $d = 13$)~\citep{CC01a},  ``adult'' ($N= 32,561$, $d= 123$), and ``law school'' ($N=20,798$, $d= 380$)~\citep{quy2021survey}.
We report empirical results on iterations and CPU time in
Figure~\ref{fig:min-max}.
We observe that our RQGN methods converge much faster than EG.

\begin{figure}[!tbp]
\centering
\begingroup
\setlength{\tabcolsep}{1pt}
\renewcommand{\arraystretch}{0.96}
\begin{tabular}{@{}ccc@{}}
\includegraphics[width=0.327\textwidth,trim=11pt 6pt 11pt 6pt,clip]{heart.res.pdf} &
\includegraphics[width=0.327\textwidth,trim=11pt 6pt 11pt 6pt,clip]{adult.res.pdf} &
\includegraphics[width=0.327\textwidth,trim=11pt 6pt 11pt 6pt,clip]{lawschool.res.pdf} \\[-1pt]
{\footnotesize (a) ``heart'' (iteration)} & {\footnotesize (b) ``adult'' (iteration)} & {\footnotesize (c) ``law school'' (iteration)} \\[10pt]
\includegraphics[width=0.327\textwidth,trim=11pt 6pt 11pt 6pt,clip]{heart.time.pdf} &
\includegraphics[width=0.327\textwidth,trim=11pt 6pt 11pt 6pt,clip]{adult.time.pdf} &
\includegraphics[width=0.327\textwidth,trim=11pt 6pt 11pt 6pt,clip]{lawschool.time.pdf} \\[-1pt]
{\footnotesize (d) ``heart'' (time)} & {\footnotesize (e) ``adult'' (time)} & {\footnotesize (f) ``law school'' (time)}
\end{tabular}
\endgroup
\caption{Empirical results of iteration counts and CPU time (seconds) against $\|\F(\x)\|$ for the square system with nonlinear mapping \eqref{eq:overfit}.}
\label{fig:min-max}
\end{figure}




%In addition, RQGN with $k=10$ is faster than RQGN with $k=1$ (algorithms in the conference version), because it takes advantage of the block updates.
% Figures~\ref{fig:eigen_find}--\ref{fig:min-max} include the complete CPU time curves, which show the same qualitative advantage as the iteration curves. Appendix~\ref{app:additional_experiments} gives additional comparisons with generalized and randomized Broyden methods, including the effect of the initial estimators.


\section{Conclusion and Future Work}
\label{sec:conclu}
In this paper, we propose randomized quasi-Gauss--Newton (RQGN) methods for solving both overdetermined and underdetermined nonlinear equations.
We establish local convergence guarantees to show that our RQGN enjoys condition-number-free linear rates for the overdetermined case and superlinear rates for the underdetermined equation.
Notably, our convergence results do not depend on the sufficiently accurate initial Jacobian estimator that is required by Broyden's methods.
The empirical results validate the efficiency of our proposed RQGN methods.
In future work, it is interesting to study the global convergence of quasi-Newton methods for general nonlinear equations, extending the existing results for convex optimization and monotone variational inequalities~\citep{jin2024non,jin2024nonasymptotic,rodomanov2024global,jiang2024online,jiang2024improved,kamzolov2023accelerated,agafonov2024exploring,wang2024global}. 

% \acks{[Please add acknowledgments here and disclose all funding that supported this work and all competing interests before submission.]}

\appendix

\section{Proofs for Sections~\ref{sec:pre} and \ref{sec:rqn} }
\label{app:proofs_preliminaries}
We provide detailed proofs for the propositions in Sections~\ref{sec:pre} and \ref{sec:rqn}.

\subsection{Proof of Proposition~\ref{prop:gram_lip}}
\label{app:proof_prop_gram_lip}
\begin{proof}
The $L$-Lipschitz continuity of $\F(\,\cdot\,)$ follows directly from~equation \eqref{eq:lip}.
For the partial relation and the Lipschitz continuity of $\mH(\cdot)$, we focus on the case of $n\geq d$. 
Under Assumption~\ref{ass:Jlip}, we achieve
\begin{align*}
\|\H(\x)\|
&=\|\J(\x)^{\top}\J(\x)\|
\leq\|\J(\x)^{\top}\|\|\J(\x)\|
\leq L^2,
\end{align*}
which means $\H(\x)\preceq L^2\I_m$.
Besides, it also holds
\begin{align*}
    \|\H(\x)-\H(\y)\|\leq \big\|\J(\x)^{\top}\big\|\Norm{\J(\x)-\J(\y)}+\big\|\J(\x)^{\top}-\J(\y)^{\top}\big\|\Norm{\J(\y)}\leq 2LL_2\|\x-\y\|.
\end{align*}
The proof for the case $n<d$ is almost identical to that of $n\geq d$, so we omit it here.
\end{proof}

\subsection{Proof of Proposition~\ref{prop:gram_lower}}
\label{app:proof_prop_gram_lower}
\begin{proof}
Under Assumption~\ref{ass:nonsingular}, it holds
$\sigma_m(\H(\x))=\sigma_m^{2}(\J(\x))\geq\mu^2$,    
which implies $\mu^2\I_m\preceq\H(\x)$. 
\end{proof}

% \section{Proofs for Section~\ref{sec:rqn}}
% \label{app:proofs_rqgn}

\subsection{Proof of Proposition~\ref{prop:Gram_concor}}
\label{app:proof_prop_gram_concordance}
\begin{proof}
From Proposition~\ref{prop:gram_lip}, we have
    \begin{align*}
    -2L_2L \|\x-\y\|\cdot\I_m \preceq   \H(\y)-\H(\x)\preceq 2L_2L \|\x-\y\|\cdot\I_m.
    \end{align*}
    By Proposition~\ref{prop:gram_lower}, we have
    \begin{align*}
     -\frac{2L_2L}{\mu^2}\|
    \x-\y\| \H(\y)\preceq   \H(\y)-\H(\x)\preceq \frac{2L_2L}{\mu^2}\|\x-\y\|\H(\x),
    \end{align*}
    rearranging above inequalities directly indicates equation \eqref{eq:Gram_bound_0}.
\end{proof}

\section{Proofs for Section~\ref{sec:local_convergence}}
\label{app:proofs}
We first provide some basic results for later proofs.

\begin{lemma}
\label{lm:matrix_order_equivalence}
Let $\A\in\BR^{m\times m}$ be positive definite and
$\B\in\BR^{m\times m}$ be symmetric.
For $a_1,a_2\in\BR$ such that $a_1\leq a_2$, the following two statements are equivalent: 
\begin{itemize}
\item $a_1\A\preceq\B\preceq a_2\A$;
\item $a_1\I_m\preceq\A^{-1/2}\B\A^{-1/2}\preceq a_2\I_m$
\end{itemize}
\end{lemma}
\begin{proof}
Suppose the first statement $a_1\A\preceq\B\preceq a_2\A$ holds.
For all $\v\in\BR^m$, applying this relation to vector $\A^{-1/2}\v$ indicates
\begin{align*}
a_1\|\v\|^2
&=a_1(\A^{-1/2}\v)^{\top}\A(\A^{-1/2}\v)\leq
(\A^{-1/2}\v)^{\top}\B(\A^{-1/2}\v)\\
&=\v^{\top}\A^{-1/2}\B\A^{-1/2}\v
\leq a_2(\A^{-1/2}\v)^{\top}\A(\A^{-1/2}\v)
=a_2\|\v\|^2.
\end{align*}
Hence, we achieve $a_1\I_m\preceq\A^{-1/2}\B\A^{-1/2}\preceq a_2\I_m$.

Conversely, suppose the second statement
$a_1\I_m\preceq\A^{-1/2}\B\A^{-1/2}\preceq a_2\I_m$ holds.
For all~$\v\in\BR^m$, applying this relation to vector $\A^{1/2}\v$ indicates
\begin{align*}
a_1\v^{\top}\A\v
&=a_1\|\A^{1/2}\v\|^2
\leq(\A^{1/2}\v)^{\top}\A^{-1/2}\B\A^{-1/2}(\A^{1/2}\v)\\
&=\v^{\top}\B\v
\leq a_2\|\A^{1/2}\v\|^2
=a_2\v^{\top}\A\v.
\end{align*}
Hence, we achieve
$a_1\A\preceq\B\preceq a_2\A$.
\end{proof}

We next state the following lemma.
\begin{lemma}[{\citet[Lemma~26]{lin2022explicit}}]
\label{lm:lin_lemma26}
Let $\{X_t\}_{t\geq0}$ be a sequence of nonnegative random variables.
Suppose that there exist constants $a\geq0$ and $\tau>1$ such that
\begin{align*}
\mathbb E[X_t]
&\leq a\biggl(1-\frac{1}{\tau}\biggr)^t
\end{align*}
for every $t\geq0$. Then, for all $p\in(0,1)$, with probability at
least $1-p$, we have
\begin{align*}
X_t
&\leq\frac{a\tau^2}{p}
\biggl(1-\frac{1}{1+\tau}\biggr)^t
\end{align*}
simultaneously for every $t\geq0$.
\end{lemma}

\subsection{Proof of Lemma~\ref{lm:Gram_approximate}}
\label{app:proof_lemma_gram_approximate}
\begin{proof}
We first use induction to prove 
\begin{align}\label{eq:induction-G-H}
    \H_t\preceq\G_t
\end{align}
holds for all $t\geq 0$.
The induction base in the case of $t=0$ directly holds by our assumption.
For the induction step, we suppose the partial relation \eqref{eq:induction-G-H} holds for all $t=\hat t$. 
Note that Proposition~\ref{prop:Gram_concor} implies
\begin{align}
\label{eq:gram_transport_lm}
\frac{1}{1+Mr_{\hat t}}\H_{\hat t}
\preceq \H_{{\hat t}+1}
\preceq(1+Mr_{\hat t})\H_{\hat t}.
\end{align}
Hence, the induction hypothesis $\H_{\hat t}\preceq\G_{\hat t}$ implies
\begin{align*}
\H_{{\hat t}+1}\overset{\eqref{eq:gram_transport_lm}}\preceq(1+Mr_{\hat t})\H_{\hat t}
\preceq(1+Mr_{\hat t})\G_{\hat t}=\widetilde{\G}_{\hat t}.
\end{align*}
Applying Lemma~\ref{lm:QN_order} to the update \eqref{eq:update_G}  yields $\H_{\hat t+1}\preceq\G_{\hat t+1}$, which finished the induction.

% For any positive semidefinite matrix $\A$, we have
% $\A\preceq\tr{\A}\I$ because its largest eigenvalue is no larger than the sum of its eigenvalues.
We then show the relation $\mG_t\preceq(1+\delta_t)\H_t$.
The relation \eqref{eq:induction-G-H} implies
\begin{align*}
\H_t^{-1/2}(\G_t-\H_t)\H_t^{-1/2}\succeq\0.
\end{align*}
Recall that all positive semi-definite matrix $\A\in\BR^{m\times m}$ satisfies
$\A\preceq\tr{\A}\I_m$. 
Applying this fact with $\mA=\H_t^{-1/2}(\G_t-\H_t)\H_t^{-1/2}$, we achieve
\begin{align}\label{eq:delta-GH-1}
\begin{split}    
\G_t
&=\H_t^{1/2}\Bigl(\I_m+\H_t^{-1/2}(\G_t-\H_t)\H_t^{-1/2}\Bigr)\H_t^{1/2}\\
&\preceq\H_t^{1/2}\big(\I_m+{\rm tr}\big(\H_t^{-1/2}(\G_t-\H_t)\H_t^{-1/2}\big)\I_m\big)\H_t^{1/2}\\
&=\big(1+{\rm tr}\big(\H_t^{-1/2}(\G_t-\H_t)\H_t^{-1/2}\big)\big)\H_t\\
&=\big(1+\tr{\H_t^{-1}(\G_t-\H_t)}\big)\H_t \\
&=\big(1+\sigma_{\H_t}(\G_t)\big)\H_t.
\end{split}
\end{align}
Moreover, the partial relation \eqref{eq:induction-G-H} indicates
\begin{align}\label{eq:G-H-I}
\G_t-\H_t\preceq\|\G_t-\H_t\|\I_m     
\end{align}
and
\begin{align}\label{eq:G-H-tr}
\|\G_t-\H_t\|\leq\tr{\G_t-\H_t}=\tau_{\H_t}(\G_t).     
\end{align}
Additionally, Propositions~\ref{prop:gram_lip} and~\ref{prop:gram_lower} imply
\begin{align}\label{eq:H-mu-L}
\mu^2\I_m\preceq\H_t\preceq L^2\I_m.    
\end{align}
Hence, it holds
\begin{align}\label{eq:delta-GH-2}
\begin{split}
\G_t
&\,=\,\H_t+(\G_t-\H_t)\\
&\overset{\eqref{eq:G-H-I}}\preceq\H_t+\|\G_t-\H_t\|\I_m\\
&\overset{\eqref{eq:H-mu-L}}\preceq\left(1+\frac{\|\G_t-\H_t\|}{\mu^2}\right)\H_t \\
&\overset{\eqref{eq:G-H-tr}}\preceq\Bigl(1+\frac{\tau_{\H_t}(\G_t)}{\mu^2}\Bigr)\H_t\\
&\,\preceq\left(1+
\frac{mL^2\tau_{\H_t}(\G_t)}
{\mu^2\tr{\H_t}}\right)\H_t,
\end{split}
\end{align}
where the last step is based on the fact
\begin{align*}
    \tr{\H_t} \overset{\eqref{eq:H-mu-L}}\leq \tr{L^2\mI_m} = mL^2.
\end{align*}
Combining the choices of $\delta_t$ in the statement and results
\eqref{eq:delta-GH-1} and \eqref{eq:delta-GH-2},
we achieve 
\begin{align*}
\H_t\preceq\G_t\preceq(1+\delta_t)\H_t     
\end{align*}
for each of Options I--III.




We now consider the upper bound for $\mathbb E_t[\delta_{t+1}]$.
% Taking inverses in the first inequality of~\eqref{eq:gram_transport_lm} gives
% \begin{align*}
% \H_{t+1}^{-1}
% &\preceq(1+Mr_t)\H_t^{-1}.
% \end{align*}
% Since $\G_t\succeq\vzero$, taking the congruence transformation with
% $\G_t^{1/2}$ and then taking traces give
% \begin{align*}
% \tr{\H_{t+1}^{-1}\G_t}
% &=\tr{\G_t^{1/2}\H_{t+1}^{-1}\G_t^{1/2}}\\
% &\leq(1+Mr_t)\tr{\G_t^{1/2}\H_t^{-1}\G_t^{1/2}}\\
% &=(1+Mr_t)\tr{\H_t^{-1}\G_t}.
% \end{align*}
For Options I and II, the choice of~$\delta_t=\sigma_{\H_t}(\G_t)$ gives
\begin{align}
\label{eq:defdelta_tequal}
    \tr{\H_t^{-1}\G_t}
&=\tr{\H_t^{-1}(\G_t-\H_t)}+\tr{\I_m}
=\delta_t+m.
\end{align}
Consequently, it holds
\begin{align}
    \label{eq:sigma_HttildeGt}
    \begin{split}
\sigma_{\H_{t+1}}(\widetilde{\G}_t)
&={\rm tr}\big(\H_{t+1}^{-1}\widetilde{\G}_t\big)-m\\
&=\tr{\H_{t+1}^{-1}\cdot(1+Mr_t)\G_t}-m\\
&=(1+Mr_t)\tr{\H_{t+1}^{-1}\G_t}-m\\
&\!\overset{\eqref{eq:Gram_bound_0}}{\leq}
(1+Mr_t)^2\tr{\H_t^{-1}\G_t}-m\\
&\!\!\overset{\eqref{eq:defdelta_tequal}}{=}(1+Mr_t)^2(\delta_t+m)-m\\
&=(1+Mr_t)^2\delta_t
  +m\bigl((1+Mr_t)^2-1\bigr)\\
&\leq(1+Mr_t)^2\delta_t
  +2mMr_t(1+Mr_t)^2\\
&=(1+Mr_t)^2(\delta_t+2mMr_t),
    \end{split}
\end{align}
where the last inequality is due to the fact that $(1+a)^2-1\leq2a(1+a)^2$ holds for all $a\geq 0$.
Finally, Lemma~\ref{lm:QN_one_iter} gives
\begin{align*}
\mathbb E_t[\delta_{t+1}]
&\leq(1-\gamma)\sigma_{\H_{t+1}}(\widetilde{\G}_t)\overset{\eqref{eq:sigma_HttildeGt}}{\leq} (1-\gamma)(1+Mr_t)^2(\delta_t+2mMr_t),
\end{align*}
where $\gamma=\mu^2/(mL^2)$ for Option I and 
$\gamma=1/m$ for Option II.
This proves~equation \eqref{eq:delta-recur} with $\tilde c=1$ for Options I and II.

For Option III, Lemma~\ref{lm:QN_one_iter} indicates
\begin{align}
    \label{eq:tau_Ht+1_G}
    \begin{split}
\tau_{\H_{t+1}}(\widetilde{\G}_t)
&={\rm tr}\big(\widetilde{\G}_t-\H_{t+1}\big)\\
&=\tr{(1+Mr_t)\G_t-\H_{t+1}}\\
&\!\overset{\eqref{eq:Gram_bound_0}}{\leq}
\tr{(1+Mr_t)\G_t-\frac{1}{1+Mr_t}\H_t}\\
&=(1+Mr_t)\tau_{\H_t}(\G_t)
 +\left(1+Mr_t-\frac{1}{1+Mr_t}\right)\tr{\H_t}\\
&\overset{\text{($\dag$)}}{\leq}(1+Mr_t)\tau_{\H_t}(\G_t)+2Mr_t\tr{\H_t}\\
&\leq(1+Mr_t)
 \left(\frac{\tau_{\H_t}(\G_t)}{\tr{\H_t}}+2Mr_t\right)
 \tr{\H_t}\\
&\overset{\eqref{eq:Gram_bound_0}}{\leq}(1+Mr_t)^2
 \left(\frac{\tau_{\H_t}(\G_t)}{\tr{\H_t}}+2Mr_t\right)
 \tr{\H_{t+1}},
\end{split}
\end{align}
where the step ($\dag$) is due to the fact that $1+a-1/(1+a)
\leq2a$ holds for all $a\geq 0$.
Consequently, Lemma~\ref{lm:QN_one_iter} gives
\begin{align}\label{eq:tau_t+1bound}
\begin{split}    
\mathbb E_t[\tau_{\H_{t+1}}(\G_{t+1})] 
&\leq \left(1-\frac1m\right) 
\tau_{\H_{t+1}}(\widetilde{\G}_t) \\
&\!\overset{\eqref{eq:tau_Ht+1_G}}{\leq} 
\left(1-\frac1m\right)(1+Mr_t)^2
\left(\frac{\tau_{\H_t}(\G_t)}{\tr{\H_t}}+2Mr_t\right)
\tr{\H_{t+1}}.
\end{split}
\end{align}
According the choice of $\delta_t={mL^2\tau_{\H_t}(\G_t)}/(\mu^2\tr{\H_t})$ for Option III, we achieve
\begin{align*}
\mathbb E_t[\delta_{t+1}]
&=\frac{mL^2}{\mu^2\tr{\H_{t+1}}}
\mathbb E_t[\tau_{\H_{t+1}}(\G_{t+1})]\\
&\!\!\overset{\eqref{eq:tau_t+1bound}}{\leq}\frac{mL^2}{\mu^2\tr{\H_{t+1}}}
\left(1-\frac1m\right)(1+Mr_t)^2
\left(\frac{\tau_{\H_t}(\G_t)}{\tr{\H_t}}+2Mr_t\right)
\tr{\H_{t+1}}\\
&=\left(1-\frac1m\right)(1+Mr_t)^2
\left(\frac{mL^2\tau_{\H_t}(\G_t)}
{\mu^2\tr{\H_t}}+\frac{2mL^2Mr_t}{\mu^2}\right)\\
&=\left(1-\frac1m\right)(1+Mr_t)^2
\left(\delta_t+\frac{2mL^2Mr_t}{\mu^2}\right),
\end{align*}
which proves equation~\eqref{eq:delta-recur} with $\gamma=1/m$ and $\tilde c=L^2/\mu^2$ for Option III.

Finally, applying Lemma~\ref{lm:matrix_order_equivalence} in the view of
$\H_0\preceq\G_0\preceq\eta_0\H_0$ gives
\begin{align*}  
\mI_m\preceq\H_0^{-1/2}\G_0\H_0^{-1/2}\preceq\eta_0\mI_m,
\end{align*}
that is
\begin{align*}
\0
&\preceq\H_0^{-1/2}(\G_0-\H_0)\H_0^{-1/2}
\preceq(\eta_0-1)\I_m.
\end{align*}
Therefore, we achieve
\begin{align*}
\sigma_{\H_0}(\G_0)
={\rm tr}\big(\H_0^{-1/2}(\G_0-\H_0)\H_0^{-1/2}\big)
\leq m(\eta_0-1),
\end{align*}
and
\begin{align*}
\frac{mL^2\tau_{\H_0}(\G_0)}
{\mu^2\tr{\H_0}}
=\frac{mL^2\tr{\G_0-\H_0}}{\mu^2\tr{\H_0}}
\leq\frac{mL^2(\eta_0-1)\tr{\H_0}}{\mu^2\tr{\H_0}}
=\frac{mL^2(\eta_0-1)}{\mu^2}.
\end{align*}
\end{proof}


\subsection{Proof of Lemma~\ref{lm:linear-quadra_n<d}}
\label{app:proof_lemma_linear_quadratic_under}
\begin{proof}
Recall that for the case of $n<d$,
we have defined   
\begin{align*}
\H(\x)=\J(\x)\J(\x)^{\top}
\qquad\text{and}\qquad
\lambda_{\mathrm{under}}(\x) =\sqrt{\F(\x)^{\top}(\H(\x))^{-1}\F(\x)}.
\end{align*}
For all $\vx$ satisfies condition \eqref{eq:G_condi_2},
it holds
\begin{align*}
\|\F(\x)\|
&=\big\|\H(\x)^{1/2}\H(\x)^{-1/2}\F(\x)\big\|
\leq\big\|\H(\x)^{1/2}\big\|\lambda_{\mathrm{under}}(\x)
\\&=\Norm{\J(\x)}\lambda_{\mathrm{under}}(\x)\overset{\eqref{eq:lip}}{\leq}L\lambda_{\mathrm{under}}(\x)
\overset{\eqref{eq:G_condi_2}}{\leq}\frac{\mu^2}{2L_2},
\end{align*}
which implies 
\begin{align}\label{eq:x-in-under}
 \x\in\Omega_{\mathrm{under}}   
\end{align}
Moreover, the condition $\H(\x)\preceq\G$  gives
\begin{align}
\label{eq:boundr_lambda}
\|\x_{+}-\x\|^2
&\overset{\eqref{eq:qNupdate}}{=}
\F(\x)^{\top}\G^{-1}\H(\x)\G^{-1}\F(\x)\notag\leq\F(\x)^{\top}\G^{-1}\F(\x)\\
&\leq\F(\x)^{\top}\H(\x)^{-1}\F(\x)
 =(\lambda_{\mathrm{under}}(\x))^2.
\end{align}
On the other hand, it holds
\begin{align}
\label{eq:Fx_+Fx}
\begin{split}
&\F(\x_{+})\\
&=\F(\x)+\int_0^1
\J\bigl(\x+z(\x_{+}-\x)\bigr)(\x_{+}-\x)\,{\rm d}z\\
&=\F(\x)+\J(\x)(\x_{+}-\x)+\int_0^1\Bigl(\J\bigl(\x+z(\x_{+}-\x)\bigr)-\J(\x)\Bigr)
(\x_{+}-\x)\,{\rm d}z\\
&\overset{\eqref{eq:qNupdate}}{=}
\F(\x)-\J(\x)\J(\x)^{\top}\G^{-1}\F(\x)+\int_0^1\Bigl(\J\bigl(\x+z(\x_{+}-\x)\bigr)-\J(\x)\Bigr)
(\x_{+}-\x)\,{\rm d}z\\
&\overset{\eqref{eq:H-Gram}}{=}
(\I_n-\H(\x)\G^{-1})\F(\x)+\R_F,
\end{split}
\end{align}
where we define
\begin{align*}
\R_F
\triangleq\int_0^1\Bigl(\J\bigl(\x+z(\x_{+}-\x)\bigr)-\J(\x)\Bigr)
(\x_{+}-\x)\,{\rm d}z.    
\end{align*}
% Thus, the preceding expansion becomes
% \begin{align*}
% \F(\x_{+})=(\I-\H(\x)\G^{-1})\F(\x)+\R_F.
% \end{align*}
We also have
\begin{align}
\label{eq:R_Fnormbound}
\begin{split}
\|\R_F\|
&\leq\int_0^1
\Norm{\J\bigl(\x+z(\x_{+}-\x)\bigr)-\J(\x)}
\|\x_{+}-\x\|\,{\rm d}z\\
&\overset{\eqref{eq:lip}}{\leq}
L_2\|\x_{+}-\x\|^2\int_0^1z\,{\rm d}z
=\frac{L_2}{2}\|\x_{+}-\x\|^2
\overset{\eqref{eq:boundr_lambda}}{\leq}
\frac{L_2}{2}(\lambda_{\mathrm{under}}(\x))^2.
\end{split}
\end{align}
The expansion of $\F(\x_{+})$ gives
\begin{align}
\label{eq:HFx+bound}
\H(\x)^{-1/2}\F(\x_{+})
&\overset{\eqref{eq:Fx_+Fx}}{=}\bigl(\I_n-\H(\x)^{1/2}\G^{-1}\H(\x)^{1/2}\bigr)
\H(\x)^{-1/2}\F(\x)+\H(\x)^{-1/2}\R_F.
\end{align}
The result \eqref{eq:x-in-under}  and Proposition~\ref{prop:gram_lower} indicate
\begin{align}\label{eq:H12mu}
\big\|\H(\x)^{-1/2}\big\|\leq\frac{1}{\mu}.    
\end{align}
Taking inverses in~\eqref{eq:G_condi_2} gives
\begin{align*}
\frac{1}{\eta}\H(\x)^{-1}
&\preceq\G^{-1}\preceq\H(\x)^{-1}.
\end{align*}
Applying Lemma~\ref{lm:matrix_order_equivalence} to above inequality with
$\A=\H(\x)^{-1}$ and $\B=\G^{-1}$ gives
\begin{align*}
\frac{1}{\eta}\I_n
&\preceq\H(\x)^{1/2}\G^{-1}\H(\x)^{1/2}\preceq\I_n.
\end{align*}
Therefore, we have
\begin{align*}
\0
&\preceq\I_n-\H(\x)^{1/2}\G^{-1}\H(\x)^{1/2}
\preceq\left(1-\frac{1}{\eta}\right)\I_n,
\end{align*}
which implies
\begin{align}
\label{eq:normI-HGH}
\Norm{\I_n-\H(\x)^{1/2}\G^{-1}\H(\x)^{1/2}}
&\leq1-\frac{1}{\eta}.
\end{align}
Combining these bounds with the expansion of $\F(\x_{+})$ gives
\begin{align}
\label{eq:HFx+}
\Norm{\H(\x)^{-1/2}\F(\x_{+})}
&\!\!\overset{\eqref{eq:HFx+bound}}{\leq}\Norm{\I_n-\H(\x)^{1/2}\G^{-1}\H(\x)^{1/2}}
\Norm{\H(\x)^{-1/2}\F(\x)}
+\Norm{\H(\x)^{-1/2}}\|\R_F\|\notag\\
&\!\!\overset{\eqref{eq:H12mu}}\leq\Norm{\I_n-\H(\x)^{1/2}\G^{-1}\H(\x)^{1/2}}
\lambda_{\mathrm{under}}(\x)+\frac{\|\R_F\|}{\mu}\notag\\
&\!\!\!\!\!\!\overset{\eqref{eq:normI-HGH},\,\eqref{eq:R_Fnormbound}
}{\leq}\left(1-\frac{1}{\eta}\right)\lambda_{\mathrm{under}}(\x)
  +\frac{L_2}{2\mu}(\lambda_{\mathrm{under}}(\x))^2
\end{align}
and
\begin{align*}
\|\F(\x_{+})\|
&=\Norm{\H(\x)^{1/2}\H(\x)^{-1/2}\F(\x_{+})}\\
&\leq L\Norm{\H(\x)^{-1/2}\F(\x_{+})}\\
&\overset{\eqref{eq:HFx+}}{\leq}
L\lambda_{\mathrm{under}}(\x)
+\frac{LL_2}{2\mu}(\lambda_{\mathrm{under}}(\x))^2
\overset{\eqref{eq:G_condi_2}}{\leq}
\frac{\mu^2}{4L_2}+\frac{\mu^3}{32LL_2}
\leq\frac{\mu^2}{2L_2},
\end{align*}
where the first inequality is based on Proposition \ref{prop:gram_lip}.
Thus, both $\x$ and $\x_{+}$ belong to $\Omega_{\mathrm{under}}$.

Furthermore, the definition of $\lambda_{\mathrm{under}}(\cdot)$ gives
\begin{align*}
\lambda_{\mathrm{under}}(\x_{+})
&=\sqrt{\F(\x_{+})^{\top}\H(\x_{+})^{-1}\F(\x_{+})}\\
&\overset{\eqref{eq:Gram_bound_0}}{\leq}\sqrt{1+M\|\x_{+}-\x\|}
\sqrt{\F(\x_{+})^{\top}\H(\x)^{-1}\F(\x_{+})} \\
&=
\sqrt{1+M\|\x_{+}-\x\|} \,\big\|\H(\x)^{-1/2}\F(\x_{+})\big\|\\
&\!\!\!\!\!\!\!\overset{\eqref{eq:boundr_lambda},\,\eqref{eq:HFx+}}{\leq}
\sqrt{1+M\lambda_{\mathrm{under}}(\x)}\,
 \left(\left(1-\frac{1}{\eta}\right)\lambda_{\mathrm{under}}(\x)
 +\frac{L_2}{2\mu}(\lambda_{\mathrm{under}}(\x))^2\right)\\
&\leq(1+M\lambda_{\mathrm{under}}(\x))
 \left(\left(1-\frac{1}{\eta}\right)\lambda_{\mathrm{under}}(\x)
 +\frac{L_2}{2\mu}(\lambda_{\mathrm{under}}(\x))^2\right)\\
&=\left(1-\frac{1}{\eta}\right)\lambda_{\mathrm{under}}(\x)
+\left(\left(1-\frac{1}{\eta}\right)M+\frac{L_2}{2\mu}
+\frac{ML_2}{2\mu}\lambda_{\mathrm{under}}(\x)\right)
(\lambda_{\mathrm{under}}(\x))^2\\
&\!\!\overset{\eqref{eq:G_condi_2}}{\leq}\!\!
\left(1-\frac{1}{\eta}\right)\lambda_{\mathrm{under}}(\x)
+\left(\frac{2LL_2}{\mu^2}+\frac{L_2}{2\mu}
+\frac{L_2}{4\mu}\right)(\lambda_{\mathrm{under}}(\x))^2\\
&\leq\left(1-\frac{1}{\eta}\right)\lambda_{\mathrm{under}}(\x)
 +\frac{3LL_2}{\mu^2}(\lambda_{\mathrm{under}}(\x))^2,
\end{align*}
where $M=2LL_2/\mu^2$ follows Proposition \ref{prop:Gram_concor}.
This finishes the proof.
\end{proof}

\subsection{Proof of Theorem~\ref{thm:linear_under}}
\label{app:proof_theorem_linear_under}
\begin{proof}
For the ease of presentation, we denote
\begin{align*}
\lambda_{\mathrm{under},t}\triangleq\lambda_{\mathrm{under}}(\x_t).    
\end{align*}
We use induction to prove  $\x_t\in\Omega_{\mathrm{under}}$ and equation \eqref{eq:induc_lambda_n<d} holds for all $t\geq 0$.
For the induction base, the initialization in Line~\ref{line:input} of
Algorithm~\ref{alg:RQGN} and the definition of $\eta_0$ in equation \eqref{eq:dfn-eta0} give~$\H_0\preceq\G_0\preceq\eta_0\H_0$, which proves \eqref{eq:induc_lambda_n<d} in the case of $t=0$.
Additionally, condition~\eqref{eq:linear_initial_condi_2} gives
\begin{align*}
\lambda_{\mathrm{under},0}
&\leq\frac{\mu^2}{20\eta_0LL_2}
\leq\frac{\mu^2}{4LL_2}.
\end{align*}
Thus, applying Lemma~\ref{lm:linear-quadra_n<d} with $\eta=\eta_0$ gives $\x_0\in\Omega_{\mathrm{under}}$. 

For the induction step, we suppose  $\x_t\in\Omega_{\mathrm{under}}$ and equation~\eqref{eq:induc_lambda_n<d} hold for all $t\leq \hat t$.
Since the induction hypothesis implies $\lambda_{\mathrm{under},\hat t}\leq\lambda_{\mathrm{under},0}$, Lemma~\ref{lm:linear-quadra_n<d} implies 
\begin{align}\label{eq:r_hat_t_bound}
\x_{\hat t+1}\in\Omega_{\mathrm{under}}
\qquad\text{and}\qquad
r_{\hat t}\leq \lambda_{\mathrm{under},\hat t}
\end{align}
Furthermore, applying Proposition~\ref{prop:Gram_concor} to $\x_{\hat t}$ and
$\x_{\hat t+1}$ gives
\begin{align*}
\H_{\hat t+1}
&\overset{\eqref{eq:Gram_bound_0}}{\preceq}
(1+Mr_{\hat t})\H_{\hat t}
\overset{\eqref{eq:induc_lambda_n<d}}{\preceq}
(1+Mr_{\hat t})\G_{\hat t}
=\widetilde{\G}_{\hat t}
\end{align*}
and
\begin{align*}
\widetilde{\G}_{\hat t}
&=(1+Mr_{\hat t})\G_{\hat t}\\
&\overset{\eqref{eq:induc_lambda_n<d}}{\preceq}
\eta_0(1+Mr_{\hat t})
\exp\biggl(2\sum_{i=0}^{\hat t-1}Mr_i\biggr)\H_{\hat t}\\
&\overset{\eqref{eq:Gram_bound_0}}{\preceq}
\eta_0(1+Mr_{\hat t})^2
\exp\biggl(2\sum_{i=0}^{\hat t-1}Mr_i\biggr)\H_{\hat t+1}\\
&\preceq\eta_0\exp\biggl(2\sum_{i=0}^{\hat t}Mr_i\biggr)
\H_{\hat t+1},
\end{align*}
where the last step is based on the fact $1+x\leq {\rm e}^{x}$ with $x=Mr_{\hat{t}}$.
The above two results imply we can apply Lemma~\ref{lm:QN_order} with $\mG=\widetilde{\G}_{\hat t}$, $\mH=\H_{\hat t+1}$, and $\eta=\eta_0\exp\big(2\sum_{i=0}^{\hat t}Mr_i\big)$ to~achieve
\begin{align}
    \label{eq:G_tH_tleq2}
\H_{{\hat t}+1} &\preceq\G_{{\hat t}+1}
 \preceq\eta_0\exp\biggl(2\sum_{i=0}^{\hat t}Mr_i\biggr)\H_{{\hat t}+1}.
\end{align}
We also have
\begin{align}
    \label{eq:sum_r_t_bound}
\sum_{i=0}^{t-1}r_i
&\overset{\eqref{eq:r_hat_t_bound}}{\leq}
\sum_{i=0}^{t-1}\lambda_{\mathrm{under},i}
\overset{\eqref{eq:induc_lambda_n<d}}{\leq}
 \sum_{i=0}^{t-1}\biggl(1-\frac{1}{2\eta_0}\biggr)^i
 \lambda_{\mathrm{under},0}
 \leq2\eta_0\lambda_{\mathrm{under},0}
 \overset{\eqref{eq:linear_initial_condi_2}}{\leq}\frac{1}{5M},
\end{align}
where $M=2LL_2/\mu^2$.
Therefore, it holds
    \begin{align*}
\H_{{\hat t}+1}
&\preceq\G_{{\hat t}+1}
\overset{\eqref{eq:G_tH_tleq2}}{\preceq}
\eta_0\exp\biggl(2M\sum_{i=0}^{\hat t}r_i\biggr)\H_{\hat t+1}
\overset{\eqref{eq:sum_r_t_bound}}{\preceq}
\eta_0{\rm e}^{2/5}\H_{\hat t+1}
\preceq\frac{3\eta_0}{2}\H_{\hat t+1}.
\end{align*}
Additionally, Lemma~\ref{lm:linear-quadra_n<d} implies
\begin{align*}
\lambda_{\mathrm{under},\hat t+1}
&\overset{\eqref{eq:linear-quadra-n<d}}{\leq}
\biggl(1-\frac{2}{3\eta_0}
+\frac{3LL_2}{\mu^2}\lambda_{\mathrm{under},\hat t}\biggr)
\lambda_{\mathrm{under},\hat t}\overset{\eqref{eq:induc_lambda_n<d},\,\eqref{eq:linear_initial_condi_2}}{\leq}
\biggl(1-\frac{2}{3\eta_0}+\frac{3}{20\eta_0}\biggr)
\lambda_{\mathrm{under},\hat t}\\
&~\leq\biggl(1-\frac{1}{2\eta_0}\biggr)
\lambda_{\mathrm{under},\hat t}
\overset{\eqref{eq:induc_lambda_n<d}}{\leq}
\biggl(1-\frac{1}{2\eta_0}\biggr)^{\hat t+1}
\lambda_{\mathrm{under},0}.
\end{align*}
This completes the induction.
\end{proof}



\subsection{Proof of Theorem~\ref{thm:n<dsuperlinear}}
\label{app:proof_theorem_superlinaer_under}

\begin{proof} 
We follow the notation $\lambda_{\mathrm{under},t}=\lambda_{\mathrm{under}}(\x_t)$ defined in Appendix
\ref{app:proof_theorem_linear_under}.
Recall that the definitions of $\tilde c$ and $\gamma$ in Lemma \ref{lm:Gram_approximate} satisfy
 $\tilde c n\geq1$ and $1-\gamma<1$, then
 condition~\eqref{eq:initial_condi_n<d} implies equation \eqref{eq:linear_initial_condi_2} holds.
Thus, Theorem~\ref{thm:linear_under} holds naturally that we have $\x_t\in\Omega_{\mathrm{under}}$ and
\begin{align}
\label{eq:thm_linear_condi_imply}
    \lambda_{\mathrm{under},t}\leq \left(1-\frac{1}{2\eta_0}\right)^{t}\lambda_{\mathrm{under},0}\leq \lambda_{\mathrm{under},0}\leq \frac{\mu^2}{4LL_2}
\end{align}
Therefore, Lemma~\ref{lm:Gram_approximate} gives
\begin{align}\label{eq:HGdelta}
    \H_t\preceq \G_t\preceq (1+\delta_t)\H_t.
\end{align}
Hence, we apply Lemma~\ref{lm:linear-quadra_n<d} with $\vx = \vx_t$, $\vx_+=\vx_{t+1}$, and $\eta=1+\delta_t$ to achive
\begin{align}
    \label{eq:r_tboundedbylambda_t}
r_t=\norm{\vx_{t+1}-\vx_t}\leq\lambda_{\mathrm{under},t}
\end{align}
and
% By~\eqref{eq:def_eta} and the matrix bounds in~\eqref{eq:delta-recur},
% $\eta_t-1\leq\delta_t$.
\begin{align}
\label{eq:exp_lambda_n<d}
\lambda_{\mathrm{under},t+1}
&\!\!\overset{\eqref{eq:linear-quadra-n<d}}{\leq}
\biggl(1-\frac{1}{1+\delta_t}\biggr)\lambda_{\mathrm{under},t}
+\frac{3LL_2}{\mu^2}(\lambda_{\mathrm{under},t})^2\notag\\
& \leq\delta_t\lambda_{\mathrm{under},t}
+\frac{3LL_2}{\mu^2}(\lambda_{\mathrm{under},t})^2.
\end{align}
According to Lemma~\ref{lm:Gram_approximate} with $m=n$, it holds
\begin{align}
\label{eq:exp_delta_n<d}
\begin{split}    
\mathbb E_t[\delta_{t+1}]
&\!\!\overset{\eqref{eq:delta-recur}}{\leq}
(1-\gamma)(1+Mr_t)^2(\delta_t+2\tilde c nMr_t) \\
&\!\!\overset{\eqref{eq:r_tboundedbylambda_t}}{\leq}
(1-\gamma)(1+M\lambda_{\mathrm{under},t})^2
 (\delta_t+2\tilde c nM\lambda_{\mathrm{under},t}).
 \end{split}
\end{align}
Let
\begin{align}\label{eq:dfn-xi-t}
\xi_t=\delta_t + \frac{7\tilde c nLL_2\lambda_{\mathrm{under},t}}{\mu^2}.    
\end{align}
Recall that we have $M=2LL_2/\mu^2$ and $\tilde c n\geq1$, which implies
$2\tilde c nM\leq7\tilde c nLL_2/\mu^2$ and~$3LL_2/\mu^2\leq7\tilde c nLL_2/\mu^2$.
This leads to
\begin{align}
    \label{eq:xi_t_lower_bound}
\delta_t+2\tilde{c}nM\lambda_{\mathrm{under},t}\leq\xi_t\qquad\text{and}\qquad\delta_t+ \frac{3LL_2}{\mu^2}\lambda_{\mathrm{under},t}\leq\xi_t.
\end{align}
Combining equations~\eqref{eq:exp_lambda_n<d} and \eqref{eq:exp_delta_n<d}, we have
\begingroup
\small
\begin{align}
\label{eq:xi_t_one_iter_n<d}
\begin{split}    
&\mathbb E_t[\xi_{t+1}]
\\&\overset{\eqref{eq:dfn-xi-t}}{=}
\BE_t\Bigl[\delta_{t+1}
+\frac{7\tilde c nLL_2}{\mu^2}\lambda_{\mathrm{under},t+1}\Bigr]\\
&\!\!\!\!\!\overset{\eqref{eq:exp_delta_n<d},\,\eqref{eq:exp_lambda_n<d}}{\leq}
(1-\gamma)(1+M\lambda_{\mathrm{under},t})^2
\bigl(\delta_t+2\tilde{c}nM\lambda_{\mathrm{under},t}\bigr)+\frac{7\tilde c nLL_2}{\mu^2}
\left(\delta_t+\frac{3LL_2}{\mu^2}\lambda_{\mathrm{under},t}\right)
\lambda_{\mathrm{under},t}\\
&\overset{\eqref{eq:xi_t_lower_bound}}{\leq}
(1-\gamma)(1+M\lambda_{\mathrm{under},t})^2\xi_t
 +\frac{7\tilde c nLL_2}{\mu^2}\lambda_{\mathrm{under},t}\,\xi_t \\
&~\leq(1-\gamma)\biggl((1+M\lambda_{\mathrm{under},t})^2
+\frac{7\tilde c nLL_2}{(1-\gamma)\mu^2}
\lambda_{\mathrm{under},t}\biggr)\xi_t\\
&~\leq(1-\gamma)
\exp\Biggl(\biggl(2M+\frac{7\tilde c nLL_2}
{(1-\gamma)\mu^2}\biggr)\lambda_{\mathrm{under},t}\Biggr)\xi_t \\
&~\leq(1-\gamma)
 \exp\biggl(\frac{11\tilde c nLL_2}{(1-\gamma)\mu^2}\lambda_{\mathrm{under},t}\biggr)\xi_t,
\end{split}
\end{align}
\endgroup
where we use the fact $(1+a)^2+b\leq (1+a)^2(1+b)\leq \exp(2a+b)$ for all $a,b\geq 0$ by taking~$a = (1+M\lambda_{\mathrm{under},t})\geq 0$ and $b = 7\tilde{c}nLL_2\lambda_{\mathrm{under},t}/((1-\gamma)\mu^2)\geq 0$ in the second last inequality and the definition $M=2LL_2/\mu^2\leq 2\tilde{c}n/((1-\gamma)\mu^2)$ in the last inequality.
% \begin{align}
% \label{eq:xi_t_one_iter_n<d}
% \mathbb E_t[\xi_{t+1}]
% &\!\!\!\!\!\!\overset{\eqref{eq:exp_delta_n<d},\eqref{eq:exp_lambda_n<d},\eqref{eq:xi_t_lower_bound}}{\leq}
% (1-\gamma)(1+M\lambda_{\mathrm{under},t})^2\xi_t
%  +\frac{7\tilde c nLL_2}{\mu^2}\lambda_{\mathrm{under},t}\xi_t\notag\\
% &~~=(1-\gamma)\biggl((1+M\lambda_{\mathrm{under},t})^2
% +\frac{7\tilde c nLL_2}{(1-\gamma)\mu^2}
% \lambda_{\mathrm{under},t}\biggr)\xi_t\notag\\
% &~~\leq(1-\gamma)
% \exp\Biggl(\biggl(2M+\frac{7\tilde c nLL_2}
% {(1-\gamma)\mu^2}\biggr)\lambda_{\mathrm{under},t}\Biggr)\xi_t\notag\\
% &~~\leq(1-\gamma)
%  \exp\biggl(\frac{11\tilde c nLL_2}{(1-\gamma)\mu^2}\lambda_{\mathrm{under},t}\biggr)\xi_t.
% \end{align}

Taking expectation on~\eqref{eq:xi_t_one_iter_n<d} and applying equation \eqref{eq:thm_linear_condi_imply} gives
\begin{align}
\label{eq:recur_xi_t_n<d}
\EBP{\xi_{t+1}}
&\overset{\eqref{eq:xi_t_one_iter_n<d},\, \eqref{eq:thm_linear_condi_imply}}{\leq}
(1-\gamma)
 \exp\Biggl(\frac{11\tilde c nLL_2}{(1-\gamma)\mu^2}
 \biggl(1-\frac{1}{2\eta_0}\biggr)^t\lambda_{\mathrm{under},0}\Biggr)
 \EBP{\xi_t}.
\end{align}
Recall that the geometric series satisfies
\begin{align}
\label{eq:geosum_boun}
    \sum_{i=0}^{t-1}\biggl(1-\frac{1}{2\eta_0}\biggr)^i
&\leq\sum_{i=0}^{\infty}\biggl(1-\frac{1}{2\eta_0}\biggr)^i
=2\eta_0.
\end{align}
Hence, the recursion~\eqref{eq:recur_xi_t_n<d} gives
\begin{align}
\label{eq:xi_t_bound_n<d}
\EBP{\xi_t}
&\overset{\eqref{eq:recur_xi_t_n<d}}{\leq}(1-\gamma)^t
 \prod_{i=0}^{t-1}
 \exp\Biggl(\frac{11\tilde c nLL_2}{(1-\gamma)\mu^2}
 \biggl(1-\frac{1}{2\eta_0}\biggr)^i
 \lambda_{\mathrm{under},0}\Biggr)\xi_0\notag\\
&=(1-\gamma)^t
 \exp\Biggl(\frac{11\tilde c nLL_2}{(1-\gamma)\mu^2}
 \lambda_{\mathrm{under},0}
 \sum_{i=0}^{t-1}\biggl(1-\frac{1}{2\eta_0}\biggr)^i\Biggr)
 \xi_0\notag\\
&\overset{\eqref{eq:geosum_boun}}{\leq}(1-\gamma)^t
 \exp\biggl(\frac{22\tilde c nLL_2\eta_0}{(1-\gamma)\mu^2}\lambda_{\mathrm{under},0}\biggr)
 \xi_0\notag\\
&\overset{\eqref{eq:initial_condi_n<d}}{\leq}
(1-\gamma)^t {\rm e}^{11/10}
 \biggl(\delta_0+\frac{7(1-\gamma)}{20\eta_0}\biggr)
 =(1-\gamma)^tq_{\text{\rm under}}.
\end{align}
Additionally, the definition of $\xi_t$ gives
\begin{align}\label{eq:lambda-ratio}
\frac{\lambda_{\mathrm{under},t+1}}{\lambda_{\mathrm{under},t}}
\overset{\eqref{eq:exp_lambda_n<d}}{\leq}\delta_t+\frac{3LL_2\lambda_{\mathrm{under},t}}{\mu^2}\leq\delta_t+\frac{7\tilde c nLL_2\lambda_{\mathrm{under},t}}{\mu^2}
=\xi_t,
\end{align}
The definition of $\eta_t$ in equation \eqref{eq:def_eta} and equation \eqref{eq:HGdelta} indicates 
\begin{align}\label{eq:eta-xi}
\eta_t\leq\delta_t+1\leq\xi_t+1,     
\end{align}
where the last step is based on equation \eqref{eq:lambda-ratio}.
Therefore, we have
\begin{align}\label{eq:lambda_{t+1}/lambda_t}
\EBP{\eta_t-1}
\overset{\eqref{eq:eta-xi}}\leq\EBP{\xi_t}
\overset{\eqref{eq:xi_t_bound_n<d}}{\leq}
(1-\gamma)^tq_{\text{\rm under}}
~~\text{and}~~
\EBP{\frac{\lambda_{\mathrm{under},t+1}}{\lambda_{\mathrm{under},t}}}\leq\EBP{\xi_t}
\overset{\eqref{eq:xi_t_bound_n<d}}{\leq}
(1-\gamma)^tq_{\text{\rm under}}.
\end{align}
Combing equation \eqref{eq:lambda_{t+1}/lambda_t} and Lemma \ref{lm:lin_lemma26} with
$X_t=\lambda_{\mathrm{under},t+1}/\lambda_{\mathrm{under},t}$,
$a=q_{\text{\rm under}}$, and $\tau=1/\gamma>1$, with probability at least $1-p$, we have
\begin{align}
\label{eq:lambda_t+1-recurn<d}
\frac{\lambda_{\mathrm{under},t+1}}{\lambda_{\mathrm{under},t}}
&\leq \biggl(1-\frac{\gamma}{1+\gamma}\biggr)^t\frac{q_{\text{\rm under}}}{p\gamma^2}
\end{align}
simultaneously for every $t\geq0$.
Recall that $\gamma\in(0,1)$ and $\log(1+\gamma)\geq\gamma/2$, then taking
\begin{align*}
t_0
=\biggl\lceil
\frac{\log\bigl(1+2q_{\text{\rm under}}/(p\gamma^2)\bigr)}{\log(1+\gamma)}
\biggr\rceil=\OM\biggl(\frac{\log\bigl(1+q_{\text{\rm under}}/(p\gamma^2)\bigr)}{\gamma}\biggr)
\end{align*}
gives
\begin{align}\label{eq:superlinearleq1/2n<d}
\biggl(1-\frac{\gamma}{1+\gamma}\biggr)^{t_0}\frac{q_{\text{\rm under}}}{p\gamma^2}
&\leq\frac{1}{2}.
\end{align}
Therefore, for all $t\geq1$, we have
\begin{align*}
\lambda_{\mathrm{under},t_0+t}
&=\lambda_{\mathrm{under},t_0}
\prod_{i=0}^{t-1}
\frac{\lambda_{\mathrm{under},t_0+i+1}}
{\lambda_{\mathrm{under},t_0+i}}\\
&\!\!\!\!\!\!\!\overset{\eqref{eq:lambda_t+1-recurn<d},\, \eqref{eq:superlinearleq1/2n<d}}{\leq}\lambda_{\mathrm{under},t_0}
\prod_{i=0}^{t-1}
\biggl(\frac{1}{2}
\biggl(1-\frac{\gamma}{1+\gamma}\biggr)^i\biggr)\\
&=\biggl(1-\frac{\gamma}{1+\gamma}\biggr)^{\sum_{i=0}^{t-1}i}
\biggl(\frac{1}{2}\biggr)^t\lambda_{\mathrm{under},t_0}\\
&=\biggl(1-\frac{\gamma}{1+\gamma}\biggr)^{t(t-1)/2}
\biggl(\frac{1}{2}\biggr)^t\lambda_{\mathrm{under},t_0}\\
&\!\!\overset{\eqref{eq:induc_lambda_n<d}}{\leq}
\biggl(1-\frac{\gamma}{1+\gamma}\biggr)^{t(t-1)/2}
 \biggl(\frac{1}{2}\biggr)^t
 \biggl(1-\frac{1}{2\eta_0}\biggr)^{t_0}\lambda_{\mathrm{under},0}
\end{align*}
holds with probability at least $1-p$.
\end{proof}


\subsection{Proof of Proposition~\ref{prop:local_nonsingularity}}
\label{app:proof_prop_local_nonsingularity}
\begin{proof}
For all $\x\in\Omega_{\mathrm{over}}$, the definition of
$\Omega_{\mathrm{over}}$ and Assumption~\ref{ass:nonsingular_J^*} indicate
\begin{align}
    \label{eq:local_nonsingularity_condi}
\|\x-\x_*\|
&\leq\frac{\mu_*^2}{4LL_2}
\qquad\text{and}\qquad
\H(\x_*)=\J(\x_*)^{\top}\J(\x_*)\succeq\mu_*^2\I_d.
\end{align}
Since Proposition~\ref{prop:gram_lip} gives $\H(\cdot)$ is $2LL_2$-Lipschitz continuous, it holds
\begin{align*}
\H(\x)
\succeq\H(\x_*)-2LL_2\|\x-\x_*\|\I_d\overset{\eqref{eq:local_nonsingularity_condi}}{\succeq}\mu_*^2\I_d-2LL_2\cdot\frac{\mu_*^2}{4LL_2}\I_d
=\frac{\mu_*^2}{2}\I_d.
\end{align*}
Thus, Assumption~\ref{ass:nonsingular} holds with $\Omega=\Omega_{\mathrm{over}}$ and
$\mu=\mu_*/\sqrt{2}$.
\end{proof}

\subsection{Proof of Proposition~\ref{prop:assumption-sati}}
\label{app:proof_prop}
\begin{proof}
We consider three conditions in this proposition as follows.
\paragraph{Case 1:}
 For all $\vx\in\BR^d$ such that $\|\x-\x_*\|\leq\delta$,  Proposition~\ref{prop:gram_lip} gives
\begin{align}
    \label{eq:F_i_bound}
    \begin{split}
\sum_{i=1}^n|F_i(\x)|
&\leq\|\F(\x_*)\|_1+\|\F(\x)-\F(\x_*)\|_1\\
&\leq\|\F(\x_*)\|_1+\sqrt n\|\F(\x)-\F(\x_*)\|\\
&\overset{\eqref{eq:lip}}{\leq}\|\F(\x_*)\|_1+\sqrt nL\delta.
    \end{split}
\end{align}
Using the fact $\nabla^2\phi(\x)=\H(\x)+\sum_{i=1}^nF_i(\x)\nabla^2F_i(\x)$, we achieve
\begin{align*}
&\nabla^2\phi(\x)-\nabla^2\phi(\x_*)\\
&=\H(\x)-\H(\x_*)
+\sum_{i=1}^n\bigl(F_i(\x)-F_i(\x_*)\bigr)\nabla^2F_i(\x_*)+\sum_{i=1}^nF_i(\x)
\bigl(\nabla^2F_i(\x)-\nabla^2F_i(\x_*)\bigr)
\end{align*}
which implies
\begin{align}
\label{eq:nabla_2phi_bound}
\begin{split}    
&\|\nabla^2\phi(\x)-\nabla^2\phi(\x_*)\|\\
&\leq\|\H(\x)-\H(\x_*)\|
+\sum_{i=1}^n|F_i(\x)-F_i(\x_*)|\,\|\nabla^2F_i(\x_*)\| \\
&\quad+\sum_{i=1}^n|F_i(\x)|\,
\|\nabla^2F_i(\x)-\nabla^2F_i(\x_*)\|\\
&\overset{\eqref{eq:lip}}{\leq}\left(2LL_2+L\sum_{i=1}^n\|\nabla^2F_i(\x_*)\|
+L_3\sum_{i=1}^n|F_i(\x)|\right)\|\x-\x_*\|\\
&\overset{\eqref{eq:F_i_bound}}{\leq}\left(2LL_2+L\sum_{i=1}^n\|\nabla^2F_i(\x_*)\|
+L_3\sqrt nL\delta+L_3\|\F(\x_*)\|_1\right)\|\x-\x_*\|\\
&=2\rho\|\x-\x_*\|,
\end{split}
\end{align}
where $\rho$ follows equation \eqref{eq:rho-cond-1}.
Here, the second inequality also uses the facts that the Gram matrix $\H(\cdot)$ is $2LL_2$-Lipschitz continuous from Proposition \ref{prop:gram_lip} and the condition that each~$\nabla^2F_i(\cdot)$ is $L_3$-Lipschitz continuous. 
Since we have $\nabla\phi(\x_*)=\J(\x_*)^{\top}\F(\x_*)=\0$, it holds
\begin{align}
\label{eq:sati-1}
\begin{split}
&\Norm{\J(\x)^{\top}\F(\x)-(\mH(\x_*)+\mQ(\x_*))(\x-\x_*)} \\
&= \Norm{\nabla\phi(\vx)-\nabla\phi(\vx^*)-\nabla^2\phi(\x^*)(\x-\x^*)}\\
&=\Norm{\int_{0}^{1}
\bigl(\nabla^2\phi(\x_*+s(\x-\x_*))-\nabla^2\phi(\x_*)\bigr)
(\x-\x_*)\,{\rm d}s}\\
&\!\!\overset{\eqref{eq:nabla_2phi_bound}}{\leq}\int_0^1 2\rho s\|\x-\x_*\|^2\,{\rm d}s
=\rho\|\x-\x_*\|^2,
\end{split}
\end{align}
which verifies Assumption~\ref{ass:x_star_philip}.

\paragraph{Case 2:} Condition \eqref{eq:phi-2-Lip-xy} implies that 
\begin{align*}
\|\nabla^2\phi(\x)-\nabla^2\phi(\x_*)\|
\leq\rho_0\|\x-\x_*\|    
\end{align*}
holds for all $\vx\in\BR^d$, where $\vx^*\in\BR^d$ follows the definition in Assumption \ref{ass:x_star_philip}.
In the view of equation \eqref{eq:nabla_2phi_bound}, we can follow the derivation of equation \eqref{eq:sati-1} by replacing $\rho$ with $\rho_0/2$ to verify 
Assumption~\ref{ass:x_star_philip} with $\rho=\rho_0/2$.

\paragraph{Case 3:} 
Note that the condition $\F(\x_*)=\0$ implies $\vx^*$ is a stationary point of $\phi(\cdot)$. 
Consequently, the form of
$\nabla^2\phi(\cdot)$ gives
\begin{align}
    \label{eq:zerox_*_property}
\nabla^2\phi(\x_*)=\H(\x_*)+\sum_{i=1}^nF_i(\x_*)\nabla^2F_i(\x_*)=\H(\x_*)=\J(\x_*)^{\top}\J(\x_*).
\end{align}
Moreover, the fundamental theorem of calculus gives
\begin{align}
    \label{eq:F_diff_bound}
    \begin{split}
&\Norm{\F(\x)-\F(\x_*)-\J(\x_*)(\x-\x_*)}\\
&=\Norm{\int_0^1
\bigl(\J(\x_*+s(\x-\x_*))-\J(\x_*)\bigr)
(\x-\x_*)\,{\rm d}s}\\
&\overset{\eqref{eq:lip}}{\leq}
\int_0^1L_2s\|\x-\x_*\|^2\,{\rm d}s
=\frac{L_2}{2}\|\x-\x_*\|^2.
    \end{split}
\end{align}
Proposition~\ref{prop:gram_lip} also gives
\begin{align}
    \label{eq:normFxbounded}
\|\F(\x)\|
&=\|\F(\x)-\F(\x_*)\|
\leq L\|\x-\x_*\|.
\end{align}
Thus, it holds
\begin{align*}
&\big\|\J(\x)^{\top}\F(\x)-\nabla^2\phi(\x_*)((\x-\x_*)\big\|\\
&\overset{\eqref{eq:zerox_*_property}}{=}\big\|\J(\x)^{\top}\F(\x)-\J(\x_*)^{\top}\F(\x_*)
-\H(\x_*)(\x-\x_*)\big\|\\
&~=\big\|\big(\J(\x)^{\top}-\J(\x_*)^{\top}\big)\F(\x)
+\J(\x_*)^{\top}
(\F(\x)-\F(\x_*)-\J(\x_*)(\x-\x_*))\big\|\\
&~\leq\big\|\J(\x)^{\top}-\J(\x_*)^{\top}\big\|\,\|\F(\x)\|
+\|\J(\x_*)\|\,
\|\F(\x)-\F(\x_*)-\J(\x_*)(\x-\x_*)\|\\
&\!\!\!\overset{\eqref{eq:lip},\eqref{eq:normFxbounded}}{\leq} LL_2\|\x-\x_*\|^2+ \|\J(\x_*)\|\,
\|\F(\x)-\F(\x_*)-\J(\x_*)(\x-\x_*)\|\\
&\!\!\!\overset{\eqref{eq:lip},\eqref{eq:F_diff_bound}}{\leq}
\left(LL_2+\frac{LL_2}{2}\right)\|\x-\x_*\|^2
=\frac{3LL_2}{2}\|\x-\x_*\|^2,
\end{align*}
which verifies Assumption~\ref{ass:x_star_philip} with $\rho = {3LL_2}/{2}$.
\end{proof}

\subsection{Proof of Lemma~\ref{lm:pd_approximate_phi}}
\label{app:proof_lemma_pd_approximate_phi}
\begin{proof}
We first prove equation~\eqref{eq:approx_phi_positive}.
The condition $c\in(1,1/\theta)$ indicates $(c-1)/c<1$, 
which implies for all $\vx\in\BR^d$ such that $\|\x-\x_*\|\leq(c-1)\mu_*^2/(4cLL_2)$, it holds
\begin{align*}
\x\in\Omega_{\mathrm{over}}.     
\end{align*}
Since matrix $\Q(\x_*)$ is symmetric and Assumption~\ref{ass:nonsingular_J^*} gives $\H(\x_*)\succ\0$, 
it holds 
\begin{align*}
\big|\v^{\top}\H(\x_*)^{-1/2}\Q(\x_*)
\H(\x_*)^{-1/2}\v\big|
&\leq\big\|\H(\x_*)^{-1/2}\Q(\x_*)\H(\x_*)^{-1/2}\big\|
\|\v\|^2\overset{\eqref{eq:condiDelta}}=\theta\|\v\|^2.
\end{align*}
This implies
\begin{align*}
-\theta\I_d
&\preceq\H(\x_*)^{-1/2}\Q(\x_*)\H(\x_*)^{-1/2}
\preceq\theta\I_d.
\end{align*}
Combining the above result and Lemma~\ref{lm:matrix_order_equivalence} with
$\A=\H(\x_*)$, $\B=\Q(\x_*)$, $a_1=-\theta$, and~$a_2=\theta$, we achieve
\begin{align}
    \label{eq:Qx_star_bound_H_x_star}
-\theta \H(\x^*)\preceq \Q(\x^*)\preceq \theta \H(\x^*).
\end{align}
Therefore, Propositions \ref{prop:Gram_concor} and \ref{prop:local_nonsingularity} imply
\begin{align*}
-\theta\bigl(1+M\|\x-\x_*\|\bigr)\H(\x)
\overset{\eqref{eq:Gram_bound_0}}{\preceq}-\theta\H(\x_*)
\overset{\eqref{eq:Qx_star_bound_H_x_star}}{\preceq}\Q(\x_*)
\overset{\eqref{eq:Qx_star_bound_H_x_star}}{\preceq}\theta\H(\x_*)\overset{\eqref{eq:Gram_bound_0}}{\preceq}\theta\bigl(1+M\|\x-\x_*\|\bigr)\H(\x)
\end{align*}
for all $\vx\in\BR^d$, where $M=2LL_2/\mu^2$ and $\mu=\mu_*/\sqrt{2}$.
Additionally, for all $c>1$, the condition of $\|\x-\x_*\|\leq(c-1)\mu_*^2/(4cLL_2)$  in the statement implies
\begin{align*}
1+M\|\x-\x_*\|
&\leq1+\frac{2LL_2}{\mu^2}\cdot
\frac{(c-1)\mu_*^2}{4cLL_2}
=2-\frac{1}{c}\leq c,
\end{align*}
which indicates
\begin{align}
    \label{eq:Qx_star_bound}
-c\theta\H(\x)
&\preceq\Q(\x_*)\preceq c\theta\H(\x).
\end{align}
Recall that Propositions~\ref{prop:gram_lip} and~\ref{prop:gram_lower} give
\begin{align}\label{eq:mu^2-H-L}
\mu^2\I_d\preceq\H(\x)\preceq L^2\I_d.    
\end{align}
Thus, for matrix $\tilde{\nabla}^2\phi(\x) = \H(\x)+\Q(\x^*)$, it holds
\begin{align}
    \label{eq:tildenabla^2phi_bound}
(1-c\theta)\mu^2\I_d
\overset{\eqref{eq:mu^2-H-L}}\preceq(1-c\theta)\H(\x)
\overset{\eqref{eq:Qx_star_bound}}{\preceq}\tilde{\nabla}^2\phi(\x)\overset{\eqref{eq:Qx_star_bound}}{\preceq}(1+c\theta)\H(\x)
\overset{\eqref{eq:mu^2-H-L}}\preceq(1+c\theta)L^2\I_d
\end{align}
for all $c\in(1,1/\theta)$,
which proves~equation \eqref{eq:approx_phi_positive}.

We then prove equation~\eqref{eq:concor_phi}.
Recall that Proposition \ref{prop:gram_lip} indicates $\H(\cdot)$ is $2LL_2$-Lipschitz continuous, which leads to
\begin{align}
\label{eq:Hy-Hxboundebydistance}
-2LL_2\|\y-\x\|\I_d
&\preceq\H(\y)-\H(\x)
=\tilde{\nabla}^2\phi(\y)-\tilde{\nabla}^2\phi(\x)
\preceq2LL_2\|\y-\x\|\I_d.
\end{align}
We also have
\begin{align}
\label{eq:distanceboundbytildeHessian}
2LL_2\|\y-\x\|\I_d
\overset{\eqref{eq:tildenabla^2phi_bound}}{\preceq}\frac{2LL_2\|\y-\x\|}{(1-c\theta)\mu^2}
\tilde{\nabla}^2\phi(\x)=\widehat M\|\y-\x\|\tilde{\nabla}^2\phi(\x),
\end{align}
where $\widehat M= 2LL_2/((1-c\theta)\mu^2)$.
Therefore, it holds
\begin{align*}
\tilde{\nabla}^2\phi(\y)
&\overset{\eqref{eq:Hy-Hxboundebydistance}}{\preceq}\tilde{\nabla}^2\phi(\x)+2LL_2\|\y-\x\|\I_d
\overset{\eqref{eq:distanceboundbytildeHessian}}{\preceq}(1+\widehat M\|\y-\x\|)\tilde{\nabla}^2\phi(\x)
\end{align*}
and
\begin{align*}
\tilde{\nabla}^2\phi(\y)
&\overset{\eqref{eq:Hy-Hxboundebydistance}}{\succeq}\tilde{\nabla}^2\phi(\x)-2LL_2\|\y-\x\|\I_d
\overset{\eqref{eq:distanceboundbytildeHessian}}{\succeq}(1-\widehat M\|\y-\x\|)\tilde{\nabla}^2\phi(\x),
\end{align*}
which proves~\eqref{eq:concor_phi}.
\end{proof}

\subsection{Proof of Lemma~\ref{lm:GNlinearquadra}}
\label{app:proof_lemma_linear_quadratic_over}
\begin{proof}
For the ease of presentation, we denote
\begin{align}\label{eq:dfn-R}
\R\triangleq\nabla\phi(\x)-\nabla^2\phi(\x_*)(\x-\x_*).    
\end{align}
The expressions of $\tilde{\nabla}^2\phi(\cdot)$ and ${\nabla}^2\phi(\cdot)$ implies
\begin{align*}
\tilde{\nabla}^2\phi(\x)=\nabla^2\phi(\x_*)+\H(\x)-\H(\x_*).  
\end{align*}
Combing with the update rule~\eqref{eq:qNupdate} and the definition of $\mR$, we achieve
\begin{align*}
\x_{+}-\x_*
&=\x-\x_*-\G^{-1}\nabla\phi(\x)\\
&=(\I_d-\G^{-1}\tilde{\nabla}^2\phi(\x))(\x-\x_*)
  +\G^{-1}(\H(\x)-\H(\x_*))(\x-\x_*)-\G^{-1}\R.
\end{align*}
Left multiplying matrix $(\tilde{\nabla}^2\phi(\x))^{1/2}$ on both sides of above equation gives
\begin{align}
\label{eq:splithatlambda}
(\tilde{\nabla}^2\phi(\x))^{1/2}(\x_{+}-\x_*)
&=\u+\v-\w,
\end{align}
where
\begin{align}
\u
&\triangleq\bigl(\I_d-\bigl(\tilde{\nabla}^2\phi(\x)\bigr)^{1/2}
\G^{-1}\bigl(\tilde{\nabla}^2\phi(\x)\bigr)^{1/2}\bigr)
\bigl(\tilde{\nabla}^2\phi(\x)\bigr)^{1/2}(\x-\x_*), \label{dfn:proof-u}\\
\v
&\triangleq\bigl(\tilde{\nabla}^2\phi(\x)\bigr)^{1/2}\G^{-1}
(\H(\x)-\H(\x_*))(\x-\x_*), \label{dfn:proof-v}\\
\w
&\triangleq\bigl(\tilde{\nabla}^2\phi(\x)\bigr)^{1/2}\G^{-1}\R. \label{dfn:proof-w}
\end{align}
We provide upper bounds for $\u$, $\vv$ and $\vw$ as follows.
\begin{itemize}
\item \textbf{The upper bound for $\|\u\|$.}
The condition~\eqref{eq:condi_x_linear_quadra} for matrix $\G$ gives
\begin{align*}
\frac{1}{\eta}(\H(\x))^{-1}
&\preceq\G^{-1}\preceq(\H(\x))^{-1}.
\end{align*}
According to Lemma \ref{lm:pd_approximate_phi},
taking inverse on equation~\eqref{eq:approx_phi_positive} gives
\begin{align*}
(1-c\theta)\bigl(\tilde{\nabla}^2\phi(\x)\bigr)^{-1}
\preceq (\H(\x))^{-1}
\preceq (1+c\theta)\bigl(\tilde{\nabla}^2\phi(\x)\bigr)^{-1},
\end{align*}
where $c\in(1,1/\theta)$ and $\theta$ follows definition \eqref{eq:condiDelta}.
The above results yields
\begin{align}
\label{eq:boundGinverse}
\frac{1-c\theta}{\eta}\bigl(\tilde{\nabla}^2\phi(\x)\bigr)^{-1}
\preceq \frac{1}{\eta}(\H(\x))^{-1}\preceq\G^{-1}\preceq (\H(\x))^{-1}
\preceq(1+c\theta)\bigl(\tilde{\nabla}^2\phi(\x)\bigr)^{-1}.
\end{align}
Applying Lemma~\ref{lm:matrix_order_equivalence}
to equation \eqref{eq:boundGinverse} 
with $\A=(\tilde{\nabla}^2\phi(\x))^{-1}$, 
$\B=\G^{-1}$, $a_1=(1-c\theta)/\eta$, and $a_2=1+c\theta$
gives
\begin{align}
\label{eq:upperphiGphi}
\frac{1-c\theta}{\eta}\I_d
\preceq\bigl(\tilde{\nabla}^2\phi(\x)\bigr)^{1/2}\G^{-1}
\bigl(\tilde{\nabla}^2\phi(\x)\bigr)^{1/2}
\preceq(1+c\theta)\I_d.
\end{align}
Recall that the setting $\eta\geq 1$ implies $c\theta\leq1-(1-c\theta)/\eta$, which leads to
\begin{align}\label{eq:bound_u}
\begin{split}
\|\u\|
& \overset{\eqref{dfn:proof-u},\,\eqref{eq:our_measure}}\leq\Bigl\|\I_d-\big(\tilde{\nabla}^2\phi(\x)\big)^{1/2}\G^{-1}
\big(\tilde{\nabla}^2\phi(\x)\big)^{1/2}\Bigr\|
\lambda_{\mathrm{over}}(\x) \\
&~~\,\overset{\eqref{eq:upperphiGphi}}{\leq}~~\,
\left(1-\frac{1-c\theta}{\eta}\right)\lambda_{\mathrm{over}}(\x).
\end{split}
\end{align}
\item \textbf{The upper bound for $\|\v\|$.}
The definition of $\lambda_{\mathrm{over}}(\cdot)$ in equation
\eqref{eq:our_measure} and Lemma~\ref{lm:pd_approximate_phi}
imply
\begin{align}
\label{eq:distance_lambda}
\|\x-\x_*\|
\leq\Norm{\bigl(\tilde{\nabla}^2\phi(\x)\bigr)^{-1/2}}
\Norm{\bigl(\tilde{\nabla}^2\phi(\x)\bigr)^{1/2}(\x-\x_*)}
\overset{\eqref{eq:approx_phi_positive}}{\leq}
\frac{\lambda_{\mathrm{over}}(\x)}{\sqrt{1-c\theta}\mu},
\end{align}
where $c\in(1,1/\theta)$ and $\theta$ follows definition \eqref{eq:condiDelta}.
Based on the $2LL_2$-Lipschitz continuity of $\H(\cdot)$ from
Proposition~\ref{prop:gram_lip} and Lemma~\ref{lm:pd_approximate_phi}, we obtain
\begin{align}
\|\v\|
&\overset{\eqref{dfn:proof-v}}\leq\Norm{\bigl(\tilde{\nabla}^2\phi(\x)\bigr)^{1/2}\G^{-1}
\bigl(\tilde{\nabla}^2\phi(\x)\bigr)^{1/2}}
\Norm{\bigl(\tilde{\nabla}^2\phi(\x)\bigr)^{-1/2}}\cdot\Norm{\H(\x)-\H(\x_*)}\|\x-\x_*\|\nonumber\\
&\,\,\leq \Norm{\bigl(\tilde{\nabla}^2\phi(\x)\bigr)^{1/2}\G^{-1}
\bigl(\tilde{\nabla}^2\phi(\x)\bigr)^{1/2}}
\Norm{\bigl(\tilde{\nabla}^2\phi(\x)\bigr)^{-1/2}} \cdot 2LL_2\|\x-\x^*\|^2\nonumber
\\
&\!\!\!\!\!\overset{\eqref{eq:upperphiGphi}, \eqref{eq:approx_phi_positive}}{\leq}
\frac{2(1+c\theta)LL_2}{\sqrt{1-c\theta}\mu}\|\x-\x_*\|^2\nonumber\\
&\!\!\overset{\eqref{eq:distance_lambda}}{\leq}
\frac{2(1+c\theta)LL_2}{(1-c\theta)^{3/2}\mu^3}
\bigl(\lambda_{\mathrm{over}}(\x)\bigr)^2.\label{eq:bound_v}
\end{align}
\item \textbf{The upper bound for $\|\w\|$.}
Recall that equation \eqref{eq:boundGinverse} gives
\begin{align*}
\frac{1-c\theta}{\eta}\G\preceq\tilde{\nabla}^2\phi(\x)
\preceq(1+c\theta)\G,
\end{align*}
which implies
\begin{align}
\label{eq:boundGinverse2}
\frac{1-c\theta}{\eta}\G^{-1}\preceq\G^{-1}\tilde{\nabla}^2\phi(\x)\G^{-1}
&\preceq(1+c\theta)\G^{-1}.
\end{align}
On the other hand,
Proposition \ref{prop:gram_lip} and Lemma \ref{lm:GNlinearquadra} give
\begin{align}
\label{eq:Glowerbound}
    \mu^2\I_d
&\preceq\H(\x)\preceq\G.
\end{align}
Combining above two results, we obtain
\begin{align}
    \label{eq:boundGinversephiGinverse}
\Norm{\bigl(\tilde{\nabla}^2\phi(\x)\bigr)^{1/2}\G^{-1}}^2
&=\Norm{\G^{-1}\tilde{\nabla}^2\phi(\x)\G^{-1}}\overset{\eqref{eq:boundGinverse2}}{\leq}(1+c\theta)\Norm{\G^{-1}}
\overset{\eqref{eq:Glowerbound}}{\leq}\frac{1+c\theta}{\mu^2}.
\end{align}
Thus, the definition of $\mR$ in equation \eqref{eq:dfn-R} implies
\begin{align}
    \label{eq:bound_w}
\begin{split}
\|\w\|
&\overset{\eqref{dfn:proof-w}}\leq\Norm{\bigl(\tilde{\nabla}^2\phi(\x)\bigr)^{1/2}\G^{-1}}\Norm{\R}\\
&\!\!\!\!\!\!\overset{\eqref{eq:boundGinversephiGinverse},\,\eqref{eq:x_star_philip}}{\leq}
\frac{\sqrt{1+c\theta}\,\rho}{\mu}\|\x-\x_*\|^2
\overset{\eqref{eq:distance_lambda}}{\leq}
\frac{\sqrt{1+c\theta}\,\rho}{(1-c\theta)\mu^3}
\bigl(\lambda_{\mathrm{over}}(\x)\bigr)^2.
\end{split}
\end{align}
\end{itemize}
We then define
\begin{align*}
K
&\triangleq\frac{2(1+c\theta)LL_2}{(1-c\theta)^{3/2}\mu^3}
+\frac{\sqrt{1+c\theta}\,\rho}{(1-c\theta)\mu^3}.
\end{align*}
Plugging the above upper bounds for $\u$, $\vv$ and $\vw$ into~equation \eqref{eq:splithatlambda} gives
\begin{align}
\label{eq:phix_+-x}
\Norm{\big(\tilde{\nabla}^2\phi(\x)\big)^{1/2}(\x_{+}-\x_*)}
&\overset{\eqref{eq:bound_u},\,\eqref{eq:bound_v},\,\eqref{eq:bound_w}}{\leq}\left(1-\frac{1-c\theta}{\eta}\right)\lambda_{\mathrm{over}}(\x)
+K\bigl(\lambda_{\mathrm{over}}(\x)\bigr)^2.
\end{align}
On the other hand, Lemma \ref{lm:pd_approximate_phi} gives
\begin{align}
\label{eq:lower_x_+-x^*over}
\!\!\!\big\|\big(\tilde{\nabla}^2\phi(\x)\big)^{1/2}
(\x_{+}-\x_*)\big\|^2
&=(\x_{+}-\x_*)^{\top}\tilde{\nabla}^2\phi(\x)(\x_{+}-\x_*)\overset{\eqref{eq:approx_phi_positive}}{\geq}
(1-c\theta)\mu^2\|\x_{+}-\x_*\|^2.
\end{align}
Recall that definitions of $K$ and $C_1$ give $K\leq C_1$.
Hence, we achieve equation \eqref{eq:overfit_distance_lambda} as follows
\begin{align*}
\|\x_{+}-\x_*\|
&\overset{\eqref{eq:lower_x_+-x^*over}}{\leq}\frac{1}{\sqrt{1-c\theta}\mu}
\Norm{\bigl(\tilde{\nabla}^2\phi(\x)\bigr)^{1/2}
(\x_{+}-\x_*)}\\
&\overset{\eqref{eq:phix_+-x}}{\leq}
\frac{1}{\sqrt{1-c\theta}\mu}
\left(\left(1-\frac{1-c\theta}{\eta}\right)\lambda_{\mathrm{over}}(\x)
+K(\lambda_{\mathrm{over}}(\x))^2\right)\\
&\,~\leq\,~\frac{1}{\sqrt{1-c\theta}\mu}
\left(\left(1-\frac{1-c\theta}{\eta}\right)\lambda_{\mathrm{over}}(\x)
+C_1(\lambda_{\mathrm{over}}(\x))^2\right).
\end{align*}

We then prove equation \eqref{eq:movedistance_lambda} as follows
\begin{align}\label{eq:x_+-xlambda}
\begin{split}    
\|\x_{+}-\x\|
&\,\,\overset{\eqref{eq:qNupdate}}{=}~\, 
\Norm{\G^{-1}\J(\x)^{\top}\F(\x)} \\
&\,~=\,~\Norm{\G^{-1}\bigl(\J(\x)^{\top}\F(\x)
-\J(\x_*)^{\top}\F(\x_*)\bigr)} \\
& \overset{\eqref{eq:Glowerbound}}\leq\frac{\Norm{\J(\x)^{\top}(\F(\x)-\F(\x_*))}
 +\Norm{(\J(\x)^{\top}-\J(\x_*)^{\top})\F(\x_*)}}{\mu^2} \\
&\,\,\overset{\eqref{eq:lip}}{\leq}\,\,
\frac{L^2+L_2\|\F(\x_*)\|}{\mu^2}\|\x-\x_*\| \\
&\overset{\eqref{eq:distance_lambda}}{\leq}
\frac{L^2+L_2\|\F(\x_*)\|}{\sqrt{1-c\theta}\mu^3}
\lambda_{\mathrm{over}}(\x).
\end{split}
\end{align}

Recall that $\widehat M=2LL_2/((1-c\theta)\mu^2)$ and
$D=LL_2(L^2+L_2\|\F(\x_*)\|)$, which implies
\begin{align}
    \label{eq:Mx+-xbound}
    \begin{split}
\widehat M\|\x_{+}-\x\|
&\overset{\eqref{eq:x_+-xlambda}}{\leq}
\frac{2D\lambda_{\mathrm{over}}(\x)}{(1-c\theta)^{3/2}\mu^5} \\
&\,\overset{\eqref{eq:condi_x_linear_quadra}}{\leq}\,
\frac{2D}{(1-c\theta)^{3/2}\mu^5}
\cdot\frac{(1-c\theta)^{3/2}\mu_*^5}{16D}\\
&~\,=~\,\frac{\mu_*^5}{8\mu^5}
=\frac{1}{\sqrt{2}}<1,
\end{split}
\end{align}
where the last equality uses the fact $\mu=\mu_*/\sqrt{2}$.
Applying Lemma \ref{lm:pd_approximate_phi} with $\y=\x_{+}$ gives
\begin{align*}
\tilde{\nabla}^2\phi(\x_{+})
&\overset{\eqref{eq:concor_phi}}{\succeq}
\bigl(1-\widehat M\|\x_{+}-\x\|\bigr)
\tilde{\nabla}^2\phi(\x)\\
&\overset{\eqref{eq:approx_phi_positive}}{\succeq}
\bigl(1-\widehat M\|\x_{+}-\x\|\bigr)
    (1-c\theta)\mu^2\I_d\\
&\!\overset{\eqref{eq:Mx+-xbound}}{\succeq}\biggl(1-\frac{1}{\sqrt{2}}\biggr)
(1-c\theta)\mu^2\I_d\succ\0.
\end{align*}
Thus, the measure $\lambda_{\mathrm{over}}(\cdot)$ is well-defined, and we have
\begin{align*}
\bigl(\lambda_{\mathrm{over}}(\x_{+})\bigr)^2
&=(\x_{+}-\x_*)^{\top}\tilde{\nabla}^2\phi(\x_{+})(\x_{+}-\x_*)\\
&\!\!\overset{\eqref{eq:concor_phi}}{\leq}\!\!
\bigl(1+\widehat M\|\x_{+}-\x\|\bigr)
(\x_{+}-\x_*)^{\top}\tilde{\nabla}^2\phi(\x)(\x_{+}-\x_*)\\
&=\bigl(1+\widehat M\|\x_{+}-\x\|\bigr)
\Norm{\bigl(\tilde{\nabla}^2\phi(\x)\bigr)^{1/2}
(\x_{+}-\x_*)}^2.
\end{align*}
Note that $\lambda_{\mathrm{over}}(\x_{+})\geq0$ and
$1+\widehat M\|\x_{+}-\x\|\geq1$. Therefore, taking square roots on above inequality proves equation \eqref{eq:overfit_linear_quadra} as follows
\begin{align*}
\lambda_{\mathrm{over}}(\x_{+})
&\,~\leq~\,\sqrt{1+\widehat M\|\x_{+}-\x\|}
\Norm{\bigl(\tilde{\nabla}^2\phi(\x)\bigr)^{1/2}
(\x_{+}-\x_*)}\\
&\overset{\eqref{eq:Mx+-xbound}}{\leq}
\sqrt{1+\frac{2D\lambda_{\mathrm{over}}(\x)}{(1-c\theta)^{3/2}\mu^5}
}
\Norm{\bigl(\tilde{\nabla}^2\phi(\x)\bigr)^{1/2}
(\x_{+}-\x_*)}\\
&\overset{\eqref{eq:phix_+-x}}{\leq}
\sqrt{1+\frac{2D\lambda_{\mathrm{over}}(\x)}{(1-c\theta)^{3/2}\mu^5}}\,\left(\left(1-\frac{1-c\theta}{\eta}\right)\lambda_{\mathrm{over}}(\x)
+ K\bigl(\lambda_{\mathrm{over}}(\x)\bigr)^2\right)\\
&~\leq\,~\left(1+\frac{2D\lambda_{\mathrm{over}}(\x)}{(1-c\theta)^{3/2}\mu^5}\right)
\left(\left(1-\frac{1-c\theta}{\eta}\right)\lambda_{\mathrm{over}}(\x)
+K\bigl(\lambda_{\mathrm{over}}(\x)\bigr)^2\right)\\
&~=\left(1-\frac{1-c\theta}{\eta}\right)\lambda_{\mathrm{over}}(\x) + \left(\frac{2D}{(1-c\theta)^{3/2}\mu^5}\left(1-\frac{1-c\theta}{\eta}\right)+K\right)
\bigl(\lambda_{\mathrm{over}}(\x)\bigr)^2 \\
&~~~~~~~~~~+ \frac{2DK(\lambda_{\mathrm{over}}(\vx))^{3}}{(1-c\theta)^{3/2}\mu^5}
\\
&~\leq\,~\left(1-\frac{1-c\theta}{\eta}\right)\lambda_{\mathrm{over}}(\x)
+\left(\frac{2D}{(1-c\theta)^{3/2}\mu^5}+2K\right)
\bigl(\lambda_{\mathrm{over}}(\x)\bigr)^2\\
&~\leq\,~\left(1-\frac{1-c\theta}{\eta}\right)\lambda_{\mathrm{over}}(\x)
+C_1\bigl(\lambda_{\mathrm{over}}(\x)\bigr)^2,
\end{align*}
where the second last inequality is based on the result $2D\lambda_{\mathrm{over}}(\vx)/((1-c\theta)^{3/2}\mu^5)\leq 1$ achieved by the condition on $\lambda_{\mathrm{over}}(\x)$ shown in \eqref{eq:condi_x_linear_quadra}
and the last inequality is based on definitions of~$C_1$ and $K$ which imply $2D/((1-c\theta)^{3/2}\mu^5)+2K\leq C_1$.
\end{proof}

\subsection{Proof of Theorem~\ref{thm:linear1n>d}}
\label{app:proof_theorem_linear_over}
\begin{proof}
For the ease of presentation, we denote
\begin{align*}
\lambda_{\mathrm{over},t}\triangleq\lambda_{\mathrm{over}}(\x_t)
\end{align*}
in the following analysis.

% Recall that $D=L^3L_2+LL_2^2\|\F(\x_*)\|$, then the definition of $C_1$ gives
% $C_1\geq2D/((1-c\theta)^{3/2}\mu^5)$ and
% $C_1\geq2cLL_2/((c-1)\mu^3)$.

We now use induction to prove equation \eqref{eq:induc_lambda} and the result
\begin{align}
\label{eq:inducx_tdistance}
\|\x_t-\x_*\|
\leq\frac{(c-1)\mu_*^2}{4cLL_2}
\end{align}
hold for all $t\geq0$.
For the induction base, 
the definition of $\eta_0$ gives
\begin{align*}
\H_0\preceq\G_0\preceq\eta_0\H_0,
\end{align*}
which implies equation \eqref{eq:induc_lambda} holds for $t=0$.
Additionally, the initial condition \eqref{eq:linear_initial_condi} directly gives~equation \eqref{eq:inducx_tdistance} for $t=0$.
For the induction step, we suppose equations 
\eqref{eq:inducx_tdistance} and \eqref{eq:induc_lambda} hold for all $i=0,1,\dots,t$.
Then equation \eqref{eq:inducx_tdistance} implies that $\x_i\in\Omega_{\mathrm{over}}$ for all $i=0,1,\dots,t$, which gives
\begin{align}\label{eq:lambda_over_no_grow}
\lambda_{\mathrm{over},t}
\overset{\eqref{eq:induc_lambda}}{\leq}\lambda_{\mathrm{over},0}
\overset{\eqref{eq:linear_initial_condi}}{\leq}
\frac{1-c\theta}{10\eta_0 C_1}
\overset{\eqref{eq:mu-C1}}{\leq} 
\frac{(1-c\theta)^{3/2}\mu_*^5}{16D}. 
\end{align}
Applying Lemma~\ref{lm:GNlinearquadra} with $\vx=\x_t$  and $\vx_+=\vx_{t+1}$ gives
\begin{align}\label{eq:lambda_over_t+1_bound}
\begin{split}
\lambda_{\mathrm{over},t+1}
&~~\overset{\eqref{eq:overfit_linear_quadra}}{\leq}~~
\biggl(1-\frac{1-c\theta}{\eta_t}
+C_1\lambda_{\mathrm{over},t}\biggr)\lambda_{\mathrm{over},t}\\
&\!\!\overset{\eqref{eq:induc_lambda},\,\eqref{eq:lambda_over_no_grow}}{\leq}\!\!\biggl(1-\frac{2(1-c\theta)}{3\eta_0}
+\frac{1-c\theta}{10\eta_0}\biggr)\lambda_{\mathrm{over},t}\\
&~~~\leq~~~\biggl(1-\frac{1-c\theta}{2\eta_0}\biggr)
\lambda_{\mathrm{over},t}\\
&~~\,\overset{\eqref{eq:induc_lambda}}{\leq}~~\biggl(1-\frac{1-c\theta}{2\eta_0}\biggr)^{t+1}
\lambda_{\mathrm{over},0}.
\end{split}
\end{align}
It also implies
\begin{align*}
\|\x_{t+1}-\x_*\|
&\,\,\overset{\eqref{eq:overfit_distance_lambda}}{\leq}\,
\frac{1}{\sqrt{1-c\theta}\mu}
\left(\left(1-\frac{1-c\theta}{\eta_t}\right)
\lambda_{\mathrm{over},t}+C_1\lambda_{\mathrm{over},t}^2\right)\\
&\overset{\eqref{eq:lambda_over_t+1_bound}}{\leq}
\frac{1}{\sqrt{1-c\theta}\mu}
\left(1-\frac{1-c\theta}{2\eta_0}\right)^{t+1}
\lambda_{\mathrm{over},0}\\
&\,~\leq~\frac{\lambda_{\mathrm{over},0}}
{\sqrt{1-c\theta}\mu}\overset{\eqref{eq:linear_initial_condi}}{\leq}\frac{\sqrt{1-c\theta}}{10\eta_0 C_1 \mu}\\
&\,\,\overset{\eqref{eq:mu-C1}}{\leq}
\frac{\sqrt{1-c\theta}\,(c-1)\mu^2}{20c\eta_0LL_2}
\leq\frac{(c-1)\mu_*^2}{4cLL_2},
\end{align*}
which means
\begin{align*}
\x_{t+1}\in\Omega_{\mathrm{over}}.    
\end{align*}
Furthermore, 
the induction hypothesis on equation \eqref{eq:induc_lambda} and
Proposition~\ref{prop:Gram_concor} with $\vx=\vx_t$ and~$\vy=\vx_{t+1}$ imply
\begin{align*}
\widetilde{\G}_t
&=(1+Mr_t)\G_t\overset{\eqref{eq:induc_lambda}}{\preceq}
\eta_0(1+Mr_t)
\exp\biggl(2\sum_{i=0}^{t-1}Mr_i\biggr)\H_t\\
&\overset{\eqref{eq:Gram_bound_0}}{\preceq}
\eta_0(1+Mr_t)^2
\exp\biggl(2\sum_{i=0}^{t-1}Mr_i\biggr)\H_{t+1}\preceq\eta_0\exp\biggl(2\sum_{i=0}^{t}Mr_i\biggr)\H_{t+1},
\end{align*}
where the last step is based on the fact $1+x\leq {\rm e}^{x}$ with $x=Mr_{\hat{t}}$.
It also implies
\begin{align*}
\H_{t+1}
&\overset{\eqref{eq:Gram_bound_0}}{\preceq}
(1+Mr_t)\H_t
\overset{\eqref{eq:induc_lambda}}{\preceq}(1+Mr_t)\G_t
=\widetilde{\G}_t.
\end{align*}
Therefore, we can Lemma~\ref{lm:QN_order} with $\mG=\widetilde{\G}_t$ and $\mH=\mH_{t+1}$ to achive
\begin{align}
\label{eq:over_G_t+1_H_t+1}
\H_{t+1}
&\preceq\G_{t+1}
\preceq\eta_0\exp\biggl(2\sum_{i=0}^{t}Mr_i\biggr)\H_{t+1}.
\end{align}
% Using $M=2LL_2/\mu^2$,
% $D=LL_2(L^2+L_2\|\F(\x_*)\|)$, and~\eqref{eq:movedistance_lambda},
% we have
% Applying the third inequality displayed at the beginning of the proof for every
% $0\leq i\leq t$, the induction hypothesis and the initial condition give
Additionally, Lemma \ref{lm:GNlinearquadra} with $\vx=\vx_t$ and $\vx_+=\vx_{t+1}$ implies
\begin{align}
M\sum_{i=0}^{t}r_i
&\overset{\eqref{eq:movedistance_lambda}}{\leq}
\frac{2D}{\sqrt{1-c\theta}\mu^5}
\sum_{i=0}^{t}\lambda_{\mathrm{over},i}\notag\\
&\overset{\eqref{eq:induc_lambda}}{\leq}
\frac{2D\lambda_{\mathrm{over},0}}{\sqrt{1-c\theta}\mu^5}
\sum_{i=0}^{t}\biggl(1-\frac{1-c\theta}{2\eta_0}\biggr)^i\notag\\
&~\leq\frac{4D\lambda_{\mathrm{over},0}\eta_0}{(1-c\theta)^{3/2}\mu^5}
\overset{\eqref{eq:linear_initial_condi}}{\leq}
\frac{2D}{5\sqrt{1-c\theta}\mu^5C_1}
\leq\frac{1-c\theta}{5}
\leq\frac{1}{5}.
\label{eq:sumr_ibound}
\end{align}
Consequently, we have
\begin{align*}
\G_{t+1}
\overset{\eqref{eq:over_G_t+1_H_t+1}}{\preceq}\eta_0\exp\biggl(2M\sum_{i=0}^{t}r_i\biggr)\H_{t+1}
\overset{\eqref{eq:sumr_ibound}}{\preceq}\eta_0{\rm e}^{2/5}\H_{t+1}
\preceq\frac{3\eta_0}{2}\H_{t+1}.
\end{align*}
This completes the induction. 
\end{proof}

\subsection{Proof of Theorem~\ref{thm:two-stage-n>d}}
\label{app:proof_theorem_two_stage_over}
\begin{proof}
We follow the notation  $\lambda_{\mathrm{over},t}=\lambda_{\mathrm{over}}(\x_t)$ defined in Appendix \ref{app:proof_theorem_linear_over}.
Condition~\eqref{eq:inicondi_final_n>d} means
the initial condition of Theorem~\ref{thm:linear1n>d} holds, then we achieve
\begin{align}
\label{eq:linear_path_over}
\lambda_{\mathrm{over},t}\overset{\eqref{eq:induc_lambda}}\leq \biggl(1-\frac{1-c\theta}{2\eta_0}\biggr)^t
\lambda_{\mathrm{over},0}
\qquad \text{and} \qquad
\|\x_t-\x_*\|\overset{\eqref{eq:inicondi_final_n>d}}{\leq}\frac{(c-1)\mu_*^2}{4cLL_2}.
\end{align}
Additionally, Lemma~\ref{lm:Gram_approximate} with $\Omega=\Omega_{\mathrm{over}}$ and $m=d$ gives
\begin{align}\label{eq:thm_delta_t_recur}
\H_t\preceq\G_t\preceq(1+\delta_t)\H_t
\qquad\text{and}\qquad
\mathbb E_t[\delta_{t+1}]
\leq(1-\gamma)(1+Mr_t)^2(\delta_t+2\tilde c dMr_t).
\end{align}
% As verified in the induction proof of Theorem~\ref{thm:linear1n>d},
% condition~\eqref{eq:linear_initial_condi} and~\eqref{eq:linear_path_over}
% satisfy the distance and $\lambda_{\mathrm{over}}$ conditions in~\eqref{eq:condi_x_linear_quadra}.
% Together with the matrix bounds above,
Note that equations \eqref{eq:linear_path_over} and \eqref{eq:thm_delta_t_recur} indicate conditions of
Lemma~\ref{lm:GNlinearquadra} holds by taking $\vx=\vx_t$, $\vx_+=\vx_{t+1}$, $\mG=\mG_t$, and $\eta=1+\delta_t$, then we achieve
\begin{align}
\label{eq:Mr_tboundedbylambda_t}
Mr_t
&\overset{\eqref{eq:movedistance_lambda}}{\leq}
\frac{2D\lambda_{\mathrm{over},t}}{\sqrt{1-c\theta}\mu^5}.
\end{align}
and
\begin{align}\label{eq:lambda_over_t+1boundbylambda_t}
\begin{split}    
\lambda_{\mathrm{over},t+1}
&\overset{\eqref{eq:overfit_linear_quadra}}{\leq}
\left(1-\frac{1-c\theta}{1+\delta_t}
+C_1\lambda_{\mathrm{over},t}\right)\lambda_{\mathrm{over},t} \\
&~=~ \left(\frac{\delta_t+c\theta}{1+\delta_t}
+C_1\lambda_{\mathrm{over},t}\right)\lambda_{\mathrm{over},t} \\
&~\leq\,\,\bigl(\delta_t
+C_1\lambda_{\mathrm{over},t}+c\theta\bigr)\lambda_{\mathrm{over},t}.
\end{split}
\end{align}
Therefore, the definition $C_2\triangleq C_1+4\tilde{c}dD/\big(\sqrt{1-c\theta}\mu^5\big)$ leads to
\begin{align}\label{eq:E_tdeltat+1}
\begin{split}    
\mathbb E_t[\delta_{t+1}]
&\overset{\eqref{eq:thm_delta_t_recur}}{\leq}
(1-\gamma)(1+Mr_t)^2(\delta_t+2\tilde c dMr_t)\\
&\overset{\eqref{eq:Mr_tboundedbylambda_t}}{\leq}
(1-\gamma)\biggl(1+\frac{2D\lambda_{\mathrm{over},t}}{\sqrt{1-c\theta}\mu^5}
\biggr)^2
\biggl(\delta_t+\frac{4\tilde c dD\lambda_{\mathrm{over},t}}{\sqrt{1-c\theta}\mu^5}
\biggr)\\
&~\,=~\,(1-\gamma)\biggl(1+\frac{2D\lambda_{\mathrm{over},t}}{\sqrt{1-c\theta}\mu^5}
\biggr)^2
\bigl(\delta_t+(C_2-C_1)\lambda_{\mathrm{over},t}\bigr).
\end{split}
\end{align}
We denote
\begin{align}\label{eq:xi-t}
    \xi_t\triangleq\delta_t+C_2\lambda_{\mathrm{over},t} \geq \delta_t+C_1\lambda_{\mathrm{over},t},
\end{align}
then it holds
\begin{align}\label{eq:overfit_delta_recur}
\begin{split}    
\mathbb E_t[\delta_{t+1}]
&\overset{\eqref{eq:E_tdeltat+1}}{\leq}(1-\gamma)\left(1+\frac{2D\lambda_{\mathrm{over},t}}{\sqrt{1-c\theta}\mu^5}
\right)^2\bigl(\delta_t+(C_2-C_1)\lambda_{\mathrm{over},t}\bigr)\\
&~\,\leq\,~(1-\gamma)
\left(1+\frac{C_2\lambda_{\mathrm{over},t}}{2}\right)^2
\left(\delta_t+C_2\lambda_{\mathrm{over},t}\right)\\
&\overset{\eqref{eq:xi-t}}{=}(1-\gamma)
\left(1+\frac{C_2\lambda_{\mathrm{over},t}}{2}\right)^2\xi_t,  
\end{split}
\end{align}
where the second inequality is based on the definition of $C_2$ and the setting $\tilde c \geq 1$ which result~$2D/(\sqrt{1-c\theta}\mu^5)\leq C_2/2$. 
Additionally, equation~\eqref{eq:lambda_over_t+1boundbylambda_t} gives
\begin{align}
\label{eq:overfit_lambda_recur}
\begin{split}
C_2\lambda_{\mathrm{over},t+1}
 \overset{\eqref{eq:lambda_over_t+1boundbylambda_t}}{\leq} C_2\lambda_{\mathrm{over},t}
\bigl(\delta_t+C_1\lambda_{\mathrm{over},t}\bigr)
+c\theta C_2\lambda_{\mathrm{over},t}  
 \overset{\eqref{eq:xi-t}}\leq C_2\lambda_{\mathrm{over},t}\xi_t
+c\theta C_2\lambda_{\mathrm{over},t}.
\end{split}
\end{align}
Recall that we have defined $\nu\triangleq\min\{\gamma,(1-c\theta)/(2\eta_0)\}<\gamma<1$, then combining results of~\eqref{eq:overfit_delta_recur} and \eqref{eq:overfit_lambda_recur} gives
\begin{subequations}
\begin{align}
\label{eq:xi_one_step_over_raw}
\mathbb E_t[\xi_{t+1}]
&\!\overset{\eqref{eq:overfit_delta_recur},\,\eqref{eq:overfit_lambda_recur}}{\leq}\Bigl((1-\gamma)
\Bigl(1+\frac{C_2\lambda_{\mathrm{over},t}}{2}\Bigr)^2
+C_2\lambda_{\mathrm{over},t}\Bigr)\xi_t
+c\theta C_2\lambda_{\mathrm{over},t}\\
\label{eq:xi_one_step_over}
&~~~~\leq (1-\nu)
\Bigl(\Bigl(1+\frac{C_2\lambda_{\mathrm{over},t}}{2}\Bigr)^2
+\frac{C_2\lambda_{\mathrm{over},t}}{1-\nu}\Bigr)\xi_t
+c\theta C_2\lambda_{\mathrm{over},t}\notag\\
&~~~~\leq (1-\nu)
\Bigl(1+\frac{C_2\lambda_{\mathrm{over},t}}{2}\Bigr)^2\Bigl(1
 + \frac{C_2\lambda_{\mathrm{over},t}}{1-\nu}\Bigr)\xi_t
 + c\theta C_2\lambda_{\mathrm{over},t}\notag\\
&~~~~\leq(1-\nu)\exp\Bigl(C_2\lambda_{\mathrm{over},t}\Bigr)
\exp\Bigl(\frac{C_2\lambda_{\mathrm{over},t}}{1-\nu}\Bigr)\xi_t
+c\theta C_2\lambda_{\mathrm{over},t}\notag\\
&~~~~=(1-\nu)
\exp\Bigl(\frac{(2-\nu)C_2\lambda_{\mathrm{over},t}}{1-\nu}
\Bigr)\xi_t
+c\theta C_2\lambda_{\mathrm{over},t},
\end{align}
\end{subequations}
where the last second inequality is based on the result $1+s\leq {\rm e}^s$ with~$s=C_2\lambda_{{\rm over},t}$
and~$s={C_2\lambda_{\mathrm{over},t}}/(1-\nu)$.


Taking expectation on~result \eqref{eq:xi_one_step_over} and applying~equation \eqref{eq:linear_path_over} give
\begin{align}
    \label{eq:xi_t+1_recur_xi_t}
    \begin{split}
\mathbb E[\xi_{t+1}]
&\overset{\eqref{eq:xi_one_step_over},\,\eqref{eq:linear_path_over}}{\leq}
(1-\nu)
\exp\Biggl(\frac{(2-\nu)C_2}{1-\nu}
\biggl(1-\frac{1-c\theta}{2\eta_0}\biggr)^t
\lambda_{\mathrm{over},0}\Biggr)\mathbb E[\xi_t]
\\
&~~~~~~~~~~~~~~~~~+c\theta C_2\biggl(1-\frac{1-c\theta}{2\eta_0}\biggr)^t
\lambda_{\mathrm{over},0}.
    \end{split}
\end{align}
Expanding the term $\mathbb E[\xi_t]$ according to above recursion gives
\begin{align*}
\mathbb E[\xi_{t+1}]
&\leq(1-\nu)^2
\exp\Biggl(\frac{(2-\nu)C_2}{1-\nu}\lambda_{\mathrm{over},0}
\Bigl(\bigl(1-\frac{1-c\theta}{2\eta_0}\bigr)^t
+\bigl(1-\frac{1-c\theta}{2\eta_0}\bigr)^{t-1}\Bigr)\Biggr)
\mathbb E[\xi_{t-1}]\\
&\quad+c\theta C_2(1-\nu)
\biggl(1-\frac{1-c\theta}{2\eta_0}\biggr)^{t-1}
\lambda_{\mathrm{over},0}
\exp\Biggl(\frac{(2-\nu)C_2}{1-\nu}
\biggl(1-\frac{1-c\theta}{2\eta_0}\biggr)^t
\lambda_{\mathrm{over},0}\Biggr)\\
&\quad+c\theta C_2\biggl(1-\frac{1-c\theta}{2\eta_0}\biggr)^t
\lambda_{\mathrm{over},0}.
\end{align*}
Repeating this substitution gives
\begingroup
\small
\begin{align}\label{eq:finialxi_t_recur}
\begin{split}
&\!\!\!\mathbb E[\xi_t]
\leq(1-\nu)^t\xi_0
\exp\Biggl(\frac{(2-\nu)C_2\lambda_{\mathrm{over},0}}{1-\nu}
\sum_{j=0}^{t-1}\biggl(1-\frac{1-c\theta}{2\eta_0}\biggr)^j\Biggr)\\
&~~+c\theta C_2\lambda_{\mathrm{over},0}
\sum_{i=0}^{t-1}(1-\nu)^{t-1-i}
\biggl(1-\frac{1-c\theta}{2\eta_0}\biggr)^i
\exp\Biggl(\frac{(2-\nu)C_2\lambda_{\mathrm{over},0}}{1-\nu}
\sum_{j=i+1}^{t-1}\biggl(1-\frac{1-c\theta}{2\eta_0}\biggr)^j\Biggr)
    \end{split}
\end{align}
\endgroup
for all $t\geq0$.
In addition, we have
\begin{align*}
\frac{(2-\nu)C_2}{1-\nu}
\lambda_{\mathrm{over},0}
\sum_{j=i}^{t-1}\biggl(1-\frac{1-c\theta}{2\eta_0}\biggr)^j
\leq\frac{2(2-\nu)C_2\eta_0}{(1-\nu)(1-c\theta)}
\lambda_{\mathrm{over},0}\overset{\eqref{eq:inicondi_final_n>d}}{\leq}
\frac{(2-\nu)(1-\gamma)}{5(1-\nu)(2-\gamma)}
\leq\frac{1}{5}
\end{align*}
for all  $i=0,\dots,t$, 
where the last inequality uses the setting $\nu\leq\gamma$.
Consequently, we have
\begin{align}
\label{eq:exponential_bound}
\exp\Biggl(\frac{(2-\nu)C_2}{1-\nu}
\lambda_{\mathrm{over},0}
\sum_{j=i}^{t-1}\biggl(1-\frac{1-c\theta}{2\eta_0}\biggr)^j\Biggr)
&\leq {\rm e}^{1/5}.
\end{align}
Recall that $\eta_0\geq1$ and $\nu\leq(1-c\theta)/(2\eta_0)$, then we have
\begin{align*}
c\theta=1-(1-c\theta)\leq1 - \frac{1-c\theta}{2\eta_0} \leq 1-\nu,
\end{align*}
which implies
\begin{align}
\label{eq:sumweightedbound}
\sum_{i=0}^{t-1}c\theta(1-\nu)^{t-1-i}
\biggl(1-\frac{1-c\theta}{2\eta_0}\biggr)^i
&\leq\sum_{i=0}^{t-1}(1-\nu)^{t-i}(1-\nu)^i
=t(1-\nu)^t.
\end{align}
Combing above reulst, we achive
\begin{align}
\mathbb E[\xi_t]
&\overset{\eqref{eq:finialxi_t_recur},\, \eqref{eq:exponential_bound}}{\leq} {\rm e}^{1/5}(1-\nu)^t\xi_0
 +{\rm e}^{1/5}c\theta C_2\lambda_{\mathrm{over},0}
 \sum_{i=0}^{t-1}(1-\nu)^{t-1-i}
 \biggl(1-\frac{1-c\theta}{2\eta_0}\biggr)^i\notag\\
&~~~~\,=\,\,{\rm e}^{1/5}\Biggl((\delta_0+C_2\lambda_{\mathrm{over},0})(1-\nu)^t
+C_2\lambda_{\mathrm{over},0}\sum_{i=0}^{t-1}c\theta(1-\nu)^{t-1-i}
\biggl(1-\frac{1-c\theta}{2\eta_0}\biggr)^i\Biggr)\notag\\
&\,\,\,~~\overset{\eqref{eq:sumweightedbound}}{\leq} {\rm e}^{1/5}
\bigl(\delta_0+(t+1)C_2\lambda_{\mathrm{over},0}\bigr)(1-\nu)^t\notag\\
&~~~~\overset{\eqref{eq:inicondi_final_n>d}}{\leq}
{\rm e}^{1/5}\left(\delta_0+
\frac{(1-\gamma)(1-c\theta)(t+1)}{10(2-\gamma)\eta_0}\right)
(1-\nu)^t\notag\\
\label{eq:xi_decay_over}
&~~\,\,\,\,\,\leq\,\,q_{\text{\rm over}}(t+1)(1-\nu)^t.
\end{align}
The relation of $\lambda_{\mathrm{over},t+1}$ and $\lambda_{\mathrm{over},t}$ in~\eqref{eq:overfit_lambda_recur} gives
\begin{align}
\label{eq:ratio_one_step_over}
\frac{\lambda_{\mathrm{over},t+1}}
{\lambda_{\mathrm{over},t}}
\leq c\theta+\xi_t.
\end{align}
Note that $0\leq\eta_t-1\leq\delta_t\leq\xi_t$, then taking expectation and using~\eqref{eq:xi_decay_over} give
\begin{align*}
\EBP{\eta_t-1}
\leq\EBP{\xi_t}
\overset{\eqref{eq:xi_decay_over}}{\leq}
q_{\text{\rm over}}(t+1)(1-\nu)^t\\
\end{align*}
and
\begin{align*}
\EBP{\frac{\lambda_{\mathrm{over},t+1}}{\lambda_{\mathrm{over},t}}}
\overset{\eqref{eq:ratio_one_step_over}}{\leq}
c\theta+\EBP{\xi_t}
\overset{\eqref{eq:xi_decay_over}}{\leq}
c\theta+q_{\text{\rm over}}(t+1)(1-\nu)^t,
\end{align*}
which proves~equation \eqref{eq:eta-ratio-over}.

It remains to prove the high-probability result.
We define the event $\fE_j$ as the inequality
\begin{align*}
    \xi_j>\frac{4q_{\text{\rm over}}}{p\nu^2}
\left(1-\frac{\nu}{2}\right)^j
\end{align*}
holds. Then we have
\begin{align*}
\mathbb P\Bigg(\bigcup_{j=0}^{+\infty}\fE_j\Bigg)
& \leq \sum_{j=0}^{\infty}\mathbb P\left(\fE_j\right) = \sum_{j=0}^{+\infty}\mathbb P\left(\xi_j>\frac{4q_{\text{\rm over}}}{p\nu^2}
\left(1-\frac{\nu}{2}\right)^j\right) \\
& \leq
\frac{p\nu^2}{4q_{\text{\rm over}}}
\sum_{j=0}^{+\infty}\frac{\mathbb E[\xi_j]}{(1-\nu/2)^j}\overset{\eqref{eq:xi_decay_over}}{\leq}
\frac{p\nu^2}{4}\sum_{j=0}^{\infty}(j+1)
\left(\frac{1-\nu}{1-\nu/2}\right)^j\\
& \leq\frac{p\nu^2}{4}\sum_{j=0}^{+\infty}(j+1)
\left(1-\frac{\nu}{2}\right)^j=p,
\end{align*}
where the first inequality is based on union bound;
the second inequality is based on Markov's inequality; 
the last inequality follows $1-\nu\leq(1-\nu/2)^2$;
and the equality uses the result~$\sum_{j=0}^{\infty}(j+1)z^j=1/(1-z)^2$ with $z=1-\nu/2$.
This implies
\begin{align*}
\mathbb P\left(~\bigcap_{j=0}^{+\infty}\overline{\fE_j}~\right) 
=\mathbb P\left(~\overline{\bigcup_{j=0}^{+\infty}\fE_j}~\right) \geq 1-p.
\end{align*}
Therefore, with probability at least $1-p$, the inequality
\begin{align}
\label{eq:xi_j_union_bound}
\xi_j
&\leq\frac{4q_{\text{\rm over}}}{p\nu^2}
\left(1-\frac{\nu}{2}\right)^j
\end{align}
holds for all $j\geq0$.
We take
\begin{align*}
t_0
=\Biggl\lceil
\frac{2}{\nu}
\log\Bigl(1+\frac{8q_{\text{\rm over}}}{p(1-c\theta)\nu^2}\Bigr)
\Biggr\rceil
=\OM\Biggl(\frac{\log\bigl(1+q_{\text{\rm over}}/(p(1-c\theta)\nu^2)\bigr)}{\nu}\Biggr),
\end{align*}
then the fact $1-\nu/2\leq{\rm e}^{-\nu/2}$ and the definition of $q_{\rm over}$ leads to
\begin{align}
\label{eq:qover_phase_bound}
\frac{4q_{\text{\rm over}}}{p\nu^2}
\left(1-\frac{\nu}{2}\right)^{t_0}
\leq\frac{4q_{\text{\rm over}}}{p\nu^2}
\exp\left(-\frac{\nu t_0}{2}\right)
\leq\frac{4q_{\text{\rm over}}/(p\nu^2)}
{1+8q_{\text{\rm over}}/(p(1-c\theta)\nu^2)}
\leq\frac{1-c\theta}{2}.
\end{align}
Therefore, with probability at least $1-p$, the inequality
% \begin{align}\label{eq:xi-t0-i}
% c\theta + \xi_{t_0+i}
% &\overset{\eqref{eq:xi_j_union_bound}}{\leq}
% c\theta + \frac{4q_{\text{\rm over}}}{p\nu^2}
% \left(1-\frac{\nu}{2}\right)^{t_0+i}
% \overset{\eqref{eq:qover_phase_bound}}{\leq}
% c\theta +  \frac{1-c\theta}{2}\left(1-\frac{\nu}{2}\right)^i
% \leq c\theta +  \frac{1-c\theta}{2} = \frac{1+c\theta}{2}
% \end{align}
\begin{align}\label{eq:ratio_one_step_over_bound}
\begin{split}    
\frac{\lambda_{\mathrm{over},t_0+i+1}}
{\lambda_{\mathrm{over},t_0+i}}
& \overset{\eqref{eq:ratio_one_step_over}}{\leq}
c\theta + \xi_{t_0+i}
\overset{\eqref{eq:xi_j_union_bound}}{\leq}
c\theta + \frac{4q_{\text{\rm over}}}{p\nu^2}
\left(1-\frac{\nu}{2}\right)^{t_0+i} \\
& \overset{\eqref{eq:qover_phase_bound}}{\leq}
c\theta +  \frac{1-c\theta}{2}\left(1-\frac{\nu}{2}\right)^i
\leq c\theta +  \frac{1-c\theta}{2} = \frac{1+c\theta}{2}
\end{split}
\end{align}
holds for all $i\geq 0$.
This implies for all $t\geq1$, the upper bound
\begin{align*}
\lambda_{\mathrm{over},t_0+t}
&=\lambda_{\mathrm{over},t_0}
\prod_{i=0}^{t-1}
\frac{\lambda_{\mathrm{over},t_0+i+1}}
{\lambda_{\mathrm{over},t_0+i}}\\
&\!\!\overset{\eqref{eq:ratio_one_step_over_bound}}{\leq}\lambda_{\mathrm{over},t_0}
\prod_{i=0}^{t-1}\frac{1+c\theta}{2}\\
&=\biggl(\frac{1+c\theta}{2}\biggr)^t
\lambda_{\mathrm{over},t_0}\\
&\!\!\overset{\eqref{eq:linear_path_over}}{\leq}
\biggl(\frac{1+c\theta}{2}\biggr)^t
\biggl(1-\frac{1-c\theta}{2\eta_0}\biggr)^{t_0}
\lambda_{\mathrm{over},0}.
\end{align*}
holds with probability at least $1-p$, which proves equation \eqref{eq:n>dtwostage}.
\end{proof}

\subsection{Proof of Corollary~\ref{col:theta=0n>d}}
\label{app:proof_coro_n>d}
\begin{proof}
We follow notations 
$\lambda_{\mathrm{over},t}=\lambda_{\mathrm{over}}(\x_t)$ and
$\xi_t=\delta_t+C_2\lambda_{\mathrm{over},t}$ defined in Appendix~\ref{app:proof_theorem_two_stage_over}.
Note that equation~\eqref{eq:linear_path_over} with $\theta=0$ gives
\begin{align}
\label{eq:linear_path_theta_zero}
\lambda_{\mathrm{over},t}
\leq\biggl(1-\frac{1}{2\eta_0}\biggr)^t
\lambda_{\mathrm{over},0}
\qquad\text{and}\qquad
\|\x_t-\x_*\|\leq\frac{(c-1)\mu_*^2}{4cLL_2}.
\end{align}
Additionally, equation \eqref{eq:xi_one_step_over_raw} with $\theta=0$  gives
\begin{align}
\label{eq:xi_one_step_theta_zero}
\begin{split}    
\mathbb E_t[\xi_{t+1}]
&\leq\left((1-\gamma)
\left(1+\frac{C_2\lambda_{\mathrm{over},t}}{2}\right)^2
+C_2\lambda_{\mathrm{over},t}\right)\xi_t \\
&\leq (1-\gamma)\left(1+\frac{C_2\lambda_{\mathrm{over},t}}{2}\right)^2\left(1+\frac{C_2\lambda_{\mathrm{over},t}}{1-\gamma}\right)\xi_t\\
&\leq(1-\gamma)
\exp\left(\frac{(2-\gamma)C_2\lambda_{\mathrm{over},t}}{1-\gamma}
\right)\xi_t,
\end{split}
\end{align}
where we use inequality $1+x\leq {\rm e}^x$ with $x=C_2\lambda_{\mathrm{over},t}/2$ and $x=C_2\lambda_{\mathrm{over},t}/(1-\gamma)$ in the last step.
Setting $\theta=0$ in~\eqref{eq:overfit_lambda_recur} gives
\begin{align}
\label{eq:lambda_one_step_theta_zero}
\lambda_{\mathrm{over},t+1}
&\leq\xi_t\lambda_{\mathrm{over},t}.
\end{align}
Taking expectation on~\eqref{eq:xi_one_step_theta_zero} and using~\eqref{eq:linear_path_theta_zero}, we have
\begin{align}\label{eq:xi_decay_theta_zero}
\begin{split}    
\EBP{\xi_t}
&\!\!\overset{\eqref{eq:xi_one_step_theta_zero}}{\leq} (1-\gamma)\exp\Biggl(\frac{(2-\gamma)C_2}{1-\gamma}\lambda_{\mathrm{over},t-1}\Biggr)\EBP{\xi_{t-1}}\\
&\!\!\overset{\eqref{eq:xi_one_step_theta_zero}}{\leq} (1-\gamma)^t \exp\Biggl(\frac{(2-\gamma)C_2}{1-\gamma}\sum_{i=0}^{t-1}\lambda_{\mathrm{over},i}\Biggr)\xi_0\\
&\!\!\overset{ \eqref{eq:linear_path_theta_zero}}{\leq}
(1-\gamma)^t\exp\Biggl(\frac{(2-\gamma)C_2}{1-\gamma}
\sum_{i=0}^{t-1}\biggl(1-\frac{1}{2\eta_0}\biggr)^i
\lambda_{\mathrm{over},0}\Biggr)
(\delta_0+C_2\lambda_{\mathrm{over},0})\\
&\,\leq\,\,(1-\gamma)^t
\exp\Biggl(\frac{2(2-\gamma)C_2\eta_0\lambda_{\mathrm{over},0}}{1-\gamma}\Biggl)
(\delta_0+C_2\lambda_{\mathrm{over},0})\\
&\overset{\eqref{eq:inicondi_final_n>d}}{\leq}
(1-\gamma)^tq_{\text{\rm over}}.
\end{split}
\end{align}
Hence, we apply Lemma~\ref{lm:lin_lemma26} with
$X_t=\xi_t$, $a=q_{\text{\rm over}}$, and $\tau=1/\gamma>1$, which implies with probability at least $1-p$, the inequality
\begin{align}
\label{eq:xi_high_probability_theta_zero}
\xi_t
&\leq\biggl(1-\frac{\gamma}{1+\gamma}\biggr)^t
\frac{q_{\text{\rm over}}}{p\gamma^2}
\end{align}
holds for all $t\geq0$.
This implies
\begin{align}
\label{eq:lambda_t+1-recur_over}
\lambda_{\mathrm{over},t+1}\overset{\eqref{eq:lambda_one_step_theta_zero}}{\leq} \xi_t\lambda_{\mathrm{over},t}
\overset{
\eqref{eq:xi_high_probability_theta_zero}}{\leq}
\biggl(1-\frac{\gamma}{1+\gamma}\biggr)^t
\frac{q_{\text{\rm over}}\lambda_{\mathrm{over},t}}{p\gamma^2}.
\end{align}
Taking
\begin{align*}
t_0
=\biggl\lceil
\frac{\log\bigl(1+2q_{\text{\rm over}}/(p\gamma^2)\bigr)}{\log(1+\gamma)}
\biggr\rceil=\OM\biggl(\frac{\log\bigl(1+q_{\text{\rm over}}/(p\gamma^2)\bigr)}{\gamma}\biggr)
\end{align*}
gives 
\begin{align}
\label{eq:superlinearleq1/2_over}
\biggl(1-\frac{\gamma}{1+\gamma}\biggr)^{t_0}
\frac{q_{\text{\rm over}}}{p\gamma^2}
&\leq\frac{1}{2}.
\end{align}
Then for all $i\geq0$, it holds
\begin{align*}
\lambda_{\mathrm{over},t_0+i+1}
&\overset{\eqref{eq:lambda_t+1-recur_over}}{\leq}
\biggl(1-\frac{\gamma}{1+\gamma}\biggr)^{t_0+i}
\frac{q_{\text{\rm over}}\lambda_{\mathrm{over},t_0+i}}{p\gamma^2}
\\
&~\,=\,\,\,\biggl(1-\frac{\gamma}{1+\gamma}\biggr)^i
\biggl(\biggl(1-\frac{\gamma}{1+\gamma}\biggr)^{t_0}
\frac{q_{\text{\rm over}}}{p\gamma^2}\biggr)
\lambda_{\mathrm{over},t_0+i}\\
&\overset{\eqref{eq:superlinearleq1/2_over}}{\leq}
\frac{1}{2}\biggl(1-\frac{\gamma}{1+\gamma}\biggr)^i
\lambda_{\mathrm{over},t_0+i}
\end{align*}
with probability at least $1-p$. 
Consequently, for all $t\geq1$, the upper bound
\begin{align*}
\lambda_{\mathrm{over},t_0+t}
&=\lambda_{\mathrm{over},t_0}
\prod_{i=0}^{t-1}
\frac{\lambda_{\mathrm{over},t_0+i+1}}
{\lambda_{\mathrm{over},t_0+i}}\\
&\leq\lambda_{\mathrm{over},t_0}
\prod_{i=0}^{t-1}\left(\frac{1}{2}
\left(1-\frac{\gamma}{1+\gamma}\right)^i\right)\\
&=\left(1-\frac{\gamma}{1+\gamma}\right)^{t(t-1)/2}
\biggl(\frac{1}{2}\biggr)^t\lambda_{\mathrm{over},t_0}\\
&\!\!\!\overset{\eqref{eq:linear_path_theta_zero}}{\leq}\!\!
\biggl(1-\frac{\gamma}{1+\gamma}\biggr)^{t(t-1)/2}
\biggl(\frac{1}{2}\biggr)^t
\biggl(1-\frac{1}{2\eta_0}\biggr)^{t_0}
\lambda_{\mathrm{over},0}
\end{align*}
holds with probability at least $1-p$.
\end{proof}

\section{Implementation of RQGN with F-BFGS}
\label{app:fbfgs}

We give the detailed implementation of RQGN with the fast BFGS update (Option II) in Algorithm~\ref{alg:fbfgs}, 
which can be regarded as Algorithm~\ref{alg:RQGN} with $\G_t=(\L_t^{\top}\L_t)^{-1}$. In particular, the fast BFGS update satisfies
\begin{align*}
\G_{t+1}
&=\fbfgs(\widetilde{\G}_t,\H(\x_{t+1}),\u_t) \\
&=\bfgs(\widetilde{\G}_t,\H(\x_{t+1}),\widetilde{\L}_t^{\top}\u_t),
\end{align*}
where $\widetilde{\G}_t=\big(\widetilde{\L}_t^{\top}\widetilde{\L}_t\big)^{-1}$.
Following \citet[Proposition~5.7]{liu2023symmetric} or \citet[Equation~(9.11)]{gower2017randomized}, the update on $\mG_{t+1}$ can be implemented by the closed-form iteration 
\begin{align*}
\L_{t+1}={}&\widetilde{\L}_t+
\big(\u_t(\u_t^\top\u_t)^{-1/2}
-\widetilde{\L}_t\H(\x_{t+1})\widetilde{\L}_t^\top\u_t
(\u_t^\top\widetilde{\L}_t\H(\x_{t+1})\widetilde{\L}_t^\top\u_t)^{-1/2}\big)\\
&~\quad\quad\times
\big(\u_t^\top\widetilde{\L}_t\H(\x_{t+1})\widetilde{\L}_t^\top\u_t\big)^{-1/2}
\u_t^\top\widetilde{\L}_t,
\end{align*}
which requires the cost of $\OM(m^2)$ flops per iteration.

\begin{algorithm}[H]
\caption{RQGN with the fast BFGS update}
\label{alg:fbfgs}
\begin{algorithmic}[1]
\STATE \textbf{Input:} $\x_0\in\Omega$, $\L_0\in\RB^{m\times m}$ such that $(\L_0^{\top}\L_0)^{-1}\succeq\H(\x_0)$, and $M\geq0$ \\[0.15cm]
\STATE \textbf{for} $t=0,1,\ldots$ \\[0.15cm]
\STATE \quad $\displaystyle
\x_{t+1}=\begin{cases}
\x_t-\J(\x_t)^{\top}\L_t^{\top}\L_t\F(\x_t),& n<d,\\
\x_t-\L_t^{\top}\L_t\J(\x_t)^{\top}\F(\x_t),& n\geq d;
\end{cases}$ \\[0.15cm]
\STATE \quad $r_t=\|\x_{t+1}-\x_t\|$ and $\widetilde{\L}_t=\L_t/\sqrt{1+Mr_t}$ \\[0.15cm]
\STATE \quad $[\u_t]_i\overset{\mathrm{i.i.d.}}{\sim}\mathcal N(0,1)$ \\[0.15cm]
\STATE \quad $\displaystyle
\begin{aligned}
\L_{t+1}={}&\widetilde{\L}_t+
\big(\u_t(\u_t^\top\u_t)^{-1/2}
-\widetilde{\L}_t\H(\x_{t+1})\widetilde{\L}_t^\top\u_t
(\u_t^\top\widetilde{\L}_t\H(\x_{t+1})\widetilde{\L}_t^\top\u_t)^{-1/2}\big)\\
&~\quad\quad\times
\big(\u_t^\top\widetilde{\L}_t\H(\x_{t+1})\widetilde{\L}_t^\top\u_t\big)^{-1/2}
\u_t^\top\widetilde{\L}_t
\end{aligned}$
\STATE \textbf{end for}
\end{algorithmic}
\end{algorithm}

% \section{Additional Numerical Experiments}
% \label{app:additional_experiments}

% Existing quasi-Newton convergence results typically require the initial estimator to be close to the exact Jacobian or its inverse. We compare this behavior with RQGN for the underdetermined eigenvalue problem and the minimax problem studied in Section~\ref{sec:exp}.

% \subsection{Underdetermined Nonlinear Equations}
% \label{app:additional_under}

% We compare RQGN-SR1 with the generalized Broyden method (G-Broyden) of \citet{vater2024convergence} on the eigenvalue problem with $n=300$.

% \begin{algorithm}[H]
% \caption{Generalized Broyden method (G-Broyden)}
% \label{alg:gbroyden}
% \begin{algorithmic}[1]
% \STATE \textbf{Input:} $\x_0\in\RB^d$, $\eta>0$, and $\mathbf P_0\in\RB^{d\times n}$
% \STATE \textbf{for} $t=0,1,\ldots$
% \STATE \quad $\x_{t+1}=\x_t-\eta\mathbf P_t\F(\x_t)$
% \STATE \quad $\s_t=\x_{t+1}-\x_t$ and $\y_t=\F(\x_{t+1})-\F(\x_t)$
% \STATE \quad $\displaystyle \mathbf P_{t+1}=\mathbf P_t+
% \frac{(\s_t-\mathbf P_t\y_t)\s_t^{\top}\mathbf P_t}{\s_t^{\top}\mathbf P_t\y_t}$
% \STATE \textbf{end for}
% \end{algorithmic}
% \end{algorithm}

% We use two initializations. In Setting~1, RQGN uses $\G_0=\J(\x_0)\J(\x_0)^{\top}$ and G-Broyden uses $\mathbf P_0=\J(\x_0)^{\dag}$. In Setting~2, RQGN uses $\G_0=L^2\I_n$ and G-Broyden uses $\mathbf P_0=[\0,L\I_n]^{\dag}$. We tune $\eta\in\{0.01,0.1,1\}$ for G-Broyden and tune $L,M\in\{1,10,100,1000,10000\}$ for RQGN.

% As Figure~\ref{fig:additional-eigen} shows, the two methods have comparable iteration counts under the exact initialization, although G-Broyden is slightly faster because it avoids Jacobian-vector products. Under the inexact initialization, G-Broyden may fail to converge, whereas RQGN-SR1 retains its superlinear behavior. This agrees with the initialization-robustness theory in the main text.

% \subsection{Minimax Optimization}
% \label{app:additional_minimax}

% On the ``law school'' data set, we compare RQGN-SR1 with G-Broyden (Algorithm~\ref{alg:gbroyden}) and the randomized Broyden method~\citep{ye2021greedy}.

% \begin{algorithm}[H]
% \caption{Randomized Broyden method (R-Broyden)}
% \label{alg:rbroyden}
% \begin{algorithmic}[1]
% \STATE \textbf{Input:} $\x_0\in\RB^d$ and $\B_0\in\RB^{d\times d}$
% \STATE \textbf{for} $t=0,1,\ldots$
% \STATE \quad $\x_{t+1}=\x_t-\B_t^{-1}\F(\x_t)$
% \STATE \quad Choose $\e_i$ uniformly from $\{\e_1,\ldots,\e_d\}$
% \STATE \quad $\B_{t+1}=\B_t+(\J(\x_{t+1})-\B_t)\e_i\e_i^{\top}$
% \STATE \textbf{end for}
% \end{algorithmic}
% \end{algorithm}

% In Setting~1, RQGN, G-Broyden, and randomized Broyden use $\G_0=\J(\x_0)^{\top}\J(\x_0)$, $\mathbf P_0=\J(\x_0)^{-1}$, and $\B_0=\J(\x_0)$, respectively. In Setting~2, they use $\G_0=L^2\I_d$, $\mathbf P_0=\widetilde{\J}/L$, and $\B_0=L\widetilde{\J}$, respectively, where $\widetilde{\J}=\operatorname{diag}(1,\ldots,1,-1)$. We tune $\eta\in\{0.01,0.1,1\}$ for G-Broyden and tune $L,M\in\{1,10,100,1000,10000\}$.

% Figure~\ref{fig:additional-minimax} shows that all three methods have comparable iteration counts under exact initialization, while RQGN incurs more CPU time because of its Jacobian-vector products. Under inexact initialization, RQGN-SR1 substantially outperforms both baselines.

% \begin{figure}[p]
% \centering
% \begingroup
% \setlength{\tabcolsep}{1pt}
% \renewcommand{\arraystretch}{0.96}
% \begin{tabular}{@{}cc@{}}
% \includegraphics[width=0.327\textwidth,trim=11pt 6pt 11pt 6pt,clip]{app_exp/Eigen_300_exact_Jaco.res.pdf} &
% \includegraphics[width=0.327\textwidth,trim=11pt 6pt 11pt 6pt,clip]{app_exp/Eigen_300_exact_Jaco.time.pdf} \\[-1pt]
% {\footnotesize (a) Setting 1 (iteration)} & {\footnotesize (b) Setting 1 (time)} \\[2pt]
% \includegraphics[width=0.327\textwidth,trim=11pt 6pt 11pt 6pt,clip]{app_exp/Eigen_300_inexact_Jaco.res.pdf} &
% \includegraphics[width=0.327\textwidth,trim=11pt 6pt 11pt 6pt,clip]{app_exp/Eigen_300_inexact_Jaco.time.pdf} \\[-1pt]
% {\footnotesize (c) Setting 2 (iteration)} & {\footnotesize (d) Setting 2 (time)}
% \end{tabular}
% \endgroup
% \caption{Iteration counts and CPU time (seconds) versus $\|\F(\x)\|$ for the underdetermined eigenvalue problem, using different initial estimators.}
% \label{fig:additional-eigen}
% \end{figure}

% \begin{figure}[p]
% \centering
% \begingroup
% \setlength{\tabcolsep}{1pt}
% \renewcommand{\arraystretch}{0.96}
% \begin{tabular}{@{}cc@{}}
% \includegraphics[width=0.327\textwidth,trim=11pt 6pt 11pt 6pt,clip]{app_minimax/lawschool_exact_Jaco.res.pdf} &
% \includegraphics[width=0.327\textwidth,trim=11pt 6pt 11pt 6pt,clip]{app_minimax/lawschool_exact_Jaco.time.pdf} \\[-1pt]
% {\footnotesize (a) Setting 1 (iteration)} & {\footnotesize (b) Setting 1 (time)} \\[2pt]
% \includegraphics[width=0.327\textwidth,trim=11pt 6pt 11pt 6pt,clip]{app_minimax/lawschool_inexact_Jaco.res.pdf} &
% \includegraphics[width=0.327\textwidth,trim=11pt 6pt 11pt 6pt,clip]{app_minimax/lawschool_inexact_Jaco.time.pdf} \\[-1pt]
% {\footnotesize (c) Setting 2 (iteration)} & {\footnotesize (d) Setting 2 (time)}
% \end{tabular}
% \endgroup
% \caption{Iteration counts and CPU time (seconds) versus $\|\F(\x)\|$ on the ``law school'' data set, using different initial estimators.}
% \label{fig:additional-minimax}
% \end{figure}

\FloatBarrier


\bibliographystyle{plainnat}
\bibliography{ref}
\end{document}
